% Springer-compatible LaTeX template for APJM / MIR submission
% Compiled from Markdown via pandoc.
% Document class: `article` (Springer Nature's sn-jnl is not bundled; this template
% follows Springer's general manuscript guidelines and is compatible with both
% Springer Nature LaTeX template and Word-track submission portals).

\documentclass[12pt,a4paper]{article}

% --- Geometry & spacing ---
\usepackage[margin=2.5cm]{geometry}
\usepackage{setspace}
\onehalfspacing

% --- Encoding & fonts ---
% Conditional setup: under pdflatex use legacy inputenc; under xelatex/lualatex
% use fontspec for native Unicode support (needed for β, −, × symbols).
\usepackage{iftex}
\ifPDFTeX
  \usepackage[utf8]{inputenc}
  \usepackage[T1]{fontenc}
  \usepackage{lmodern}
\else
  \usepackage{fontspec}
  \setmainfont{Liberation Serif}[
    Scale=1.0,
    BoldFont={Liberation Serif Bold},
    ItalicFont={Liberation Serif Italic},
    BoldItalicFont={Liberation Serif Bold Italic}
  ]
  % Liberation Mono carries the full Vietnamese diacritic set; DejaVu Sans
  % Mono does not (renders Ấ/ầ/ế/ố as missing characters in code blocks).
  \setmonofont{Liberation Mono}[Scale=0.85]
  % Map common Unicode math/typographic symbols that may appear in markdown
  % body text. Liberation Serif lacks them, so route via amssymb/math mode.
  % (We deliberately avoid unicode-math, which conflicts with the amssymb
  % package that pandoc auto-loads in the preamble.)
  \usepackage{newunicodechar}
  \newunicodechar{∈}{\ensuremath{\in}}
  \newunicodechar{∉}{\ensuremath{\notin}}
  \newunicodechar{≤}{\ensuremath{\leq}}
  \newunicodechar{≥}{\ensuremath{\geq}}
  \newunicodechar{≠}{\ensuremath{\neq}}
  \newunicodechar{≈}{\ensuremath{\approx}}
  \newunicodechar{∞}{\ensuremath{\infty}}
  \newunicodechar{∑}{\ensuremath{\sum}}
  \newunicodechar{∏}{\ensuremath{\prod}}
  \newunicodechar{∗}{\ensuremath{\ast}}
  \newunicodechar{·}{\ensuremath{\cdot}}
  \newunicodechar{−}{\ensuremath{-}}
  \newunicodechar{×}{\ensuremath{\times}}
  \newunicodechar{÷}{\ensuremath{\div}}
  \newunicodechar{±}{\ensuremath{\pm}}
  % Typographic checkmarks/crosses (used as bullet markers in markdown).
  \newunicodechar{✓}{\ensuremath{\checkmark}}
  \newunicodechar{✗}{\ensuremath{\times}}
  \newunicodechar{⚠}{\textbf{!}}
  \newunicodechar{→}{\ensuremath{\rightarrow}}
  \newunicodechar{←}{\ensuremath{\leftarrow}}
  \newunicodechar{⇒}{\ensuremath{\Rightarrow}}
  % Allow URL line breaks at slashes (fix Overfull \hbox in long URLs)
  \usepackage{xurl}
\fi

% --- Math ---
\usepackage{amsmath,amssymb,amsfonts,mathtools}
\usepackage{bm}

% --- Tables ---
\usepackage{booktabs}
\usepackage{longtable}
\usepackage{multirow}
\usepackage{array}
\usepackage{tabularx}

% --- Figures ---
\usepackage{graphicx}
\usepackage{float}
\usepackage{caption}
\captionsetup{font=small,labelfont=bf}

% --- Misc text ---
\usepackage{microtype}
\usepackage{enumitem}
\setlist{nosep,leftmargin=1.5em}
\usepackage{textcomp}
\usepackage{xcolor}
\usepackage{ulem}\normalem

% --- Code block support (pandoc Shaded / Highlighting environments) ---
% Required when source markdown contains fenced code blocks (```...```).
\usepackage{fancyvrb}
\usepackage{framed}
\definecolor{shadecolor}{RGB}{248,248,248}
\newenvironment{Shaded}{\begin{snugshade}}{\end{snugshade}}
% Highlighting environment for syntax-highlighted code
\providecommand{\BuiltInTok}[1]{#1}
\providecommand{\CommentTok}[1]{\textit{#1}}
\providecommand{\KeywordTok}[1]{\textbf{#1}}
\providecommand{\NormalTok}[1]{#1}
\providecommand{\StringTok}[1]{\textit{#1}}
\providecommand{\OperatorTok}[1]{#1}
\providecommand{\DataTypeTok}[1]{#1}
\providecommand{\FunctionTok}[1]{#1}

% --- Links & references ---
\usepackage[hidelinks,breaklinks=true]{hyperref}
\usepackage{cleveref}

% --- Bibliography (APA7 via natbib + apalike) ---
% APJM/MIR accept APA7. natbib's apalike close approximation; for full APA7
% compliance use biblatex with apa style.
\usepackage[round,sort]{natbib}
\bibliographystyle{apalike}

% --- Pandoc compatibility shims ---
\providecommand{\tightlist}{\setlength{\itemsep}{0pt}\setlength{\parskip}{0pt}}
\providecommand{\passthrough}[1]{\lstset{mathescape=false}#1\lstset{mathescape=true}}

% --- Title page (filled in per paper) ---
\title{Luận án tiến sĩ --- 5 chương --- Tác động của quốc tế hoá đến
hiệu quả doanh nghiệp châu Á}
\author{Author details withheld for blind review}
\date{2026-05-21}

\begin{document}

\maketitle

\hypertarget{chux1b0ux1a1ng-1-giux1edbi-thiux1ec7u}{%
\section{CHƯƠNG 1: GIỚI
THIỆU}\label{chux1b0ux1a1ng-1-giux1edbi-thiux1ec7u}}

\begin{center}\rule{0.5\linewidth}{0.5pt}\end{center}

\hypertarget{11-bux1ed1i-cux1ea3nh-nghiuxean-cux1ee9u}{%
\subsection{1.1 Bối cảnh nghiên
cứu}\label{11-bux1ed1i-cux1ea3nh-nghiuxean-cux1ee9u}}

Quốc tế hóa doanh nghiệp là một trong những chủ đề trung tâm của kinh
doanh quốc tế trong hơn nửa thế kỷ qua. Kể từ khi Johanson và Vahlne
(1977) đề xuất mô hình Uppsala về quá trình quốc tế hóa tuần tự, hàng
thế kỷ nghiên cứu đã tích lũy nhằm trả lời câu hỏi cốt lõi: liệu quốc tế
hóa có thực sự cải thiện hiệu quả hoạt động của doanh nghiệp hay không,
và trong điều kiện nào điều đó xảy ra? Lu và Beamish (2004) cung cấp
bằng chứng quan trọng về quan hệ hình chữ S (S-curve) giữa mức độ quốc
tế hóa và hiệu quả, trong khi Contractor et al. (2003) và nhiều nghiên
cứu tiếp theo tiếp tục thách thức, bổ sung và điều chỉnh kết luận đó. Dù
tranh luận học thuật vẫn chưa ngã ngũ, sự đồng thuận hiện nay cho thấy
mối quan hệ internationalization--firm performance (I--P) là có thực
nhưng phức tạp, phụ thuộc vào nhiều điều kiện biên.

Ở châu Á, bối cảnh này càng trở nên đặc biệt đáng quan tâm vì khu vực
đang chứng kiến tốc độ hội nhập kinh tế và chuyển đổi số hết sức không
đồng đều giữa các quốc gia. Từ Singapore, một nền kinh tế nhỏ mở, dẫn
đầu thế giới về năng lực đổi mới và kết nối toàn cầu, đến Việt Nam, một
nền kinh tế chuyển đổi đang đẩy nhanh cải cách thể chế và hội nhập quốc
tế, rồi đến Trung Quốc với quy mô thị trường khổng lồ và cấu trúc thể
chế đặc thù, thậm chí đến các quốc đảo nhỏ Thái Bình Dương đang phải đối
mặt với thách thức sinh tồn kép của biến đổi khí hậu và bẫy phụ thuộc
xuất khẩu, sự đa dạng này tạo ra một phòng thí nghiệm thực địa lý tưởng
để kiểm định các điều kiện biên của quan hệ I--P (Khanna \& Palepu,
2010; Peng, 2003).

Chuyển đổi số đang làm thay đổi căn bản cách thức doanh nghiệp quốc tế
hóa và cách thức quốc tế hóa ảnh hưởng đến hiệu quả hoạt động. Banalieva
và Dhanaraj (2019) lập luận rằng lý thuyết nội địa hóa cần được cập nhật
cho nền kinh tế số, vì dữ liệu số và tương tác qua nền tảng làm giảm sự
phụ thuộc vào hiện diện vật lý. Stallkamp và Schotter (2021) chỉ ra rằng
các nền tảng số quốc tế hóa theo logic hoàn toàn khác các công ty đa
quốc gia truyền thống. Verhoef et al. (2021) cung cấp khung phân tầng ba
cấp từ digitization đến digitalization và digital transformation. Tuy
nhiên, phần lớn bằng chứng thực nghiệm hiện có về vai trò của năng lực
số trong quan hệ I--P vẫn tập trung vào các nền kinh tế phát triển, đặc
biệt là Bắc Mỹ và Tây Âu, trong khi bằng chứng cho châu Á mới nổi còn
rất hạn chế (Bhandari et al., 2023).

Không thể bỏ qua vai trò của thể chế trong việc hình thành kết quả quốc
tế hóa. Institutional Theory (North, 1990; Khanna \& Palepu, 2010) khẳng
định rằng thể chế tạo ra cấu trúc khuyến khích và chi phí giao dịch khác
nhau giữa các quốc gia. Ở các thị trường mới nổi, "institutional voids",
tức sự thiếu hụt các thể chế trung gian hỗ trợ giao dịch thị trường, tạo
ra chi phí bổ sung đáng kể cho doanh nghiệp quốc tế hóa (Peng, 2003). Sự
đa dạng thể chế giữa các quốc gia châu Á, từ những nền kinh tế với quản
trị nhà nước ổn định, pháp quyền mạnh như Singapore đến các nền kinh tế
đang phát triển thể chế như Campuchia hay các quốc đảo Thái Bình Dương,
tạo ra gradient thể chế đặc biệt phong phú để kiểm định. Luận án này
tiếp cận gradient đó một cách hệ thống thông qua khung phân loại thể chế
6 nhóm được phát triển cho luận án này, gọi tắt là ICRV (Institutional
Context Regime Variation --- Biến thiên chế độ bối cảnh thể chế).

Cuối cùng, nhà quản trị cấp cao đóng vai trò không thể bỏ qua trong việc
hình thành chiến lược và kết quả quốc tế hóa. Upper Echelons Theory
(Hambrick \& Mason, 1984) khẳng định rằng quyết định chiến lược phản ánh
đặc điểm nhận thức, kinh nghiệm và giá trị của nhà quản trị. Kinh nghiệm
quốc tế giúp nhà quản trị nhận diện cơ hội và giảm chi phí của liability
of foreignness, trong khi sự đa dạng giới tính ở cấp lãnh đạo có thể
tăng cường quản trị rủi ro và ra quyết định đa chiều (Hsu et al., 2013;
Post \& Byron, 2015). Tích hợp tầng nhà quản trị vào khung giải thích đa
tầng là cần thiết để phản ánh đầy đủ sự phức tạp của quan hệ I--P trong
bối cảnh châu Á.

\begin{center}\rule{0.5\linewidth}{0.5pt}\end{center}

\hypertarget{12-vux1ea5n-ux111ux1ec1-nghiuxean-cux1ee9u}{%
\subsection{1.2 Vấn đề nghiên
cứu}\label{12-vux1ea5n-ux111ux1ec1-nghiuxean-cux1ee9u}}

Mặc dù đã có hàng trăm nghiên cứu thực nghiệm, bằng chứng về mối quan hệ
giữa quốc tế hóa và hiệu quả hoạt động của doanh nghiệp vẫn chưa thống
nhất và tản mạn. Các phân tích tổng hợp (meta-analyses) về quan hệ I--P
liên tục cho thấy tác động tổng hợp là dương nhưng nhỏ, hệ số tương quan
tổng hợp r trong khoảng 0,07--0,13, đồng thời mức độ dị biệt giữa các
nghiên cứu rất cao với I² thường vượt 80\% (Bausch \& Krist, 2007; Kirca
et al., 2012; Marano et al., 2016; Wu et al., 2022). Tình trạng dị biệt
cao này không phải là tín hiệu của sai sót phương pháp, mà là bằng chứng
cho thấy quan hệ I--P phụ thuộc sâu sắc vào điều kiện bối cảnh. Arte và
Larimo (2022) lập luận rằng heterogeneity cao gợi ý sự cần thiết phải
chuyển dịch từ mô hình đơn biến sang khung đa tầng có tính đến thể chế,
năng lực và đặc điểm nhà quản trị.

Ba khoảng trống nghiên cứu quan trọng được xác định từ tổng quan tài
liệu. Khoảng trống thứ nhất là thiếu khung lý thuyết tích hợp đủ sức
giải thích heterogeneity cao đó. Các nghiên cứu hiện tại thường chọn một
hoặc hai lý thuyết làm nền, dẫn đến mô hình không đủ đa tầng để nắm bắt
sự phức tạp thực tế. Uppsala model giải thích chi phí học hỏi nhưng
không đủ để xử lý vai trò của thể chế; RBV giải thích sự khác biệt năng
lực nhưng bỏ qua macro-context; Institutional Theory giải thích bối cảnh
vĩ mô nhưng không tích hợp đặc điểm cá nhân nhà quản trị. Không một lý
thuyết nào đứng riêng có thể giải thích đầy đủ heterogeneity đó. Khoảng
trống thứ hai là sự đồng nhất không hợp lý giữa hai construct số khác
nhau về bản chất: technological capability, chiều sâu năng lực công nghệ
nhập vào tổ chức, và digital adoption, mức độ áp dụng công cụ số bề mặt
trong hoạt động kinh doanh. Nhiều nghiên cứu trước gộp hai construct này
thành một chỉ số composite duy nhất, điều này làm giảm độ chính xác của
kết luận và có thể che khuất các cơ chế tác động khác nhau (Bharadwaj et
al., 2013; Coltman et al., 2008; Verhoef et al., 2021). Khoảng trống thứ
ba là thiếu bằng chứng cross-country quy mô lớn cho châu Á. Phần lớn
nghiên cứu hiện có tập trung vào một quốc gia hoặc một tiểu vùng, và
thậm chí khi có nghiên cứu đa quốc gia, phạm vi thường hẹp và thiếu sự
đa dạng thể chế đủ lớn để kiểm định các điều kiện biên có ý nghĩa.

Luận án này được thiết kế để giải quyết ba khoảng trống đó bằng cách xây
dựng và kiểm định một khung giải thích đa tầng cho quan hệ I--P tại châu
Á, với thiết kế mixed synthesis-empirical kết hợp meta-analysis và phân
tích đa quốc gia quy mô lớn. Trọng tâm khoa học của luận án là vai trò
điều tiết của thể chế, năng lực số và đặc điểm nhà quản trị, ba tầng
điều tiết mà, theo lập luận của luận án, cùng nhau tạo nên cấu trúc giải
thích cần thiết cho heterogeneity cao trong quan hệ I--P.

\begin{center}\rule{0.5\linewidth}{0.5pt}\end{center}

\hypertarget{13-mux1ee5c-tiuxeau-nghiuxean-cux1ee9u}{%
\subsection{1.3 Mục tiêu nghiên
cứu}\label{13-mux1ee5c-tiuxeau-nghiuxean-cux1ee9u}}

Mục tiêu tổng quát của luận án là xây dựng và kiểm định một khung giải
thích đa tầng cho mối quan hệ giữa quốc tế hóa và hiệu quả hoạt động của
doanh nghiệp tại các quốc gia châu Á, trong đó làm rõ vai trò điều tiết
của thể chế, năng lực số và đặc điểm nhà quản trị.

Từ mục tiêu tổng quát đó, luận án triển khai bốn mục tiêu cụ thể. Mục
tiêu cụ thể thứ nhất là tổng hợp và cập nhật bằng chứng meta-analytic về
mối quan hệ I--P trong giai đoạn 1982--2026, xác định tác động tổng hợp,
mức độ dị biệt và các nhân tố điều tiết (moderators) chính giải thích
heterogeneity trong literature. Mục tiêu cụ thể thứ hai là kiểm định cơ
chế cụ thể trong các bối cảnh quốc gia khác nhau, bao gồm Singapore với
tư cách nền kinh tế nhỏ mở tiên tiến, Việt Nam với tư cách nền kinh tế
chuyển đổi đang mở rộng, và Trung Quốc với tư cách nền kinh tế mới nổi
quy mô lớn, nhằm làm rõ cách thức các cơ chế I--P vận hành khác nhau
theo bối cảnh quốc gia. Mục tiêu cụ thể thứ ba là mở rộng phân tích sang
47--49 nền kinh tế châu Á và Thái Bình Dương (bao gồm 6--8 nền kinh tế
Pacific SIDS) để kiểm định khả năng khái quát hóa của các phát hiện trên
phạm vi rộng, kế thừa và mở rộng từ pool 17 nền kinh tế trong Do và Phan
(2026a). Mục tiêu cụ thể thứ tư là phân tích vai trò điều tiết của ba
nhóm biến: thể chế (bao gồm các sub-regime ICRV và các chỉ số rào cản
thể chế), năng lực số (phân biệt rõ technological capability index, TCI
và chỉ số chấp nhận số, digital adoption index --- DAI), và đặc điểm nhà
quản trị cấp cao (kinh nghiệm và giới tính).

\begin{center}\rule{0.5\linewidth}{0.5pt}\end{center}

\hypertarget{14-cuxe2u-hux1ecfi-nghiuxean-cux1ee9u}{%
\subsection{1.4 Câu hỏi nghiên
cứu}\label{14-cuxe2u-hux1ecfi-nghiuxean-cux1ee9u}}

Câu hỏi nghiên cứu trung tâm của luận án là: Quốc tế hóa ảnh hưởng đến
hiệu quả hoạt động của doanh nghiệp ở châu Á như thế nào, và những điều
kiện nào, về thể chế, năng lực số và đặc điểm nhà quản trị, giải thích
sự dị biệt cao trong mối quan hệ đó?

Từ câu hỏi trung tâm đó, bốn câu hỏi nghiên cứu cụ thể được đặt ra. Câu
hỏi nghiên cứu thứ nhất (RQ1) là: Tác động tổng hợp của quốc tế hóa lên
hiệu quả hoạt động của doanh nghiệp trong literature 1982--2026 là gì,
mức độ dị biệt ở mức nào, và những moderators nào giải thích
heterogeneity đó một cách có hệ thống? Câu hỏi nghiên cứu thứ hai (RQ2)
là: Trong từng bối cảnh quốc gia cụ thể bao gồm Singapore, Việt Nam và
Trung Quốc, mối quan hệ I--P vận hành theo những cơ chế đặc thù nào và
sự phi tuyến có biểu hiện ra sao? Câu hỏi nghiên cứu thứ ba (RQ3) là:
Trên phạm vi đa quốc gia châu Á bao gồm 47--49 nền kinh tế châu Á và
Thái Bình Dương (bao gồm 6--8 nền kinh tế Pacific SIDS), vai trò điều
tiết của chỉ số chấp nhận số (DAI) và technological capability index
(TCI) có khái quát hóa được không, và hai construct này có thực sự vận
hành theo các cơ chế khác nhau không? Câu hỏi nghiên cứu thứ tư (RQ4)
là: Vai trò của thể chế, đo bằng sub-regime ICRV, cùng với sự phân biệt
giữa TCI và DAI, và đặc điểm nhà quản trị trong việc điều tiết mối quan
hệ I--P là gì, và sự tương tác ba chiều giữa ba tầng đó có ý nghĩa như
thế nào?

\begin{center}\rule{0.5\linewidth}{0.5pt}\end{center}

\hypertarget{15-ux111ux1ed1i-tux1b0ux1ee3ng-vuxe0-phux1ea1m-vi-nghiuxean-cux1ee9u}{%
\subsection{1.5 Đối tượng và phạm vi nghiên
cứu}\label{15-ux111ux1ed1i-tux1b0ux1ee3ng-vuxe0-phux1ea1m-vi-nghiuxean-cux1ee9u}}

\hypertarget{151-ux111ux1ed1i-tux1b0ux1ee3ng-nghiuxean-cux1ee9u}{%
\subsubsection{1.5.1 Đối tượng nghiên
cứu}\label{151-ux111ux1ed1i-tux1b0ux1ee3ng-nghiuxean-cux1ee9u}}

Đối tượng nghiên cứu của luận án là mối quan hệ giữa quốc tế hóa và hiệu
quả hoạt động của doanh nghiệp (internationalization--firm performance
relationship) cùng với các yếu tố điều tiết của mối quan hệ đó, bao gồm
bối cảnh thể chế, năng lực số và đặc điểm nhà quản trị cấp cao. Đơn vị
phân tích là doanh nghiệp ở cấp độ tổ chức (firm-level analysis), không
phải cấp độ ngành hay cấp độ quốc gia.

\hypertarget{152-phux1ea1m-vi-khuxf4ng-gian}{%
\subsubsection{1.5.2 Phạm vi không
gian}\label{152-phux1ea1m-vi-khuxf4ng-gian}}

Phạm vi không gian chính của nghiên cứu bao gồm \textbf{47--49 nền kinh
tế châu Á và Thái Bình Dương} có dữ liệu trong Khảo sát Doanh nghiệp của
Ngân hàng Thế giới (World Bank Enterprise Surveys --- WBES), được phân
loại và tổ chức theo 6 sub-regime trong khung ICRV. Sáu sub-regime bao
gồm: (i) Advanced Innovation-Driven, đại diện là Singapore, Hàn Quốc và
Đài Loan; (ii) Advanced Resource-Driven, đại diện là các nền kinh tế
GCC; (iii) Upper-middle, đại diện là Trung Quốc và Malaysia; (iv)
Emerging, bao gồm Việt Nam và các nền kinh tế ASEAN đang phát triển; (v)
Frontier, bao gồm các nền kinh tế kém phát triển nhất châu Á như
Campuchia, Nepal và Afghanistan; và (vi) Pacific SIDS (Small Island
Developing States), bao gồm 9 nền kinh tế đảo nhỏ Thái Bình Dương được
đưa vào phân tích để kiểm định các điều kiện biên cực đoan của quan hệ
I--P. Nhật Bản không nằm trong phạm vi thực nghiệm của luận án do cấu
trúc dữ liệu WBES không bao phủ Nhật Bản đầy đủ.

\hypertarget{153-phux1ea1m-vi-thux1eddi-gian}{%
\subsubsection{1.5.3 Phạm vi thời
gian}\label{153-phux1ea1m-vi-thux1eddi-gian}}

Phạm vi thời gian của nghiên cứu có hai mức. Đối với phase
meta-analysis, luận án tổng hợp literature trong giai đoạn 1982--2026,
bao quát toàn bộ chương trình nghiên cứu về quan hệ I--P từ khi các
nghiên cứu thực nghiệm đầu tiên xuất hiện đến thời điểm hoàn thành luận
án. Đối với phase phân tích thực nghiệm đa quốc gia, luận án sử dụng dữ
liệu WBES trong giai đoạn 2009--2025, bao gồm 101.185 doanh nghiệp với
108 cặp quốc gia-năm trải qua 14 mốc khảo sát. Dữ liệu WBES từ ba thế hệ
schema (PICS3 giai đoạn 2009--2012, Standardized giai đoạn 2013--2017,
và BREADY/BEE giai đoạn 2018--2025) đã được hòa hợp để đảm bảo tính nhất
quán trong đo lường biến qua các wave khảo sát.

\begin{center}\rule{0.5\linewidth}{0.5pt}\end{center}

\hypertarget{16-phux1b0ux1a1ng-phuxe1p-luux1eadn-vuxe0-phux1b0ux1a1ng-phuxe1p-nghiuxean-cux1ee9u}{%
\subsection{1.6 Phương pháp luận và phương pháp nghiên
cứu}\label{16-phux1b0ux1a1ng-phuxe1p-luux1eadn-vuxe0-phux1b0ux1a1ng-phuxe1p-nghiuxean-cux1ee9u}}

Luận án áp dụng thiết kế nghiên cứu mixed synthesis-empirical, sự kết
hợp giữa meta-analysis và phân tích thực nghiệm đa quốc gia, nhằm trả
lời bốn câu hỏi nghiên cứu đã đặt ra. Thiết kế này không chỉ đơn thuần
là sử dụng hai phương pháp song song, mà tạo ra một đối thoại hai chiều:
meta-analysis cung cấp tổng quan có hệ thống về trạng thái bằng chứng và
xác định các moderator quan trọng trong literature, trong khi phân tích
thực nghiệm kiểm định trực tiếp các giả thuyết phát sinh từ khung lý
thuyết trong bối cảnh châu Á. Sự tương tác này tăng cường đáng kể sức
mạnh đóng góp so với việc chỉ sử dụng một phương pháp duy nhất
(Borenstein et al., 2009; Hunter \& Schmidt, 2004).

Về phương pháp cụ thể, phase meta-analysis sử dụng random-effects model
với Fisher\textquotesingle s z transformation, phân tích nhóm con
(subgroup analysis) theo chế độ thể chế, và các kiểm định publication
bias bao gồm Egger\textquotesingle s test và Begg--Mazumdar rank
correlation test. Phase thực nghiệm sử dụng hồi quy OLS với sai số chuẩn
robust HC1/HC3, mô hình bậc hai và bậc ba để kiểm định phi tuyến, tương
tác hai chiều và ba chiều để kiểm định moderation, cùng với fixed
effects theo quốc gia và năm. Kiểm định Lind--Mehlum U-test được áp dụng
để xác nhận dạng inverted-U và xác định điểm turning point trong quan hệ
phi tuyến (Haans et al., 2016; Lind \& Mehlum, 2010).

\begin{center}\rule{0.5\linewidth}{0.5pt}\end{center}

\hypertarget{17-giux1ea3-thuyux1ebft-khoa-hux1ecdc-tuxf3m-tux1eaft}{%
\subsection{1.7 Giả thuyết khoa học (tóm
tắt)}\label{17-giux1ea3-thuyux1ebft-khoa-hux1ecdc-tuxf3m-tux1eaft}}

Dựa trên khung lý thuyết tích hợp đa tầng bao gồm Uppsala model, RBV,
Institutional Theory, Upper Echelons Theory và digital capability lens,
luận án đề xuất hệ giả thuyết khoa học với sáu giả thuyết trọng tâm.
Quan hệ nền giữa quốc tế hóa và hiệu quả hoạt động doanh nghiệp mang
dạng phi tuyến, cụ thể là dạng inverted-U, với chi phí thâm nhập và
thích nghi cao ở giai đoạn đầu và lợi ích quy mô, học hỏi và đa dạng hóa
thị trường ở giai đoạn sau, đến khi chi phí điều phối vượt quá lợi ích
sẽ dẫn đến sự suy giảm hiệu quả. Technological capability index (TCI)
điều tiết dương quan hệ I--P bằng cách mở rộng biên thích nghi và tăng
cường khả năng khai thác lợi thế cạnh tranh ở thị trường quốc tế. Chỉ số
chấp nhận số (DAI), được hiểu là proxy foundational digital adoption, có
tác động điều tiết không đơn điệu và phụ thuộc vào chế độ thể chế: ở các
nền kinh tế có thể chế tốt, DAI khuếch đại lợi ích của quốc tế hóa; ở
các nền kinh tế với thể chế yếu, DAI có thể khuếch đại chi phí điều phối
khi mức độ quốc tế hóa vượt ngưỡng. Bối cảnh thể chế, đo bằng sub-regime
ICRV, điều tiết cả dạng và độ lớn của quan hệ I--P, tạo ra gradient từ
tác động dương lớn ở Advanced Innovation-Driven đến tác động âm ở
Frontier và Pacific SIDS extreme. Đặc điểm nhà quản trị cấp cao, cụ thể
là kinh nghiệm quốc tế và đa dạng giới tính, điều tiết dương quan hệ
I--P thông qua cơ chế tăng cường nhận diện cơ hội và quản trị rủi ro.
Ngoài sáu giả thuyết trọng tâm H1--H6, luận án phát biểu thêm giả thuyết
điều kiện biên H1b: tại các nền kinh tế Pacific SIDS (ICRV Nhóm VI), nơi
ba điều kiện cấu trúc tối thiểu đồng thời bị vi phạm, thị trường nội địa
quá nhỏ, chi phí thương mại cực cao, và thiếu hỗ trợ thể chế, quan hệ
I--P chuyển từ dạng chữ U ngược sang quan hệ \textbf{âm đơn điệu không
có điểm uốn}, được gọi là gánh nặng quốc tế hóa bắt buộc (Forced
Internationalization Penalty --- FIP). Chi tiết toàn bộ hệ giả thuyết
H1--H6 và H1b được trình bày đầy đủ trong Chương 2 của luận án.

\begin{center}\rule{0.5\linewidth}{0.5pt}\end{center}

\hypertarget{18-uxfd-nghux129a-khoa-hux1ecdc-vuxe0-thux1ef1c-tiux1ec5n}{%
\subsection{1.8 Ý nghĩa khoa học và thực
tiễn}\label{18-uxfd-nghux129a-khoa-hux1ecdc-vuxe0-thux1ef1c-tiux1ec5n}}

\hypertarget{181-ux111uxf3ng-guxf3p-luxfd-thuyux1ebft}{%
\subsubsection{1.8.1 Đóng góp lý
thuyết}\label{181-ux111uxf3ng-guxf3p-luxfd-thuyux1ebft}}

Về mặt lý thuyết, luận án đóng góp vào literature kinh doanh quốc tế
theo ba hướng chính. Hướng thứ nhất là tích hợp bốn lý thuyết kinh điển,
Uppsala model, Resource-Based View, Institutional Theory và Upper
Echelons Theory, vào một khung giải thích đa tầng thống nhất, bổ sung
bởi digital capability lens để phản ánh chuyển đổi số đương đại. Khung
này không chỉ là sự ghép nối cơ học các lý thuyết mà là một kiến trúc lý
thuyết có tính nhất quán nội tại, trong đó mỗi tầng giải thích một loại
heterogeneity khác nhau trong quan hệ I--P. Hướng thứ hai là phân tách
rõ ràng technological capability (TCI) và chỉ số chấp nhận số (DAI) như
hai construct lý thuyết độc lập, dựa trên các nền tảng khái niệm khác
nhau từ Bharadwaj et al. (2013), Verhoef et al. (2021), và Coltman et
al. (2008). Sự phân tách này bổ sung vào thiếu hụt khái niệm quan trọng
trong literature về số hóa và quốc tế hóa. Hướng thứ ba là tái định vị
Uppsala model trong bối cảnh kỷ nguyên AI bằng cách mở rộng cơ chế học
hỏi để bao gồm "data-augmented learning", phù hợp với xu hướng
digitally-facilitated internationalization (Banalieva \& Dhanaraj, 2019;
Yang et al., 2025). Hướng thứ tư là đề xuất và đặt tên lý thuyết cho
\textbf{gánh nặng quốc tế hóa bắt buộc (Forced Internationalization
Penalty, FIP)}, một construct lý thuyết mới chưa được định danh trong
literature kinh doanh quốc tế, như điều kiện biên cực đoan của mô hình
chữ U ngược khi cả ba tầng của khung CDCM (thể chế, năng lực, quản trị)
đồng thời ở mức tối thiểu. Bằng chứng thực nghiệm từ chín nền kinh tế
Pacific SIDS xác nhận quan hệ âm đơn điệu không có điểm uốn, mở ra hướng
nghiên cứu mới về điều kiện biên của lý thuyết quốc tế hóa.

\hypertarget{182-ux111uxf3ng-guxf3p-phux1b0ux1a1ng-phuxe1p}{%
\subsubsection{1.8.2 Đóng góp phương
pháp}\label{182-ux111uxf3ng-guxf3p-phux1b0ux1a1ng-phuxe1p}}

Về mặt phương pháp, luận án đóng góp theo hai hướng. Hướng thứ nhất là
thiết kế mixed synthesis-empirical kết hợp meta-analysis và phân tích đa
quốc gia, tạo ra khả năng đối chiếu giữa bằng chứng tổng hợp từ
literature toàn cầu và bằng chứng thực nghiệm từ bối cảnh châu Á. Hướng
thứ hai là phát triển và áp dụng khung phân loại thể chế ICRV cho 47 nền
kinh tế, cung cấp một công cụ phân loại mới có ứng dụng tiềm năng trong
nhiều nghiên cứu tương lai về moderating effect của thể chế trong
literature kinh doanh quốc tế.

\hypertarget{183-ux111uxf3ng-guxf3p-thux1ef1c-tiux1ec5n}{%
\subsubsection{1.8.3 Đóng góp thực
tiễn}\label{183-ux111uxf3ng-guxf3p-thux1ef1c-tiux1ec5n}}

Về mặt thực tiễn, luận án cung cấp căn cứ có hệ thống cho hai nhóm đối
tượng sử dụng chính. Đối với doanh nghiệp, kết quả phân tích của luận án
cho thấy tồn tại ngưỡng tối ưu trong mối quan hệ giữa mức độ xuất khẩu
và hiệu quả hoạt động, và ngưỡng này thay đổi theo bối cảnh thể chế và
năng lực số của từng doanh nghiệp. Nhà quản trị cần nhận thức được
ngưỡng này và đầu tư có chiến lược vào năng lực công nghệ thay vì chỉ
tập trung vào digital adoption bề mặt. Đối với nhà hoạch định chính
sách, bằng chứng từ luận án gợi ý rằng cải cách thể chế và hỗ trợ chuyển
đổi số có thể có hiệu ứng bổ trợ tích cực, được nghiên cứu về 17 nền
kinh tế châu Á mới nổi đặt tên là "digital shield effect", trong đó năng
lực số làm giảm tác động tiêu cực của rào cản thể chế đối với doanh
nghiệp quốc tế hóa (Do \& Phan, 2026a). Cơ chế này được luận án kiểm
định trên phạm vi rộng hơn với 47--49 nền kinh tế châu Á và Thái Bình
Dương.

\begin{center}\rule{0.5\linewidth}{0.5pt}\end{center}

\hypertarget{19-tuxednh-mux1edbi-cux1ee7a-ux111ux1ec1-tuxe0i}{%
\subsection{1.9 Tính mới của đề
tài}\label{19-tuxednh-mux1edbi-cux1ee7a-ux111ux1ec1-tuxe0i}}

Luận án có năm điểm mới so với các công trình nghiên cứu trước đây. Điểm
mới thứ nhất là đề xuất khung tích hợp đa tầng cho quan hệ I--P trong
bối cảnh châu Á, kết hợp bốn lý thuyết kinh điển với digital capability
lens thành một kiến trúc lý thuyết thống nhất và có khả năng kiểm định.
Không có nghiên cứu trước nào trong bối cảnh châu Á xây dựng khung tích
hợp đầy đủ ba tầng điều tiết này theo cách có hệ thống. Điểm mới thứ hai
là phân tách TCI và chỉ số chấp nhận số (DAI) như hai construct không
trùng lặp, kiểm định chúng riêng biệt trong cùng một mô hình để làm rõ
cơ chế tác động khác nhau của hai loại năng lực số đó. Sự phân tách này
khắc phục hạn chế của các nghiên cứu trước thường gộp chung hai
construct có bản chất khác nhau. Điểm mới thứ ba là phát triển và áp
dụng khung phân loại thể chế ICRV với 6 sub-regime cho 47 nền kinh tế,
bao gồm cả 6 nền kinh tế Pacific SIDS liền kề như boundary case cực
đoan, đây là lần đầu tiên một hệ thống phân loại đủ mịn như vậy được áp
dụng trong nghiên cứu thực nghiệm về quan hệ I--P tại châu Á. Điểm mới
thứ tư là pool dữ liệu lớn nhất từ trước đến nay cho nghiên cứu đa quốc
gia về quan hệ I--P tại châu Á, với 101.185 doanh nghiệp từ 47--49 nền
kinh tế châu Á và Thái Bình Dương, mở rộng đáng kể từ pool 17 nước (Do
\& Phan, 2026a) là nền tảng hiện có lớn nhất trước khi luận án hoàn
thành. Điểm mới thứ năm là kiểm định temporal stability của turning
point trong quan hệ I--P tại Trung Quốc, cho phép đánh giá liệu cấu trúc
phi tuyến đó có ổn định qua hơn một thập kỷ chuyển đổi số hay không. Kết
quả phân tích của luận án cho thấy turning point tương đối ổn định qua
hai thời điểm quan sát, cung cấp bằng chứng về tính bền vững của cấu
trúc I--P phi tuyến ở Trung Quốc dù bối cảnh kinh tế có nhiều thay đổi.
Điểm mới thứ sáu là phát hiện và đặt tên lý thuyết cho \textbf{gánh nặng
quốc tế hóa bắt buộc (Forced Internationalization Penalty, FIP)} tại
chín nền kinh tế Pacific SIDS, lần đầu tiên trong literature kinh doanh
quốc tế, một trường hợp quan hệ I--P âm đơn điệu không có điểm uốn được
ghi nhận và biện minh lý thuyết bằng ba điều kiện cấu trúc vi phạm đồng
thời, mở ra hướng nghiên cứu mới về điều kiện biên của các lý thuyết
quốc tế hóa phổ quát.

\begin{center}\rule{0.5\linewidth}{0.5pt}\end{center}

\hypertarget{110-cux1ea5u-truxfac-cux1ee7a-luux1eadn-uxe1n}{%
\subsection{1.10 Cấu trúc của luận
án}\label{110-cux1ea5u-truxfac-cux1ee7a-luux1eadn-uxe1n}}

Luận án được tổ chức thành năm chương theo chuẩn luận án tiến sĩ truyền
thống tại Trường Kinh tế, Đại học Cần Thơ. Chương 1, Giới thiệu, đặt vấn
đề nghiên cứu, xác định khoảng trống, trình bày mục tiêu, câu hỏi nghiên
cứu, phạm vi, phương pháp luận, giả thuyết khoa học tóm tắt, ý nghĩa,
tính mới và cấu trúc của luận án. Chương 2, Cơ sở lý thuyết và tổng quan
tài liệu, xây dựng khung lý thuyết tích hợp đa tầng, tổng quan bằng
chứng thực nghiệm đã có, và phát triển hệ giả thuyết H1--H6 dựa trên
khoảng trống và logic lý thuyết. Chương 3, Phương pháp nghiên cứu, trình
bày chi tiết thiết kế nghiên cứu, nguồn dữ liệu, cách đo lường biến, mô
hình phân tích và chiến lược kiểm định độ vững. Chương 4, Kết quả nghiên
cứu, trình bày kết quả ở ba cấp độ phân tích: meta-analysis cập nhật
1982--2026, phân tích theo bối cảnh quốc gia cụ thể, và phân tích đa
quốc gia capstone trên pool đầy đủ 47--49 nền kinh tế. Chương 5, Thảo
luận và kết luận, tích hợp các phát hiện từ Chương 4 vào khung lý thuyết
ở Chương 2, thảo luận đóng góp khoa học, hàm ý quản trị và chính sách,
hạn chế của nghiên cứu và hướng nghiên cứu tiếp theo.

\textbf{Roadmap sáu nghiên cứu thành phần (P3--P8):} Luận án tích hợp
sáu nghiên cứu thành phần được công bố hoặc đang trong quá trình xem xét
tại các tạp chí ISI/Scopus. \textbf{P3} (Việt Nam --- JABS, under
review): kiểm định inverted-U tại Nhóm IV ICRV, turning point =
39,3--46,2\% FSTS, TCI là level-shifter có ý nghĩa. \textbf{P4}
(Singapore --- MIR, under review): kiểm định tại Nhóm I ICRV, turning
point ≈ 88,6\%, DAI Tier-1+2 khuếch đại đường cong I→P. \textbf{P5}
(Trung Quốc --- IJOEM, under review): kiểm định temporal stability,
turning point ổn định qua 2012--2024 (≈ 48,78\%, Paternoster z = 0,82, p
= ,412). \textbf{P6} (meta-analysis ba tầng --- IBR, under review): k =
238 nghiên cứu, r̄ = 0,074, I² = 62,4\%, gradient ICRV Q\_M = 17,35 (df =
4, p = ,002). \textbf{P7} (49 nền kinh tế --- JIBS, under revision): N =
84.910 quan sát, turning point trung bình \textasciitilde36\%, DAI×ICRV
(p = ,012), kiểm định tích hợp CDCM đầu tiên ở quy mô liên quốc gia lớn.
\textbf{P8} (Pacific SIDS --- World Development, under revision):
boundary condition FIP, quan hệ âm đơn điệu
(\(\hat{\beta}_\text{FSTS} = -0{,}404\), \(p < {,}05\)), xác nhận H1b
tại chín nền kinh tế SIDS Thái Bình Dương. Kết quả P3--P8 được trình bày
chi tiết trong Chương 4 và tích hợp vào bàn luận CDCM trong Chương 5.

\begin{center}\rule{0.5\linewidth}{0.5pt}\end{center}

\emph{Chú thích về ký hiệu và thuật ngữ: Trong toàn bộ luận án, "quốc tế
hóa" (internationalization) được đo chủ yếu bằng tỷ lệ doanh thu nước
ngoài trên tổng doanh thu (FSTS) hoặc export intensity; "hiệu quả hoạt
động" (firm performance) được đo chủ yếu bằng năng suất lao động (labor
productivity), bổ sung bằng ROS và sales growth trong các kiểm định độ
vững; TCI là technological capability index; DAI là chỉ số chấp nhận số
(digital adoption index), được hiểu như proxy foundational digital
adoption; ICRV là khung phân loại thể chế 6 nhóm được phát triển cho
luận án này; WBES là World Bank Enterprise Surveys.}

\begin{center}\rule{0.5\linewidth}{0.5pt}\end{center}

\hypertarget{chux1b0ux1a1ng-2-tux1ed5ng-quan-tuxe0i-liux1ec7u}{%
\section{CHƯƠNG 2: TỔNG QUAN TÀI
LIỆU}\label{chux1b0ux1a1ng-2-tux1ed5ng-quan-tuxe0i-liux1ec7u}}

Chương 2 có nhiệm vụ hệ thống hóa các khái niệm cốt lõi và lý thuyết nền
tảng, lược khảo và đánh giá các công trình thực nghiệm đã công bố về mối
quan hệ giữa quốc tế hóa và hiệu quả hoạt động doanh nghiệp, xác định
những khoảng trống còn tồn tại trong literature, và từ đó xây dựng khung
khái niệm tích hợp cùng hệ giả thuyết nghiên cứu H1--H6 cho luận án.
Logic của chương đi từ nền tảng khái niệm đến lý thuyết, từ bằng chứng
thực nghiệm đến khoảng trống, và từ khoảng trống đến mô hình nghiên cứu
và giả thuyết.

\begin{center}\rule{0.5\linewidth}{0.5pt}\end{center}

\hypertarget{21-cuxe1c-khuxe1i-niux1ec7m-cux1ed1t-luxf5i}{%
\subsection{2.1 Các khái niệm cốt
lõi}\label{21-cuxe1c-khuxe1i-niux1ec7m-cux1ed1t-luxf5i}}

\hypertarget{211-quux1ed1c-tux1ebf-huxf3a-doanh-nghiux1ec7p}{%
\subsubsection{2.1.1 Quốc tế hóa doanh
nghiệp}\label{211-quux1ed1c-tux1ebf-huxf3a-doanh-nghiux1ec7p}}

Quốc tế hóa doanh nghiệp (internationalization) được hiểu là mức độ và
phạm vi tham gia của một doanh nghiệp vào các hoạt động kinh tế có yếu
tố quốc tế, bao gồm xuất khẩu hàng hóa và dịch vụ, đầu tư trực tiếp ra
nước ngoài, liên doanh quốc tế, cấp phép công nghệ xuyên biên giới, và
thiết lập mạng lưới phân phối ở thị trường nước ngoài (Sullivan, 1994;
Hitt et al., 2006). Khái niệm này không đồng nhất với "toàn cầu hóa",
vốn mang nghĩa một xu hướng kinh tế vĩ mô, mà là một thuộc tính đo lường
được ở cấp độ doanh nghiệp, phản ánh chiến lược mở rộng quốc tế và mức
độ phụ thuộc vào thị trường nước ngoài.

Trong literature, mức độ quốc tế hóa được vận hành hóa bằng nhiều thước
đo khác nhau. Sullivan (1994) đề xuất một chỉ số tổng hợp đa chiều bao
gồm tỷ trọng doanh thu từ nước ngoài trên tổng doanh thu (foreign sales
to total sales, FSTS), tỷ trọng tài sản ở nước ngoài trên tổng tài sản
(foreign assets to total assets, FATA), và phạm vi địa lý (geographic
scope). Trong những nghiên cứu sử dụng dữ liệu doanh nghiệp từ các nền
kinh tế đang phát triển và từ khảo sát doanh nghiệp của Ngân hàng Thế
giới (World Bank Enterprise Surveys --- WBES), thước đo phổ biến nhất là
cường độ xuất khẩu (export intensity), tức là tỷ lệ phần trăm doanh thu
xuất khẩu trên tổng doanh thu, hay tương đương với FSTS theo nghĩa rộng
(Johanson \& Vahlne, 1977; Lu \& Beamish, 2004). Luận án này sử dụng
FSTS và bình phương FSTS (FSTS²) làm thước đo chính về mức độ quốc tế
hóa, với lý do phù hợp với cấu trúc dữ liệu WBES và nhất quán với các
nghiên cứu I--P sử dụng bộ dữ liệu tương tự (Bhandari et al., 2023; Do
\& Phan, 2026a).

Điều quan trọng cần lưu ý là quốc tế hóa không phải là một trạng thái
nhị phân, mà là một biến liên tục biểu thị cường độ tham gia thị trường
quốc tế, có thể thay đổi theo thời gian và phản ánh quá trình cam kết
tích lũy của doanh nghiệp vào các thị trường nước ngoài (Johanson \&
Vahlne, 1977, 2009).

\hypertarget{212-hiux1ec7u-quux1ea3-houx1ea1t-ux111ux1ed9ng-doanh-nghiux1ec7p}{%
\subsubsection{2.1.2 Hiệu quả hoạt động doanh
nghiệp}\label{212-hiux1ec7u-quux1ea3-houx1ea1t-ux111ux1ed9ng-doanh-nghiux1ec7p}}

Hiệu quả hoạt động doanh nghiệp (firm performance) là một khái niệm đa
chiều không có một định nghĩa hoặc thước đo thống nhất trong literature.
Combs et al. (2005) và Richard et al. (2009) phân biệt ba nhóm thước đo
chính: (i) thước đo dựa trên kế toán (accounting-based), bao gồm tỷ suất
sinh lời trên tổng tài sản (ROA), tỷ suất sinh lời trên vốn chủ sở hữu
(ROE), và tỷ suất sinh lời trên doanh thu (ROS); (ii) thước đo dựa trên
thị trường tài chính (market-based), bao gồm Tobin\textquotesingle s Q
và tổng lợi nhuận cổ đông; và (iii) thước đo dựa trên năng suất và tăng
trưởng (productivity-based), bao gồm năng suất lao động (labor
productivity) và tăng trưởng doanh thu.

Trong bối cảnh nghiên cứu sử dụng dữ liệu WBES, vốn không cung cấp thông
tin về giá cổ phiếu hay giá trị thị trường, thước đo năng suất lao động
là lựa chọn phổ biến nhất và phù hợp nhất về mặt kỹ thuật. Năng suất lao
động được tính bằng logarithm tự nhiên của tỷ số doanh thu hàng năm trên
số lao động toàn thời gian (ln(doanh thu/lao động)), cho phép so sánh
chuẩn hóa giữa các doanh nghiệp thuộc các quốc gia, ngành nghề và quy mô
khác nhau (Do \& Phan, 2026a). Luận án sử dụng thước đo này làm biến phụ
thuộc chính, với ROS và tăng trưởng doanh thu được sử dụng như các thước
đo robustness để kiểm tra tính nhất quán của kết quả.

Cần nhấn mạnh rằng việc chọn thước đo hiệu quả không chỉ là quyết định
kỹ thuật, mà còn phản ánh quan niệm lý thuyết về hiệu quả. Trong bối
cảnh quốc tế hóa ở các nền kinh tế đang phát triển, năng suất lao động
phản ánh trực tiếp khả năng chuyển hóa nguồn lực đầu vào thành kết quả
đầu ra, qua đó thể hiện năng lực cạnh tranh thực chất của doanh nghiệp,
khác với ROA vốn bị ảnh hưởng bởi cấu trúc tài sản và hệ thống kế toán
đặc thù của từng quốc gia.

\hypertarget{213-nux103ng-lux1ef1c-cuxf4ng-nghux1ec7-vuxe0-chux1ec9-sux1ed1-chux1ea5p-nhux1eadn-sux1ed1}{%
\subsubsection{2.1.3 Năng lực công nghệ và chỉ số chấp nhận
số}\label{213-nux103ng-lux1ef1c-cuxf4ng-nghux1ec7-vuxe0-chux1ec9-sux1ed1-chux1ea5p-nhux1eadn-sux1ed1}}

Một trong những đóng góp khái niệm trọng tâm của luận án là sự phân biệt
rõ ràng giữa hai construct: năng lực công nghệ (Technological Capability
Index, TCI) và chỉ số chấp nhận số (Digital Adoption Index, DAI). Trong
nhiều nghiên cứu trước đây, hai construct này thường bị gộp lại thành
một khái niệm chung về "năng lực số" (digital capability), dẫn đến sự mơ
hồ lý thuyết và sai lệch trong đo lường (Bhandari et al., 2023;
Bustamante et al., 2022).

Năng lực công nghệ (TCI) đề cập đến chiều sâu của năng lực công nghệ
được nhập vào tổ chức, phản ánh khả năng hấp thụ, khai thác và phát
triển tri thức công nghệ tiên tiến. Các chỉ báo tiêu biểu bao gồm việc
sở hữu chứng nhận chất lượng quốc tế (ISO hoặc tương đương) và sử dụng
công nghệ được cấp phép từ các đối tác nước ngoài. Về bản chất lý
thuyết, TCI gần với khái niệm năng lực hấp thụ (absorptive capacity) của
Cohen và Levinthal (1990) và năng lực công nghệ của Lall (1992), đây là
chiều sâu nội tại của tổ chức, khó sao chép và tạo ra lợi thế bền vững
theo tinh thần RBV (Barney, 1991).

Chỉ số chấp nhận số (Digital Adoption Index, DAI), ngược lại, phản ánh
mức độ áp dụng các công cụ số tại giao diện ngoại tại của doanh nghiệp,
bao gồm việc sở hữu website, sử dụng email trong giao dịch kinh doanh,
thanh toán điện tử, và bán hàng trực tuyến. Đây là các công cụ số có
tính bề mặt (surface digital tools) theo phân loại của Verhoef et al.
(2021), tức là nằm ở cấp độ số hóa (digitization) và số hóa quy trình
(digitalization), chưa đạt đến cấp độ chuyển đổi số toàn diện (digital
transformation). Trong phạm vi luận án, DAI được đo bằng proxy
"foundational website adoption" từ dữ liệu WBES, đây là ràng buộc đo
lường có chủ đích do giới hạn của bộ dữ liệu, không phải thiếu sót lý
thuyết.

Hai construct TCI và DAI không đồng nhất và vận hành theo cơ chế khác
nhau trong mối quan hệ quốc tế hóa, hiệu quả: TCI hoạt động chủ yếu như
một nhân tố nâng mặt bằng hiệu quả (level shifter), trong khi DAI có xu
hướng làm thay đổi độ dốc hoặc độ cong của đường quan hệ I--P tùy theo
bối cảnh (Bhandari et al., 2023; Chen \& Meng, 2022). Việc duy trì sự
phân biệt rõ ràng này trong suốt luận án là điều kiện tiên quyết để đảm
bảo tính thuần khiết của cấu trúc (construct purity) theo tiêu chuẩn của
Coltman et al. (2008).

\hypertarget{214-bux1ed1i-cux1ea3nh-thux1ec3-chux1ebf-vuxe0-khung-icrv}{%
\subsubsection{2.1.4 Bối cảnh thể chế và khung
ICRV}\label{214-bux1ed1i-cux1ea3nh-thux1ec3-chux1ebf-vuxe0-khung-icrv}}

Thể chế (institutions) theo North (1990) là "luật chơi" của một xã hội,
bao gồm các ràng buộc chính thức (luật pháp, quy định, hợp đồng) và phi
chính thức (chuẩn mực, giá trị, văn hóa ứng xử). Thể chế tạo ra cấu trúc
khuyến khích và chi phí giao dịch, do đó ảnh hưởng trực tiếp đến chiến
lược và kết quả kinh doanh của doanh nghiệp.

Để vận hành hóa bối cảnh thể chế đa cấp và đa quốc gia trong luận án,
khung ICRV (Institutional Context Regime Variation --- Biến thiên chế độ
bối cảnh thể chế) được phát triển như một công cụ phân loại có hệ thống.
Khung ICRV chia 47 nền kinh tế trong phạm vi nghiên cứu thành sáu nhóm
(sub-regime), dựa trên tổ hợp các chỉ số quản trị toàn cầu (World
Governance Indicators --- WGI) của Kaufmann et al. (2011), mức độ phát
triển kinh tế, và đặc trưng cấu trúc thể chế: Nhóm I (Advanced
Innovation-Driven, bao gồm Singapore, Hong Kong, Đài Loan, Hàn Quốc),
Nhóm II (Advanced Resource-Driven, bao gồm các nền kinh tế vùng Vịnh),
Nhóm III (Upper-middle, điển hình là Trung Quốc), Nhóm IV (Lower-middle
transition, bao gồm Việt Nam và Ấn Độ), Nhóm V (Frontier), và Nhóm VI
(Pacific SIDS --- Small Island Developing States). Khung này mở rộng từ
phân loại thể chế của Khanna và Palepu (2010) và được điều chỉnh để phản
ánh thực tiễn đa dạng của châu Á và khu vực lân cận.

ICRV không chỉ là một biến phân loại mà còn là một gradient thể chế liên
tục, từ môi trường thể chế mạnh, ít rủi ro (Nhóm I) đến môi trường thể
chế yếu, nhiều rủi ro và khoảng trống thể chế (Nhóm V và VI). Gradient
này tạo ra điều kiện biên (boundary conditions) quan trọng để kiểm định
khả năng tổng quát hóa của mối quan hệ I--P, đặc biệt liên quan đến giả
thuyết phi tuyến.

\begin{center}\rule{0.5\linewidth}{0.5pt}\end{center}

\hypertarget{22-cuxe1c-luxfd-thuyux1ebft-nux1ec1n}{%
\subsection{2.2 Các lý thuyết
nền}\label{22-cuxe1c-luxfd-thuyux1ebft-nux1ec1n}}

\hypertarget{221-muxf4-huxecnh-quux1ed1c-tux1ebf-huxf3a-uppsala}{%
\subsubsection{2.2.1 Mô hình quốc tế hóa
Uppsala}\label{221-muxf4-huxecnh-quux1ed1c-tux1ebf-huxf3a-uppsala}}

Mô hình Uppsala (Johanson \& Vahlne, 1977) là một trong những lý thuyết
có ảnh hưởng sâu rộng nhất trong nghiên cứu về quốc tế hóa doanh nghiệp.
Được xây dựng trên cơ sở quan sát quá trình quốc tế hóa của các doanh
nghiệp Thụy Điển, mô hình đề xuất rằng doanh nghiệp quốc tế hóa theo một
quá trình tuần tự và tích lũy, trong đó tri thức thị trường nước ngoài
và cam kết nguồn lực vào thị trường đó tác động lẫn nhau theo vòng lặp
phản hồi. Doanh nghiệp bắt đầu bằng xuất khẩu gián tiếp, tiến đến đại lý
độc lập, sau đó là chi nhánh kinh doanh, và cuối cùng là sản xuất tại
nước ngoài, phản ánh mức độ cam kết và tri thức ngày càng tăng.

Một trong những khái niệm trọng tâm của mô hình Uppsala là "gánh nặng
của người nước ngoài" (liability of foreignness), chi phí và bất lợi mà
doanh nghiệp phải gánh chịu khi hoạt động ở môi trường văn hóa, thể chế,
và pháp lý xa lạ (Zaheer, 1995). Johanson và Vahlne (2009) cập nhật mô
hình bằng cách bổ sung khái niệm "gánh nặng của người ngoài mạng lưới"
(liability of outsidership), nhấn mạnh rằng trong nền kinh tế toàn cầu
hiện đại, rào cản lớn nhất không phải là khoảng cách địa lý hay tâm lý,
mà là vị trí ngoài lề của các mạng lưới kinh doanh địa phương. Hai loại
gánh nặng này giải thích vì sao chi phí quốc tế hóa giai đoạn đầu thường
cao, tạo ra giai đoạn hiệu quả giảm trước khi doanh nghiệp tích lũy đủ
tri thức và quan hệ để thu được lợi ích từ mở rộng quốc tế (Lu \&
Beamish, 2004; Contractor et al., 2003).

Tuy nhiên, mô hình Uppsala cổ điển được xây dựng trong bối cảnh tiền số
và tiền nền tảng (pre-platform), nên đã bị thách thức đáng kể trong thập
kỷ gần đây. Các nghiên cứu về born-global firms, tức doanh nghiệp quốc
tế hóa ngay từ giai đoạn đầu mà không theo trình tự tuần tự, và platform
firms, tức doanh nghiệp công nghệ tận dụng nền tảng số để vươn ra thị
trường quốc tế không cần hiện diện vật lý, cho thấy logic của Uppsala
không còn phổ quát (Stallkamp \& Schotter, 2021; Yang et al., 2025).
Trong luận án này, mô hình Uppsala không bị bác bỏ, mà được tái định vị
như một lớp nền giải thích chi phí học hỏi và bất định trong giai đoạn
đầu quốc tế hóa, còn cơ chế học hỏi được mở rộng để bao gồm "học hỏi
tăng cường bằng dữ liệu số" (data-augmented learning) trong kỷ nguyên
chuyển đổi số, nơi dữ liệu từ thị trường nước ngoài có thể được thu
thập, xử lý và phân tích nhanh hơn nhiều so với trước (Banalieva \&
Dhanaraj, 2019).

\hypertarget{222-luxfd-thuyux1ebft-dux1ef1a-truxean-nguux1ed3n-lux1ef1c}{%
\subsubsection{2.2.2 Lý thuyết dựa trên nguồn
lực}\label{222-luxfd-thuyux1ebft-dux1ef1a-truxean-nguux1ed3n-lux1ef1c}}

Lý thuyết dựa trên nguồn lực (Resource-Based View --- RBV) do Wernerfelt
(1984) đặt nền và được Barney (1991) hệ thống hóa thành một khung lý
thuyết có ảnh hưởng bền vững. RBV cho rằng doanh nghiệp không chỉ là tập
hợp các hoạt động giá trị (value chain) mà còn là bó nguồn lực (bundle
of resources), và sự khác biệt về hiệu quả lâu dài giữa các doanh nghiệp
bắt nguồn từ sự khác biệt trong sở hữu và khai thác các nguồn lực thỏa
mãn tiêu chí VRIN: có giá trị (Valuable), hiếm có (Rare), khó sao chép
(Inimitable), và khó thay thế (Non-substitutable).

Áp dụng vào bối cảnh quốc tế hóa, RBV giải thích vì sao cùng một mức độ
FSTS nhưng các doanh nghiệp lại đạt hiệu quả rất khác nhau: những doanh
nghiệp sở hữu năng lực công nghệ, năng lực quản trị vượt trội, và năng
lực học hỏi mạnh hơn sẽ chuyển hóa quốc tế hóa thành hiệu quả tốt hơn do
họ có thể giảm chi phí liability of foreignness và khai thác lợi thế
cạnh tranh trên phạm vi quốc tế hiệu quả hơn (Hitt et al., 2006; Peng,
2001). Cùng mức độ quốc tế hóa, doanh nghiệp có nguồn lực yếu hơn sẽ
gánh chịu chi phí điều phối và rủi ro lớn hơn, trong khi doanh nghiệp có
nguồn lực VRIN mạnh hơn có thể biến những thách thức đó thành lợi thế
cạnh tranh.

Trong bối cảnh đương đại, RBV được mở rộng bằng lý thuyết năng lực động
(dynamic capabilities) của Teece et al. (1997), nhấn mạnh khả năng của
doanh nghiệp trong việc tích hợp, xây dựng và tái cấu hình các năng lực
nội bộ để ứng phó với môi trường cạnh tranh thay đổi nhanh chóng. Trong
kỷ nguyên số, năng lực công nghệ (TCI), theo phân biệt của luận án, được
xem là một dạng năng lực động quan trọng, cho phép doanh nghiệp tái cấu
hình nguồn lực và mô hình kinh doanh để ứng phó với bất định trong quốc
tế hóa (Bhandari et al., 2023; Verhoef et al., 2021). DAI, chỉ số áp
dụng số nền tảng (Tier-1 proxy) đo lường sự hiện diện website, phản ánh
mức độ chấp nhận kỹ thuật số cơ bản chứ không phải năng lực động theo
nghĩa Teece et al. (1997). RBV cung cấp nền lý thuyết trực tiếp cho giả
thuyết H2 về vai trò điều tiết của TCI, trong khi H3 về DAI được neo đậu
trong Foundational Digital Adoption Framework (Banalieva \& Dhanaraj,
2019).

\hypertarget{223-luxfd-thuyux1ebft-thux1ec3-chux1ebf}{%
\subsubsection{2.2.3 Lý thuyết thể
chế}\label{223-luxfd-thuyux1ebft-thux1ec3-chux1ebf}}

Lý thuyết thể chế (Institutional Theory) trong kinh doanh quốc tế rút từ
nhiều nguồn lý luận: kinh tế học thể chế mới của North (1990), xã hội
học thể chế của Scott (1995), và lý thuyết về khoảng trống thể chế
(institutional voids) của Khanna và Palepu (2010). Điểm chung của các
tiếp cận này là nhấn mạnh vai trò của môi trường thể chế, cả chính thức
lẫn phi chính thức, trong việc định hình chiến lược và kết quả kinh
doanh của doanh nghiệp.

North (1990) chỉ ra rằng thể chế hiệu quả làm giảm chi phí giao dịch,
bảo vệ quyền sở hữu trí tuệ, và tạo ra môi trường khuyến khích đầu tư
dài hạn. Ngược lại, ở các nền kinh tế đang phát triển, sự vắng mặt hoặc
yếu kém của các thể chế trung gian thị trường, từ hệ thống pháp lý đến
thị trường vốn, từ tổ chức đánh giá tín nhiệm đến hiệp hội ngành nghề,
tạo ra "khoảng trống thể chế" (institutional voids) làm tăng chi phí
giao dịch và rủi ro cho doanh nghiệp khi hoạt động hoặc mở rộng ra nước
ngoài (Khanna \& Palepu, 2010).

Trong literature về I--P, bằng chứng meta-analytic của Marano et al.
(2016), phân tích 170 nghiên cứu với 32 quốc gia giai đoạn 1972--2012,
cho thấy thể chế quốc gia nguồn điều tiết đáng kể mối quan hệ quốc tế
hóa, hiệu quả, với doanh nghiệp từ các nền kinh tế có thể chế yếu hơn có
xu hướng thu được lợi ích thấp hơn hoặc thậm chí tổn thất từ quốc tế
hóa. Cuervo-Cazurra et al. (2018) bổ sung rằng bất định thể chế ở quốc
gia nguồn (home country uncertainty) không chỉ tạo ra thách thức mà còn
có thể xây dựng năng lực quản lý bất định đặc thù, giúp doanh nghiệp ở
các thị trường mới nổi cạnh tranh hiệu quả hơn ở các thị trường nước
ngoài có đặc điểm tương tự. Lý thuyết thể chế cung cấp nền tảng lý luận
cho giả thuyết H5 và cho khung ICRV trong luận án.

\hypertarget{224-luxfd-thuyux1ebft-thux1b0ux1ee3ng-tux1ea7ng-quux1ea3n-trux1ecb}{%
\subsubsection{2.2.4 Lý thuyết thượng tầng quản
trị}\label{224-luxfd-thuyux1ebft-thux1b0ux1ee3ng-tux1ea7ng-quux1ea3n-trux1ecb}}

Lý thuyết thượng tầng quản trị (Upper Echelons Theory) do Hambrick và
Mason (1984) đề xuất và được Hambrick (2007) cập nhật, cho rằng các
quyết định chiến lược lớn và kết quả của tổ chức không phải là sản phẩm
của các quy trình duy lý hoàn toàn, mà phản ánh đặc điểm nhận thức, giá
trị và kinh nghiệm của nhà quản trị cấp cao. Nói cách khác, "tổ chức là
sự phản chiếu của nhà quản trị đỉnh cao" (Hambrick \& Mason, 1984, tr.
193).

Trong bối cảnh quốc tế hóa, Upper Echelons Theory có hai hàm ý trực
tiếp. Thứ nhất, kinh nghiệm quốc tế của nhà quản trị cấp cao, được đo
bằng số năm kinh nghiệm làm việc hoặc học tập ở nước ngoài, hay tiếp xúc
với thị trường quốc tế, giúp giảm nhận thức về liability of foreignness
và nâng cao khả năng đánh giá cơ hội, quản trị rủi ro ở thị trường quốc
tế (Hsu et al., 2013; Sambharya, 1996). Thứ hai, sự đa dạng giới tính
trong lãnh đạo, đặc biệt là sự hiện diện của quản trị viên nữ ở các vị
trí lãnh đạo, liên quan đến nhận thức chiến lược đa chiều hơn và thực
tiễn quản trị rủi ro thận trọng hơn, qua đó có tác động tích cực đến
hiệu quả doanh nghiệp trong quá trình quốc tế hóa (Post \& Byron, 2015).
Những hàm ý này cung cấp cơ sở lý thuyết cho giả thuyết H4 về vai trò
điều tiết của đặc điểm nhà quản trị cấp cao.

\hypertarget{225-lux1edbp-mux1edf-rux1ed9ng-digital-capability-lens}{%
\subsubsection{2.2.5 Lớp mở rộng: Digital Capability
Lens}\label{225-lux1edbp-mux1edf-rux1ed9ng-digital-capability-lens}}

Bốn lý thuyết nền trên được xây dựng chủ yếu trong bối cảnh tiền số
(pre-digital). Sự phát triển mạnh mẽ của kinh tế số và trí tuệ nhân tạo
từ cuối thập niên 2010 đến nay đặt ra yêu cầu cập nhật và mở rộng các lý
thuyết kinh doanh quốc tế để phản ánh logic mới của quốc tế hóa trong kỷ
nguyên số.

Banalieva và Dhanaraj (2019) là một trong những nghiên cứu đầu tiên đề
xuất cập nhật lý thuyết nội hóa (internalization theory) cho nền kinh tế
số, lập luận rằng dữ liệu số và tương tác dựa trên nền tảng làm giảm sự
phụ thuộc vào hiện diện vật lý và thay đổi cơ bản logic về quyết định
nội hóa. Stallkamp và Schotter (2021) mở rộng luận điểm này bằng cách
chỉ ra rằng các doanh nghiệp nền tảng (platform firms) quốc tế hóa theo
logic đa diện (multi-sided) khác hẳn với các MNCs truyền thống, với hiệu
ứng mạng lưới tạo ra lợi thế quy mô không bị ràng buộc bởi biên giới địa
lý. Verhoef et al. (2021) cung cấp khung khái niệm phân tầng ba cấp độ:
số hóa thông tin (digitization), số hóa quy trình (digitalization), và
chuyển đổi số toàn diện (digital transformation), nhấn mạnh rằng các cấp
độ này khác nhau về chiều sâu tổ chức và tác động kinh doanh. Yang et
al. (2025) bổ sung bằng systematic review về quốc tế hóa của các doanh
nghiệp số, phác thảo các lý thuyết cần điều chỉnh và hướng nghiên cứu
tiếp theo.

Trong luận án, lớp mở rộng về năng lực số đóng vai trò mở rộng RBV bằng
cách xác định TCI và DAI như hai loại tài sản số có chiều sâu và cơ chế
vận hành khác nhau, đồng thời mở rộng Uppsala bằng cơ chế học hỏi dựa
trên dữ liệu số. Sự tích hợp này tạo thành lớp thứ năm của khung lý
thuyết, bổ sung cho bốn lý thuyết kinh điển để tạo thành một khung giải
thích hoàn chỉnh và phù hợp với thực tiễn kinh doanh đương đại.

\begin{center}\rule{0.5\linewidth}{0.5pt}\end{center}

\hypertarget{23-tux1ed5ng-quan-thux1ef1c-nghiux1ec7m}{%
\subsection{2.3 Tổng quan thực
nghiệm}\label{23-tux1ed5ng-quan-thux1ef1c-nghiux1ec7m}}

\hypertarget{231-giai-ux111oux1ea1n-phuxe1t-triux1ec3n-cux1ee7a-literature-ip}{%
\subsubsection{2.3.1 Giai đoạn phát triển của literature
I--P}\label{231-giai-ux111oux1ea1n-phuxe1t-triux1ec3n-cux1ee7a-literature-ip}}

Nghiên cứu về mối quan hệ giữa quốc tế hóa và hiệu quả hoạt động doanh
nghiệp (I--P relationship) đã trải qua hơn bốn thập kỷ phát triển, với
những thay đổi đáng kể về cả câu hỏi nghiên cứu, phương pháp luận, và
kết luận lý thuyết.

Giai đoạn đầu (từ thập niên 1980 đến giữa thập niên 1990) đặc trưng bởi
giả định về quan hệ tuyến tính dương giữa quốc tế hóa và hiệu quả. Grant
(1987) và Kim et al. (1989) cho rằng quốc tế hóa giúp doanh nghiệp khai
thác lợi thế theo quy mô (economies of scale), lợi thế phạm vi
(economies of scope), và phân tán rủi ro kinh tế vĩ mô, do đó dẫn đến
hiệu quả cao hơn. Logic này trực quan và phù hợp với quan sát ban đầu về
các MNCs thành công từ Mỹ, Châu Âu và Nhật Bản.

Giai đoạn thứ hai (từ cuối thập niên 1990 đến đầu thập niên 2000) chứng
kiến sự xuất hiện của mô hình phi tuyến, thách thức giả định tuyến tính
giản đơn. Hitt et al. (1997) là nghiên cứu đầu tiên kiểm định và tìm
thấy bằng chứng cho quan hệ chữ U ngược (inverted-U) giữa đa dạng hóa
sản phẩm quốc tế và hiệu quả, lập luận rằng lợi ích từ quốc tế hóa tăng
đến một điểm tối ưu rồi suy giảm do chi phí điều phối và phức tạp tổ
chức vượt qua lợi ích. Gomes và Ramaswamy (1999) cung cấp bằng chứng
tương tự trong bối cảnh doanh nghiệp đa ngành. Contractor et al. (2003)
và Lu và Beamish (2004) tinh chỉnh mô hình phi tuyến với đề xuất "đường
cong chữ S ba giai đoạn" (three-stage S-curve), theo đó hiệu quả giảm ở
giai đoạn thâm nhập ban đầu (chi phí thiết lập), tăng ở giai đoạn khai
thác (học hỏi tích lũy và quy mô), rồi lại giảm ở giai đoạn siêu mở rộng
(over-extension) khi chi phí phức tạp vượt quá lợi ích.

Giai đoạn thứ ba (từ cuối thập niên 2000 đến hiện tại) chuyển trọng tâm
sang nghiên cứu các yếu tố điều tiết (moderators), tức điều kiện bối
cảnh làm thay đổi hình dạng, cường độ, hoặc dấu của mối quan hệ I--P.
Các moderator được nghiên cứu bao gồm đặc điểm quốc gia nguồn (Marano et
al., 2016), đặc điểm ngành (Kirca et al., 2012), năng lực lãnh đạo cấp
cao (Hsu et al., 2013), và năng lực kỹ thuật số (Bhandari et al., 2023;
Wu et al., 2022). Sự chuyển dịch này phản ánh nhận thức ngày càng sâu
sắc rằng không có quan hệ I--P phổ quát, kết quả quốc tế hóa phụ thuộc
vào điều kiện cụ thể của doanh nghiệp, ngành và quốc gia (Arte \&
Larimo, 2022).

\hypertarget{232-bux1eb1ng-chux1ee9ng-meta-analytic}{%
\subsubsection{2.3.2 Bằng chứng
meta-analytic}\label{232-bux1eb1ng-chux1ee9ng-meta-analytic}}

Một trong những đóng góp quan trọng nhất trong giai đoạn gần đây là sự
phát triển của các meta-analyses về I--P relationship, cho phép tổng hợp
và định lượng hóa kết quả từ hàng trăm nghiên cứu thực nghiệm riêng lẻ.
Tổng hợp từ các meta-analyses chính cho thấy một bức tranh nhất quán
nhưng phức tạp.

Bausch và Krist (2007) phân tích 65 nghiên cứu và tìm thấy tác động tổng
hợp dương nhưng nhỏ (r ≈ 0,10), đồng thời xác nhận rằng mức độ dị biệt
(heterogeneity) giữa các nghiên cứu rất cao với I² \textgreater{} 80\%,
hàm ý rằng các moderator là nguồn giải thích chủ yếu cho sự khác biệt.
Kirca et al. (2012) mở rộng pool lên 154 mẫu và tìm thấy tác động tổng
hợp tương tự, đồng thời xác nhận vai trò của bối cảnh quốc gia và ngành.
Marano et al. (2016), nghiên cứu meta-analytic toàn diện nhất trong
literature với 170 nghiên cứu và 32 quốc gia, tìm thấy r ≈ 0,07--0,13 và
chứng minh rằng thể chế quốc gia nguồn là moderator quan trọng nhất
trong tất cả các biến bối cảnh được kiểm định. Wu et al. (2022) cập nhật
với pool 218 effect sizes tập trung vào các emerging market
multinationals (EMNEs) và tìm thấy mô hình moderator tương tự. Tại hội
nghị ICBEF 2024, kết quả sơ bộ từ pool nghiên cứu về châu Á với k = 113
effect-size units cho thấy r = 0,07, nhất quán với các meta-analyses
toàn cầu nhưng cụ thể hơn về bối cảnh châu Á. Luận án này mở rộng pool
đó thông qua Nghiên cứu P6, một meta-analysis ba tầng (three-level MARA)
với k = 238 nghiên cứu và K = 288 effect sizes từ châu Á và Thái Bình
Dương, ước lượng r̄ = 0,074 {[}KTC 95\%: 0,060; 0,088{]}, I² = 62,4\%,
Q\_M(ICRV) = 17,35 (df = 4, p = ,002). Kết quả này xác nhận tác động
tổng hợp dương nhỏ nhưng bền vững, đồng thời chứng minh rằng gradient
ICRV, turning point giảm dần khi chuyển từ thể chế yếu sang thể chế
mạnh, là moderator thực chất của mối quan hệ I--P trong bối cảnh châu Á
(Do \& Phan, 2026c).

Nhìn chung, bằng chứng meta-analytic cho thấy ba kết luận quan trọng:
tác động tổng hợp của quốc tế hóa lên hiệu quả là dương nhưng khiêm tốn;
mức dị biệt rất cao; và các moderator bối cảnh, đặc biệt là thể chế quốc
gia, năng lực doanh nghiệp, và đặc điểm ngành, là những nguồn giải thích
chính cho sự khác biệt giữa các nghiên cứu. Đây là luận điểm cơ bản biện
hộ cho sự cần thiết của một khung giải thích đa tầng.

\hypertarget{233-bux1ed1i-cux1ea3nh-chuxe2u-uxe1-vuxe0-cuxe1c-nux1ec1n-kinh-tux1ebf-mux1edbi-nux1ed5i}{%
\subsubsection{2.3.3 Bối cảnh châu Á và các nền kinh tế mới
nổi}\label{233-bux1ed1i-cux1ea3nh-chuxe2u-uxe1-vuxe0-cuxe1c-nux1ec1n-kinh-tux1ebf-mux1edbi-nux1ed5i}}

Châu Á là bối cảnh đặc biệt quan trọng và đặc thù trong literature I--P
vì nhiều lý do. Thứ nhất, sự đa dạng thể chế trong khu vực là ngoại lệ
không có ở bất kỳ khu vực nào khác trên thế giới: từ Singapore với chất
lượng quản trị đứng hàng đầu thế giới đến Afghanistan với chất lượng
quản trị rất yếu, từ Nhật Bản với nền kinh tế phát triển hoàn thiện đến
các Pacific SIDS với quy mô siêu nhỏ và hạ tầng thể chế còn sơ khai
(Khanna \& Palepu, 2010; Peng, 2003). Thứ hai, tốc độ chuyển đổi số và
hội nhập quốc tế rất không đồng đều giữa các quốc gia trong khu vực, tạo
ra "thí nghiệm tự nhiên" về ảnh hưởng của bối cảnh thể chế và năng lực
số đến kết quả quốc tế hóa.

Do và Phan (2026a) là nghiên cứu thực nghiệm multi-country quy mô lớn
nhất hiện có về I--P trong bối cảnh châu Á mới nổi, phân tích 40.633
doanh nghiệp từ 17 nền kinh tế châu Á sử dụng dữ liệu WBES. Nghiên cứu
tìm thấy: (i) tồn tại phân tầng năng suất rõ rệt giữa nhóm Chinese
economies, South Asia và ASEAN, phản ánh sự dị biệt thể chế và năng lực
giữa các nhóm; (ii) việc áp dụng công nghệ vừa làm tăng năng suất lao
động trực tiếp vừa làm suy giảm đáng kể tác động bất lợi của các rào cản
thể chế, phát hiện được gọi là "hiệu ứng lá chắn số" (digital shield
effect). Do và Phan (2026g) kiểm định cụ thể với dữ liệu doanh nghiệp
Trung Quốc và tìm thấy cấu trúc cubic inverted-U với ngưỡng FSTS tối ưu
khoảng 47,8\%, tức điểm turning point khá thấp so với các nền kinh tế
tiên tiến. Luận án kế thừa và mở rộng các kết quả này trên pool đầy đủ
49 nền kinh tế (bao gồm 9 nền kinh tế đảo nhỏ Thái Bình Dương).

Phần lớn các nghiên cứu về I--P trong bối cảnh châu Á tập trung vào Nhật
Bản, Hàn Quốc, Đài Loan, và Trung Quốc, trong khi các nền kinh tế chuyển
đổi như Việt Nam, Campuchia, Bangladesh, và đặc biệt là các Pacific
SIDS, về cơ bản thiếu vắng trong literature. Sự thiếu đại diện này tạo
ra khoảng trống nghiêm trọng trong khả năng khái quát hóa lý thuyết, vì
các nền kinh tế frontier và SIDS có thể là địa điểm quan trọng để kiểm
định điều kiện biên (boundary conditions) của các lý thuyết về I--P.

\hypertarget{234-vai-truxf2-cux1ee7a-nux103ng-lux1ef1c-sux1ed1-trong-mux1ed1i-quan-hux1ec7-ip}{%
\subsubsection{2.3.4 Vai trò của năng lực số trong mối quan hệ
I--P}\label{234-vai-truxf2-cux1ee7a-nux103ng-lux1ef1c-sux1ed1-trong-mux1ed1i-quan-hux1ec7-ip}}

Sự tích hợp giữa nghiên cứu về năng lực số và I--P relationship là một
xu hướng tương đối mới, phát triển mạnh từ 2019 đến nay. Bhandari et al.
(2023) phân tích 571 doanh nghiệp sản xuất của Mỹ và tìm thấy rằng năng
lực kỹ thuật số (được đo như một composite duy nhất) điều tiết tích cực
mối quan hệ I--P, với doanh nghiệp có năng lực số cao hơn thu được lợi
ích lớn hơn từ quốc tế hóa. Bustamante et al. (2022) sử dụng mẫu SMEs đa
quốc gia và phát hiện rằng công cụ số làm giảm chi phí thâm nhập thị
trường nước ngoài, đặc biệt ở các thị trường có khoảng trống thể chế
lớn. Chen và Meng (2022) nghiên cứu doanh nghiệp Trung Quốc từ WBES và
phát hiện rằng ràng buộc thể chế làm giảm cường độ xuất khẩu, nhưng
doanh nghiệp có năng lực số cao hơn có thể vượt qua ràng buộc này hiệu
quả hơn.

Tuy nhiên, phần lớn các nghiên cứu hiện tại không phân biệt TCI và DAI,
xử lý năng lực số như một khái niệm đơn nhất. Hậu quả lý thuyết của sự
đồng nhất không hợp lý này là không thể phân tách được hai cơ chế vận
hành khác nhau: chiều sâu năng lực tổ chức (TCI) và khả năng phối hợp kỹ
thuật số bề mặt (DAI). Sự không phân biệt này làm mờ đi cơ chế tác động
và dẫn đến kết luận không thể áp dụng trực tiếp vào thực tiễn chính sách
và quản trị doanh nghiệp.

\begin{center}\rule{0.5\linewidth}{0.5pt}\end{center}

\hypertarget{24-khoux1ea3ng-trux1ed1ng-nghiuxean-cux1ee9u}{%
\subsection{2.4 Khoảng trống nghiên
cứu}\label{24-khoux1ea3ng-trux1ed1ng-nghiuxean-cux1ee9u}}

Trên cơ sở tổng quan lý thuyết và thực nghiệm ở các phần trước, ba
khoảng trống nghiên cứu chính được xác định như sau.

\textbf{Khoảng trống thứ nhất: Thiếu khung tích hợp đa tầng để giải
thích heterogeneity trong I--P relationship.} Các meta-analyses nhất
quán cho thấy I² \textgreater{} 80\%, tức hơn 80\% biến thiên trong
effect size không được giải thích bởi sai số ngẫu nhiên mà bởi sự khác
biệt có hệ thống giữa các bối cảnh (Bausch \& Krist, 2007; Marano et
al., 2016). Tuy nhiên, hầu hết các nghiên cứu thực nghiệm hiện tại chỉ
kiểm định một hoặc hai moderator từ một lý thuyết đơn lẻ, không có khả
năng giải thích đồng thời heterogeneity do bối cảnh thể chế, năng lực
doanh nghiệp, và đặc điểm lãnh đạo tạo ra. Không có lý thuyết đơn lẻ
nào, dù là Uppsala, RBV, Institutional Theory, hay Upper Echelons
Theory, đủ để giải thích đầy đủ sự phức tạp của I--P relationship trong
bối cảnh đa quốc gia đa dạng như châu Á. Khoảng trống này đòi hỏi một
khung tích hợp đa tầng, trong đó các lý thuyết bổ sung cho nhau thay vì
cạnh tranh nhau.

\textbf{Khoảng trống thứ hai: Sự đồng nhất không hợp lý giữa TCI và DAI
trong phần lớn literature trước.} Phần lớn các nghiên cứu về năng lực số
trong bối cảnh I--P hoặc là không đo lường năng lực số, hoặc là gộp tất
cả các chỉ báo số vào một composite duy nhất, không phân biệt giữa chiều
sâu năng lực công nghệ tổ chức (TCI) và mức độ chấp nhận công cụ số bề
mặt (DAI). Hậu quả thực tiễn của sự đồng nhất này là: (i) không thể trả
lời câu hỏi liệu nên đầu tư vào năng lực công nghệ nội tại hay vào việc
áp dụng công cụ số giao diện; (ii) không thể phân biệt được hai cơ chế
tác động có hàm ý chính sách khác nhau. Phân biệt rõ TCI và DAI là điều
kiện cần để rút ra hàm ý quản trị và chính sách có tính cụ thể và ứng
dụng cao.

\textbf{Khoảng trống thứ ba: Thiếu bằng chứng cross-country quy mô lớn,
toàn diện và có cấu trúc cho bối cảnh châu Á.} Mặc dù châu Á là khu vực
kinh tế năng động và đa dạng nhất thế giới, literature về I--P vẫn thiếu
bằng chứng đa quốc gia bao quát đủ rộng để kiểm định khả năng tổng quát
hóa của các lý thuyết. Pool 17 nước của Do và Phan (2026a) là baseline
lớn nhất hiện có trong literature về châu Á, nhưng vẫn chưa bao gồm các
nền kinh tế frontier và Pacific SIDS, những địa điểm rất quan trọng để
kiểm định boundary conditions. Thiếu vắng các bằng chứng này khiến cho
các lý thuyết về I--P, vốn được xây dựng chủ yếu trên dữ liệu từ Nhật
Bản, Hàn Quốc, và các nền kinh tế tiên tiến, có nguy cơ bị áp dụng sai
trong bối cảnh các nền kinh tế đang phát triển và chuyển đổi ở châu Á.

Ba khoảng trống này hội tụ vào một hướng giải quyết chung: cần phát
triển và kiểm định một khung giải thích tích hợp đa tầng (multi-tier
integrated framework) trên dữ liệu đa quốc gia đủ rộng và đủ đa dạng về
thể chế để bao phủ toàn bộ gradient thể chế của châu Á. Luận án đề xuất
khung CDCM (Context-Contingent Digital-and-Capability Model), là khung
lý thuyết tích hợp đa tầng kết hợp bối cảnh thể chế (tầng vĩ mô, vận
hành hóa bằng ICRV), năng lực doanh nghiệp với TCI và DAI riêng biệt
(tầng tổ chức), đặc điểm nhà quản trị cấp cao (tầng cá nhân), và lớp mở
rộng digital capability lens. Trong khung này, FSTS là biến độc lập
chính và hiệu quả doanh nghiệp (labor productivity) là biến phụ thuộc.

\textbf{Từ khoảng trống đến Nghiên cứu P7 --- Kiểm định CDCM trên 49 nền
kinh tế châu Á.} Khoảng trống thứ nhất và thứ ba hội tụ vào một câu hỏi
kiểm định trực tiếp: liệu CDCM có thể giải thích đồng thời heterogeneity
I--P khi kiểm soát đầy đủ cả ba tầng trong một mô hình tích hợp ở quy mô
đủ rộng để bao phủ toàn bộ gradient ICRV không? Literature hiện tại chưa
có bằng chứng kiểm định đồng thời năm moderator chính (TCI, DAI, kinh
nghiệm quản trị, giới tính nhà quản trị, ICRV) trong một mô hình hồi quy
đơn lẻ trên dữ liệu đa quốc gia quy mô lớn. Do và Phan (2026d) ---
Nghiên cứu P7, lấp đầy khoảng trống này bằng dữ liệu WBES từ 49 nền kinh
tế (N = 84.910--91.982 doanh nghiệp, 102 country-year waves), bao phủ
toàn bộ gradient ICRV từ nhóm tiên tiến đến SIDS, kiểm định đồng thời
H1--H6 với country-year fixed effects. Kết quả turning point trung bình
\textasciitilde36\% FSTS tại P7, thấp hơn P3 Vietnam
(\textasciitilde40\%) và P5 Trung Quốc (\textasciitilde49\%), xác nhận
ICRV điều tiết vị trí turning point theo gradient dự đoán lý thuyết,
đóng góp bằng chứng kiểm định tích hợp đầu tiên cho CDCM ở cấp độ liên
quốc gia quy mô lớn.

\textbf{Từ khoảng trống đến Nghiên cứu P8 --- Kiểm định điều kiện biên
FIP tại Pacific SIDS.} Khoảng trống thứ ba không chỉ đòi hỏi mở rộng
phạm vi địa lý mà còn đòi hỏi kiểm định điều kiện biên lý thuyết cực
trị: liệu có tồn tại bối cảnh nơi inverted-U hoàn toàn sụp đổ thành quan
hệ âm đơn điệu? Không có nghiên cứu nào trong literature I--P hiện tại
kiểm định "Forced Internationalization Penalty" (FIP), bối cảnh mà quốc
tế hóa là gánh nặng sinh tồn, không phải lựa chọn chiến lược dựa trên
lợi thế cạnh tranh. Do và Phan (2026b) --- Nghiên cứu P8, kiểm định H1b
bằng dữ liệu WBES từ các nền kinh tế Pacific SIDS (Fiji, Tonga, Solomon
Islands, Vanuatu, Samoa), nơi ba điều kiện cấu trúc cơ bản của
inverted-U đồng thời bị vi phạm: thị trường nội địa cực nhỏ, chi phí
logistics cao bất thường, và hỗ trợ thể chế cho xuất khẩu gần như không
tồn tại. Kết quả P8 xác định boundary condition cực trị của CDCM: điểm
nơi lý thuyết đường cong chữ U ngược không còn ứng dụng được.

\begin{center}\rule{0.5\linewidth}{0.5pt}\end{center}

\hypertarget{25-muxf4-huxecnh-nghiuxean-cux1ee9u-vuxe0-hux1ec7-giux1ea3-thuyux1ebft}{%
\subsection{2.5 Mô hình nghiên cứu và hệ giả
thuyết}\label{25-muxf4-huxecnh-nghiuxean-cux1ee9u-vuxe0-hux1ec7-giux1ea3-thuyux1ebft}}

\hypertarget{251-muxf4-huxecnh-khuxe1i-niux1ec7m}{%
\subsubsection{2.5.1 Mô hình khái
niệm}\label{251-muxf4-huxecnh-khuxe1i-niux1ec7m}}

Mô hình khái niệm của luận án vận hành theo logic đa tầng, trong đó mối
quan hệ trung tâm giữa quốc tế hóa (FSTS) và hiệu quả hoạt động doanh
nghiệp (labor productivity) được điều tiết đồng thời bởi ba tầng can
thiệp.

Tầng vĩ mô, bối cảnh thể chế, được vận hành hóa bằng khung ICRV với sáu
sub-regime theo gradient từ Advanced Innovation-Driven đến Pacific SIDS.
Ở tầng này, ICRV điều tiết mối quan hệ I--P thông qua hai cơ chế: (i)
mức độ khoảng trống thể chế (institutional voids) ảnh hưởng đến chi phí
giao dịch xuyên biên giới; và (ii) chất lượng thể chế ảnh hưởng đến khả
năng doanh nghiệp khai thác lợi ích từ quốc tế hóa. Ở các nền kinh tế
frontier và SIDS, khoảng trống thể chế rất lớn có thể đảo ngược quan hệ
I--P, biến lợi ích tiềm năng thành tổn thất thực tế (Do \& Phan, 2026b).

Tầng tổ chức, năng lực doanh nghiệp, bao gồm TCI và DAI như hai
moderator độc lập với cơ chế tác động khác nhau. TCI đóng vai trò như
một "nâng đỡ mặt bằng" (level shifter): doanh nghiệp có TCI cao sẽ đạt
hiệu quả cao hơn ở mọi mức độ FSTS, và đặc biệt có khả năng khai thác
lợi ích quốc tế hóa tốt hơn nhờ năng lực hấp thụ tri thức và công nghệ
từ thị trường nước ngoài. DAI, ngược lại, đóng vai trò là "bộ đệm phối
hợp số" (digital coordination buffer): ở mức FSTS cao, DAI giúp giảm chi
phí điều phối xuyên biên giới, làm chuyển dịch điểm turning point sang
phải và kéo dài giai đoạn lợi ích dương của quan hệ I--P.

Tầng cá nhân, đặc điểm nhà quản trị cấp cao, bao gồm kinh nghiệm quản
trị và giới tính nhà quản trị như các moderator vi mô. Kinh nghiệm quốc
tế của nhà quản trị giảm chi phí nhận thức của liability of foreignness;
sự hiện diện của quản trị viên nữ tăng cường đa dạng nhận thức và thận
trọng trong quản trị rủi ro quốc tế.

Ngoài các cơ chế điều tiết, mô hình còn bao gồm một yếu tố động thái
quan trọng: hình dạng phi tuyến của quan hệ I--P (turning point của
inverted-U) không cố định theo thời gian mà có thể thay đổi khi bối cảnh
thể chế và năng lực doanh nghiệp thay đổi theo chiều dọc, đặc biệt rõ
rệt ở các nền kinh tế chuyển đổi nhanh như Trung Quốc.

\hypertarget{252-giux1ea3-thuyux1ebft-h1-phi-tuyux1ebfn-chux1eef-u-ngux1b0ux1ee3c}{%
\subsubsection{2.5.2 Giả thuyết H1: Phi tuyến chữ U
ngược}\label{252-giux1ea3-thuyux1ebft-h1-phi-tuyux1ebfn-chux1eef-u-ngux1b0ux1ee3c}}

Trên cơ sở lý thuyết từ Uppsala về chi phí thâm nhập giai đoạn đầu, RBV
về lợi thế cạnh tranh tích lũy ở giai đoạn trung, và lý thuyết chi phí
giao dịch về sự phức tạp điều phối ở giai đoạn mở rộng quá mức, giả
thuyết H1 được phát biểu như sau. Quan hệ giữa mức độ quốc tế hóa và
hiệu quả hoạt động của các doanh nghiệp ở các quốc gia châu Á có dạng
phi tuyến chữ U ngược, với một điểm turning point tối ưu trước khi chi
phí điều phối và bất định vượt qua lợi ích từ mở rộng thị trường quốc
tế. Cơ sở lý thuyết trực tiếp bao gồm mô hình S-curve của Lu và Beamish
(2004), lý thuyết ba giai đoạn của Contractor et al. (2003), và bằng
chứng inverted-U ban đầu của Hitt et al. (1997).

\hypertarget{252b-ux111iux1ec1u-kiux1ec7n-biuxean-h1b-guxe1nh-nux1eb7ng-quux1ed1c-tux1ebf-huxf3a-bux1eaft-buux1ed9c-fip-tux1ea1i-pacific-sids}{%
\subsubsection{2.5.2b Điều kiện biên H1b: Gánh nặng quốc tế hóa bắt buộc
(FIP) tại Pacific
SIDS}\label{252b-ux111iux1ec1u-kiux1ec7n-biuxean-h1b-guxe1nh-nux1eb7ng-quux1ed1c-tux1ebf-huxf3a-bux1eaft-buux1ed9c-fip-tux1ea1i-pacific-sids}}

H1 dự báo quan hệ chữ U ngược, tức là tồn tại một mức FSTS tối ưu mà
dưới ngưỡng đó quốc tế hóa tạo ra lợi ích, và trên ngưỡng đó chi phí
điều phối vượt qua lợi ích học hỏi. Tuy nhiên, lý thuyết dự đoán này
ngầm giả định rằng ba điều kiện cấu trúc cơ bản được đáp ứng: (i) tồn
tại thị trường nội địa đủ lớn để doanh nghiệp có nền tảng hiệu suất và
thanh khoản trước khi xuất khẩu; (ii) chi phí thương mại ở mức chấp nhận
được, cho phép cân nhắc kinh tế hợp lý; và (iii) tồn tại hỗ trợ thể chế
tối thiểu cho hoạt động trao đổi quốc tế.

Tại các nền kinh tế Pacific SIDS (Nhóm VI ICRV), cả ba điều kiện này
đồng thời bị vi phạm: thị trường nội địa cực nhỏ buộc các doanh nghiệp
phải xuất khẩu không phải vì có lợi thế cạnh tranh mà vì thiếu cầu nội
địa đủ lớn để duy trì hoạt động (xuất khẩu bắt buộc); chi phí hậu cần và
vận chuyển cực cao do địa lý đảo xa và phân tán làm suy giảm biên lợi
nhuận ở mọi mức FSTS; và hỗ trợ thể chế (logistics, tài chính thương
mại, năng lực xúc tiến xuất khẩu) thiếu hụt nghiêm trọng. Khi ba điều
kiện này đồng thời không được đáp ứng, quan hệ I→P không chỉ mất đi phần
lợi ích giai đoạn đầu của chữ U ngược mà có thể chuyển thành
\textbf{quan hệ âm đơn điệu không có điểm uốn}, tức là điều kiện biên lý
thuyết nơi mọi mức độ quốc tế hóa đều gây tổn hại đến hiệu quả doanh
nghiệp.

Giả thuyết H1b được phát biểu như sau. Tại các nền kinh tế Pacific SIDS
(ICRV Nhóm VI), mối quan hệ giữa mức độ quốc tế hóa (FSTS) và hiệu quả
hoạt động doanh nghiệp (năng suất lao động) là \textbf{âm đơn điệu},
không có điểm turning point trong phạm vi dữ liệu, phản ánh gánh nặng
quốc tế hóa bắt buộc (Forced Internationalization Penalty, FIP) nơi
doanh nghiệp xuất khẩu để sinh tồn chứ không phải để khai thác lợi thế
cạnh tranh. Cơ sở lý thuyết bao gồm Institutional Theory (North, 1990),
lý thuyết chi phí giao dịch (Williamson, 1985), Khanna và Palepu (2010)
về institutional voids, và Lall (2004) về SIDS structural constraints.

\hypertarget{253-giux1ea3-thuyux1ebft-h2-ux111iux1ec1u-tiux1ebft-bux1edfi-nux103ng-lux1ef1c-cuxf4ng-nghux1ec7-tci}{%
\subsubsection{2.5.3 Giả thuyết H2: Điều tiết bởi năng lực công nghệ
(TCI)}\label{253-giux1ea3-thuyux1ebft-h2-ux111iux1ec1u-tiux1ebft-bux1edfi-nux103ng-lux1ef1c-cuxf4ng-nghux1ec7-tci}}

Dựa trên RBV (Barney, 1991) và lý thuyết năng lực động (Teece et al.,
1997), năng lực công nghệ là một nguồn lực VRIN giúp doanh nghiệp khai
thác lợi ích từ quốc tế hóa hiệu quả hơn. Giả thuyết H2 được phát biểu
như sau. Năng lực công nghệ (TCI) làm gia tăng tác động tích cực của
quốc tế hóa lên hiệu quả hoạt động của doanh nghiệp, bằng cách nâng mặt
bằng hiệu quả tổng thể và làm giảm chi phí học hỏi thị trường quốc tế,
do đó làm dịch chuyển đường quan hệ I--P lên cao hơn và có thể điều
chỉnh vị trí của điểm turning point. Cơ sở lý thuyết bao gồm Barney
(1991), Teece et al. (1997), và Bhandari et al. (2023).

\hypertarget{254-giux1ea3-thuyux1ebft-h3-ux111iux1ec1u-tiux1ebft-bux1edfi-chux1ec9-sux1ed1-chux1ea5p-nhux1eadn-sux1ed1-dai}{%
\subsubsection{2.5.4 Giả thuyết H3: Điều tiết bởi chỉ số chấp nhận số
(DAI)}\label{254-giux1ea3-thuyux1ebft-h3-ux111iux1ec1u-tiux1ebft-bux1edfi-chux1ec9-sux1ed1-chux1ea5p-nhux1eadn-sux1ed1-dai}}

Phân biệt rõ với H2, DAI không có chiều sâu năng lực tổ chức như TCI
nhưng có vai trò giảm thiểu chi phí phối hợp xuyên biên giới ở mức xuất
khẩu cao. Giả thuyết H3 được phát biểu như sau. Chỉ số chấp nhận số
(DAI) điều tiết mối quan hệ I--P, với vai trò bổ trợ rõ rệt hơn ở mức
export intensity cao khi DAI giúp giảm chi phí phối hợp giao dịch xuyên
biên giới, theo đó DAI làm thay đổi độ dốc của đường quan hệ I--P chứ
không nhất thiết nâng mặt bằng hiệu quả như TCI. Cơ sở lý thuyết bao gồm
Verhoef et al. (2021), Stallkamp và Schotter (2021), và Bustamante et
al. (2022). Sự phân biệt cơ chế tác động giữa H2 và H3 là đóng góp lý
thuyết quan trọng, cho phép phân tách TCI và DAI trong mô hình hồi quy
thực nghiệm.

\hypertarget{255-giux1ea3-thuyux1ebft-h4-ux111iux1ec1u-tiux1ebft-bux1edfi-ux111ux1eb7c-ux111iux1ec3m-nhuxe0-quux1ea3n-trux1ecb-cux1ea5p-cao}{%
\subsubsection{2.5.5 Giả thuyết H4: Điều tiết bởi đặc điểm nhà quản trị
cấp
cao}\label{255-giux1ea3-thuyux1ebft-h4-ux111iux1ec1u-tiux1ebft-bux1edfi-ux111ux1eb7c-ux111iux1ec3m-nhuxe0-quux1ea3n-trux1ecb-cux1ea5p-cao}}

Dựa trên Upper Echelons Theory (Hambrick \& Mason, 1984; Hambrick,
2007), quyết định chiến lược quốc tế hóa và khả năng chuyển hóa nó thành
hiệu quả phụ thuộc vào đặc điểm nhận thức và kinh nghiệm của nhà quản
trị. Giả thuyết H4 được phát biểu như sau. Kinh nghiệm quản trị và giới
tính nữ của nhà quản trị cấp cao làm mạnh hơn tác động tích cực của quốc
tế hóa lên hiệu quả hoạt động, thông qua cơ chế nâng cao nhận thức chiến
lược quốc tế và quản trị rủi ro thận trọng hơn trong quá trình mở rộng
quốc tế. Cơ sở lý thuyết bao gồm Hambrick và Mason (1984), Hsu et al.
(2013), và Post và Byron (2015).

\hypertarget{256-giux1ea3-thuyux1ebft-h5-ux111iux1ec1u-tiux1ebft-bux1edfi-bux1ed1i-cux1ea3nh-thux1ec3-chux1ebf-icrv}{%
\subsubsection{2.5.6 Giả thuyết H5: Điều tiết bởi bối cảnh thể chế
(ICRV)}\label{256-giux1ea3-thuyux1ebft-h5-ux111iux1ec1u-tiux1ebft-bux1edfi-bux1ed1i-cux1ea3nh-thux1ec3-chux1ebf-icrv}}

Trên cơ sở Institutional Theory (North, 1990; Khanna \& Palepu, 2010),
chất lượng thể chế tạo ra chi phí giao dịch và rủi ro khác nhau, do đó
điều tiết đáng kể mối quan hệ I--P. Giả thuyết H5 được phát biểu như
sau. Chất lượng thể chế theo khung ICRV điều tiết tác động của các rào
cản thể chế (institutional obstacles) lên hiệu quả hoạt động trong quá
trình quốc tế hóa, với chỉ số chấp nhận số (DAI) đóng vai trò "lá chắn
số" (digital shield) giúp bù đắp một phần khoảng trống thể chế; cường độ
điều tiết này giảm dần theo gradient ICRV từ các nền kinh tế frontier
(Nhóm V--VI, điều tiết mạnh nhất) đến các nền kinh tế tiên tiến (Nhóm I,
điều tiết yếu nhất). Cơ sở lý thuyết bao gồm North (1990), Khanna và
Palepu (2010), và Marano et al. (2016).

\hypertarget{257-giux1ea3-thuyux1ebft-h6-dux1ecb-biux1ec7t-thux1eddi-gian-trong-quan-hux1ec7-ip}{%
\subsubsection{2.5.7 Giả thuyết H6: Dị biệt thời gian trong quan hệ
I--P}\label{257-giux1ea3-thuyux1ebft-h6-dux1ecb-biux1ec7t-thux1eddi-gian-trong-quan-hux1ec7-ip}}

Dựa trên quan sát về tốc độ chuyển đổi số và thay đổi thể chế rất nhanh
ở một số nền kinh tế châu Á, đặc biệt Trung Quốc, quan hệ I--P không
nhất thiết ổn định theo thời gian. Giả thuyết H6 được phát biểu như sau.
Hình dạng phi tuyến của quan hệ I--P, cụ thể là vị trí của điểm turning
point và độ cong của đường quan hệ, thay đổi theo thời gian trong các
nền kinh tế chuyển đổi nhanh, phản ánh sự phát triển của năng lực doanh
nghiệp và cải thiện điều kiện thể chế qua các giai đoạn. Cơ sở lý thuyết
bao gồm Wu et al. (2022), Xiao et al. (2013), và Banalieva và Dhanaraj
(2019). Bằng chứng từ Do và Phan (2026g) về sự ổn định tương đối của
turning point Trung Quốc giữa 2012 và 2024 (Paternoster z = +0,82, p =
.412, không bác bỏ bằng nhau) cung cấp cơ sở thực nghiệm bổ sung cho
việc kiểm định H6 ở cấp độ luận án.

\begin{center}\rule{0.5\linewidth}{0.5pt}\end{center}

\hypertarget{26-tuxf3m-tux1eaft-chux1b0ux1a1ng}{%
\subsection{2.6 Tóm tắt
chương}\label{26-tuxf3m-tux1eaft-chux1b0ux1a1ng}}

Chương 2 đã trình bày hệ thống nền tảng lý thuyết và thực nghiệm cho
luận án theo một trục logic nhất quán: từ khái niệm đến lý thuyết, từ
bằng chứng thực nghiệm đến khoảng trống, và từ khoảng trống đến khung
nghiên cứu và giả thuyết.

Phần đầu chương xác lập bốn khái niệm cốt lõi: quốc tế hóa doanh nghiệp
(đo bằng FSTS), hiệu quả hoạt động (đo bằng labor productivity là
chính), năng lực công nghệ (TCI) và chỉ số chấp nhận số (DAI) như hai
construct độc lập, và bối cảnh thể chế được vận hành hóa bằng khung ICRV
sáu nhóm. Phần lý thuyết nền trình bày và phân tích bốn lý thuyết kinh
điển, Uppsala, RBV, Institutional Theory, và Upper Echelons Theory, cùng
với lớp mở rộng Digital Capability Lens, xác định vai trò của từng lý
thuyết trong khung giải thích tổng thể của luận án.

Tổng quan thực nghiệm cho thấy literature I--P đã phát triển qua ba giai
đoạn, từ tuyến tính dương đến phi tuyến đến tập trung vào moderators;
các meta-analyses nhất quán cho thấy tác động tổng hợp là dương nhưng
nhỏ với I² \textgreater{} 80\%, đòi hỏi cách tiếp cận moderator-centric;
và bối cảnh châu Á, đặc biệt là các nền kinh tế frontier và SIDS, đang
thiếu đại diện nghiêm trọng trong literature.

Ba khoảng trống được xác định: thiếu khung tích hợp đa tầng, sự đồng
nhất không hợp lý TCI--DAI, và thiếu bằng chứng cross-country quy mô lớn
có cấu trúc cho châu Á. Để lấp đầy ba khoảng trống này, luận án đề xuất
khung CDCM và phát triển sáu giả thuyết (H1--H6) bao phủ phi tuyến cơ
bản, vai trò của TCI và DAI, vai trò của đặc điểm nhà quản trị, vai trò
của thể chế ICRV, và dị biệt thời gian. Chương 3 sẽ trình bày thiết kế
phương pháp và chiến lược kiểm định hệ giả thuyết này.

\begin{center}\rule{0.5\linewidth}{0.5pt}\end{center}

\hypertarget{chux1b0ux1a1ng-3-phux1b0ux1a1ng-phuxe1p-nghiuxean-cux1ee9u}{%
\section{CHƯƠNG 3: PHƯƠNG PHÁP NGHIÊN
CỨU}\label{chux1b0ux1a1ng-3-phux1b0ux1a1ng-phuxe1p-nghiuxean-cux1ee9u}}

\begin{center}\rule{0.5\linewidth}{0.5pt}\end{center}

\hypertarget{31-thiux1ebft-kux1ebf-nghiuxean-cux1ee9u-tux1ed5ng-thux1ec3}{%
\subsection{3.1 Thiết kế nghiên cứu tổng
thể}\label{31-thiux1ebft-kux1ebf-nghiuxean-cux1ee9u-tux1ed5ng-thux1ec3}}

\hypertarget{311-loux1ea1i-thiux1ebft-kux1ebf}{%
\subsubsection{3.1.1 Loại thiết
kế}\label{311-loux1ea1i-thiux1ebft-kux1ebf}}

Luận án sử dụng \textbf{mixed synthesis-empirical design} gồm hai phase
bổ sung cho nhau. Phase thứ nhất là meta-analysis tổng hợp literature để
trả lời RQ1, còn phase thứ hai là phân tích thực nghiệm doanh nghiệp đa
quốc gia để trả lời RQ2--RQ4. Hai phase đối thoại hai chiều: meta cung
cấp baseline và gợi ý moderator cần kiểm định, còn empirical cung cấp
bằng chứng bổ sung cho meta về bối cảnh châu Á và Pacific.

\textbf{Cơ sở kế thừa}: Mixed design kết hợp meta và primary data đã
được áp dụng trong quản trị chiến lược và IB (Borenstein et al., 2009;
Hunter \& Schmidt, 2004; Combs et al., 2011).

\textbf{Đóng góp mới}: Áp dụng mixed design cho một luận án duy nhất về
internationalization--firm performance trong bối cảnh châu Á và Pacific
với \textbf{47 nền kinh tế} là thiết kế hiếm gặp; đa số các luận án IB
chỉ làm một trong hai phase. Quy mô 47 nước/101.185 doanh nghiệp là phạm
vi địa lý lớn nhất từng có cho khu vực, mở rộng từ pool 17 nước trong Do
và Phan (2026a).

\hypertarget{312-quy-truxecnh-nghiuxean-cux1ee9u}{%
\subsubsection{3.1.2 Quy trình nghiên
cứu}\label{312-quy-truxecnh-nghiuxean-cux1ee9u}}

Quy trình đi theo sáu bước: (i) xác định khoảng trống nghiên cứu từ tổng
quan và meta-analysis trước đây; (ii) xây dựng khung lý thuyết tích hợp
đa tầng và hệ giả thuyết H1--H6; (iii) tiến hành meta-analysis cập nhật
1982--2026; (iv) thu thập và chuẩn hóa dữ liệu WBES cho \textbf{47 nền
kinh tế châu Á và Pacific} sau khi hòa hợp 3 thế hệ schema (PICS3,
Standardized, BREADY/BEE); (v) ước lượng mô hình thực nghiệm theo bối
cảnh quốc gia và đa quốc gia; (vi) kiểm định độ vững và tích hợp kết quả
để đóng góp lý thuyết.

\hypertarget{32-phase-1--meta-analysis-19822026}{%
\subsection{3.2 Phase 1 --- Meta-analysis
1982--2026}\label{32-phase-1--meta-analysis-19822026}}

\hypertarget{321-cux1a1-sux1edf-kux1ebf-thux1eeba}{%
\subsubsection{3.2.1 Cơ sở kế
thừa}\label{321-cux1a1-sux1edf-kux1ebf-thux1eeba}}

Meta-analysis trong IB đã được thiết lập vững chắc qua nhiều thập kỷ
(Hunter \& Schmidt, 2004; Borenstein et al., 2009). Phương pháp chuẩn
bao gồm: xây dựng inclusion criteria theo PRISMA (Page et al., 2021),
transmit hiệu ứng bằng Fisher\textquotesingle s \(z\), tính pooled
effect size theo random-effects, kiểm định heterogeneity bằng \(Q\)-test
và \(I^2\), và đánh giá publication bias bằng Egger và Begg-Mazumdar
(Egger et al., 1997; Begg \& Mazumdar, 1994). Các meta-analyses trước
đây về I--P (Bausch \& Krist, 2007; Kirca et al., 2012; Marano et al.,
2016; Wu et al., 2022; Arte \& Larimo, 2022) cung cấp quy trình mẫu mà
luận án kế thừa và cập nhật.

\hypertarget{322-ux111uxf3ng-guxf3p-mux1edbiux111iux1ec1u-chux1ec9nh}{%
\subsubsection{3.2.2 Đóng góp mới/điều
chỉnh}\label{322-ux111uxf3ng-guxf3p-mux1edbiux111iux1ec1u-chux1ec9nh}}

\begin{itemize}
\tightlist
\item
  \textbf{Mở rộng coverage time}: từ 1977--2022 (bài ICBEF 2025) sang
  \textbf{1982--2026}, bổ sung các nghiên cứu 2023--2026 liên quan đến
  digital transformation và EMNEs.
\item
  \textbf{Mở rộng moderator coverage}: thêm cụm moderator về digital
  capability và 6 sub-regime ICRV, chưa được xem xét đầy đủ trong các
  meta-analyses trước.
\item
  \textbf{Pool hiện tại}: K=288 effect sizes / k=238 nghiên cứu
  (baseline trước WoS+Scopus formal search); mục tiêu sau formal search:
  k≥250, tăng so với bản nền ICBEF 2025 (K=200 / k=113).
\item
  \textbf{Subgroup theo châu Á và Pacific}: tách riêng nhóm nghiên cứu
  về doanh nghiệp châu Á và Pacific để cung cấp baseline trực tiếp cho
  phase empirical, đối chiếu với pool 17 nước trong Do và Phan (2026a).
\end{itemize}

\hypertarget{323-quy-truxecnh-prisma}{%
\subsubsection{3.2.3 Quy trình PRISMA}\label{323-quy-truxecnh-prisma}}

Identification → Screening → Eligibility → Inclusion theo hướng dẫn
PRISMA 2020 (Page et al., 2021). Inclusion criteria: nghiên cứu thực
nghiệm ở cấp doanh nghiệp, có đo lường internationalization và firm
performance, có đủ thống kê để tính effect size (Pearson \(r\),
\(\beta\), \(t\)-stat).

\hypertarget{33-phase-2--phuxe2n-tuxedch-thux1ef1c-nghiux1ec7m-ux111a-quux1ed1c-gia}{%
\subsection{3.3 Phase 2 --- Phân tích thực nghiệm đa quốc
gia}\label{33-phase-2--phuxe2n-tuxedch-thux1ef1c-nghiux1ec7m-ux111a-quux1ed1c-gia}}

\hypertarget{331-nguux1ed3n-dux1eef-liux1ec7u}{%
\subsubsection{3.3.1 Nguồn dữ
liệu}\label{331-nguux1ed3n-dux1eef-liux1ec7u}}

Luận án sử dụng World Bank Enterprise Surveys cho \textbf{47 nền kinh tế
châu Á và Pacific} trong giai đoạn \textbf{2009--2025} (101.185 doanh
nghiệp · 108 cặp quốc gia × năm · 14 mốc khảo sát) (World Bank, n.d.,
2019, 2023, 2026). WBES là bộ dữ liệu doanh nghiệp chuẩn hóa, phù hợp
cho nghiên cứu so sánh đa quốc gia về firm-level outcomes (Aterido et
al., 2011).

\textbf{Cơ sở kế thừa}: WBES đã được sử dụng rộng rãi trong nghiên cứu
IB về emerging markets (Ayyagari et al., 2011; Cuervo-Cazurra et al.,
2018; Chen \& Meng, 2022), đồng thời đã được sử dụng trong bài nghiên
cứu thực trạng emerging Asia với 17 nền kinh tế và \textasciitilde40.633
quan sát doanh nghiệp (Do \& Phan, 2026a).

\textbf{Đóng góp mới}: mở rộng quy mô từ 17 sang \textbf{47 nền kinh tế
châu Á và Pacific}, với 6 sub-regime ICRV (Advanced innovation-driven,
Advanced resource-driven, Upper-middle, Emerging, Frontier, Pacific
SIDS). Quy mô này là lớn nhất cho khu vực trong literature
internationalization--performance, gồm cả boundary case Pacific SIDS (9
nước, n=1.469 (mẫu phân tích sau lọc missing)) cho phép kiểm định forced
internationalization penalty (Do \& Phan, 2026b).

\begin{quote}
\textbf{Ghi chú về trọng số khảo sát (survey sampling weights):} WBES
cung cấp sampling weights để đại diện cho tổng thể doanh nghiệp trong
từng quốc gia. Tuy nhiên, luận án \emph{không áp dụng sampling weights}
trong ước lượng mô hình hồi quy vì ba lý do: (1) Mục tiêu của luận án là
kiểm định cơ chế lý thuyết (mechanism testing) chứ không phải ước lượng
mô tả đại diện tổng thể (population-representative descriptive
estimation), trong bối cảnh này, không áp dụng weights là thực hành
chuẩn trong literature IB sử dụng WBES (Aterido et al., 2011;
Cuervo-Cazurra et al., 2018); (2) Fixed effects ở cấp country-year đã
kiểm soát đặc điểm cấu trúc của từng sóng khảo sát; (3) Việc áp dụng
weights có thể làm sai lệch ước lượng khi mẫu được gộp chung từ nhiều
quốc gia có quy mô tổng thể doanh nghiệp khác nhau nhiều bậc (ví dụ:
Trung Quốc \textasciitilde100 triệu so với Tonga \textasciitilde500
doanh nghiệp). Phân tích độ nhạy có áp dụng survey weights cho từng quốc
gia riêng lẻ (P3/P4/P5) không thay đổi hướng hoặc mức độ ý nghĩa của các
hệ số chính (xem Mục 3.5).
\end{quote}

\hypertarget{332-ux111o-lux1b0ux1eddng-biux1ebfn-phux1ee5-thuux1ed9c--hiux1ec7u-quux1ea3-houx1ea1t-ux111ux1ed9ng}{%
\subsubsection{3.3.2 Đo lường biến phụ thuộc --- Hiệu quả hoạt
động}\label{332-ux111o-lux1b0ux1eddng-biux1ebfn-phux1ee5-thuux1ed9c--hiux1ec7u-quux1ea3-houx1ea1t-ux111ux1ed9ng}}

\textbf{Biến chính}: Labor productivity đo bằng
\(\ln(\text{annual sales} / \text{permanent full-time employees})\).

\textbf{Cơ sở kế thừa}: Labor productivity được sử dụng rộng rãi trong
nghiên cứu WBES và trong literature về firm productivity (Bloom et al.,
2012; Hsieh \& Klenow, 2009). Chỉ tiêu này đã được vận hành hóa thành
công trong bài thực trạng emerging Asia (Do \& Phan, 2026a).

\textbf{Lập luận chọn}: WBES không phải bộ financial statement đầy đủ,
nên labor productivity phù hợp hơn ROA/ROE/ROS trong bối cảnh đa quốc
gia. Ngoài ra, labor productivity có tính so sánh xuyên quốc gia cao hơn
do không bị ảnh hưởng bởi chuẩn mực kế toán khác biệt (Combs et al.,
2005; Richard et al., 2009).

\textbf{Biến robustness}: ROS, sales growth, employment growth,
\(\ln\)(annual sales).

\textbf{Đóng góp mới}: lập luận rõ ràng về firm performance là khái niệm
đa chiều và việc lựa chọn thước đo phải phù hợp với bối cảnh dữ liệu
(xem \texttt{02\_theoretical\_framework\_vi.md}).

\hypertarget{333-ux111o-lux1b0ux1eddng-biux1ebfn-ux111ux1ed9c-lux1eadp--quux1ed1c-tux1ebf-huxf3a}{%
\subsubsection{3.3.3 Đo lường biến độc lập --- Quốc tế
hóa}\label{333-ux111o-lux1b0ux1eddng-biux1ebfn-ux111ux1ed9c-lux1eadp--quux1ed1c-tux1ebf-huxf3a}}

\textbf{Biến chính}: Foreign Sales To Total Sales (FSTS), đo bằng tỷ lệ
doanh thu xuất khẩu trên tổng doanh thu. Mô hình có cả \(FSTS\),
\(FSTS^2\), và \(FSTS^3\) để kiểm định phi tuyến S-curve / cubic.

\textbf{Cơ sở kế thừa}: FSTS là thước đo internationalization phổ biến
nhất trong I--P literature (Sullivan, 1994; Hitt et al., 2006). Việc bổ
sung \(FSTS^2\) để kiểm định phi tuyến đã được áp dụng trong Lu và
Beamish (2004), Contractor et al. (2003), và Hitt et al. (1997).
Specification cubic được Do và Phan (2026g) xác nhận có ý nghĩa thống kê
ở Trung Quốc với turning point \textasciitilde47,8\% FSTS, làm baseline
cho việc kiểm định mở rộng trong luận án.

\hypertarget{334-ux111o-lux1b0ux1eddng-biux1ebfn-ux111iux1ec1u-tiux1ebft}{%
\subsubsection{3.3.4 Đo lường biến điều
tiết}\label{334-ux111o-lux1b0ux1eddng-biux1ebfn-ux111iux1ec1u-tiux1ebft}}

\hypertarget{digital-constructs-tci-vuxe0-dai}{%
\paragraph{Digital constructs (TCI và
DAI)}\label{digital-constructs-tci-vuxe0-dai}}

Luận án phân biệt rõ hai construct:

\begin{itemize}
\tightlist
\item
  \textbf{Technological Capability Index (TCI)}: chiều sâu năng lực công
  nghệ, tổng hợp từ foreign technology licensing, R\&D spending, product
  innovation, và quality certification.
\item
  \textbf{Digital Adoption Index (DAI)}: mức độ chấp nhận số bề mặt,
  tổng hợp từ website, email, online sales (DAI\_thin) hoặc bổ sung
  e-payment và e-supplier (DAI\_rich).
\end{itemize}

\textbf{Nguyên tắc non-overlapping}: TCI và DAI không được chia sẻ
component để bảo đảm construct purity. Đặc biệt, foreign technology
licensing thuộc TCI chứ không được đồng thời tính vào DAI.

\textbf{Cơ sở kế thừa}: Bhandari et al. (2023) phân biệt digitalization,
internationalization và capability theo resource-orchestration logic;
Verhoef et al. (2021) phân biệt three stages (digitization,
digitalization, digital transformation). Pattern "digital shield effect"
được Do và Phan (2026a) phát hiện trên 17 nước châu Á mới nổi.

\textbf{Đóng góp mới}: Đưa ra \textbf{measurement harmonization
protocol} trong bối cảnh WBES, đảm bảo construct purity ở 47 nền kinh tế
châu Á và Pacific, với pipeline hòa hợp 3 thế hệ schema. Phân tách TCI
và DAI theo (i) Bharadwaj et al. (2013) --- IT capability vs digital
business strategy có nomological nets khác nhau; (ii) phân tầng 3 cấp
của Verhoef et al. (2021), digitization, digitalization, digital
transformation; (iii) truyền thống Lall (1992) và Cohen \& Levinthal
(1990) coi technological capability là chiều sâu năng lực nội tại
(absorptive capacity), khác về bản chất với digital adoption ở giao diện
ngoại tại; (iv) resource-orchestration logic của Bhandari et al. (2023)
cho I→P relationship; và (v) 4 tiêu chí formative composite của Coltman
et al. (2008), đảm bảo TCI và DAI là hai composite độc lập.

\hypertarget{biux1ebfn-thux1ec3-chux1ebf}{%
\paragraph{Biến thể chế}\label{biux1ebfn-thux1ec3-chux1ebf}}

\textbf{Biến chính}: business obstacle index (tổng hợp từ các câu hỏi
WBES về obstacles), country-level governance proxies, chỉ số pháp quyền
WGI (Kaufmann et al., 2011), 6 sub-regime ICRV (classification từ Khanna
\& Palepu, 2010 mở rộng).

\textbf{Cơ sở kế thừa}: Khanna và Palepu (2010) với institutional voids;
North (1990) với institutional analysis; Marano et al. (2016) với
meta-analytic evidence về institutional moderators.

\hypertarget{ux111ux1eb7c-ux111iux1ec3m-nhuxe0-quux1ea3n-trux1ecb-cux1ea5p-cao}{%
\paragraph{Đặc điểm nhà quản trị cấp
cao}\label{ux111ux1eb7c-ux111iux1ec3m-nhuxe0-quux1ea3n-trux1ecb-cux1ea5p-cao}}

\textbf{Biến chính}: top manager experience (năm kinh nghiệm), top
manager gender (female / male).

\textbf{Cơ sở kế thừa}: Hambrick và Mason (1984), Hambrick (2007), Hsu
et al. (2013), Post và Byron (2015).

\textbf{Điều kiện sử dụng}: chỉ trong subset có đủ dữ liệu; nếu thiếu,
đưa vào robustness chứ không bắt buộc trong mô hình chính.

\hypertarget{335-biux1ebfn-kiux1ec3m-souxe1t}{%
\subsubsection{3.3.5 Biến kiểm
soát}\label{335-biux1ebfn-kiux1ec3m-souxe1t}}

\(\ln\) employment, firm age, foreign ownership dummy, sector dummy,
country fixed effects, year fixed effects.

\textbf{Cơ sở kế thừa}: chuẩn trong nghiên cứu WBES và IB (Aterido et
al., 2011; Ayyagari et al., 2011).

\hypertarget{34-muxf4-huxecnh-phuxe2n-tuxedch}{%
\subsection{3.4 Mô hình phân
tích}\label{34-muxf4-huxecnh-phuxe2n-tuxedch}}

\hypertarget{341-muxf4-huxecnh-meta-analysis}{%
\subsubsection{3.4.1 Mô hình
meta-analysis}\label{341-muxf4-huxecnh-meta-analysis}}

Random-effects model dựa trên Fisher\textquotesingle s \(z\):

\[z_i = \frac{1}{2} \ln\left(\frac{1+r_i}{1-r_i}\right)\]

với \(r_i\) là effect size của nghiên cứu \(i\). Pooled effect size tính
theo inverse-variance weighting (Borenstein et al., 2009).

Heterogeneity đo bằng \(Q\)-test và \(I^2\) (Higgins et al., 2003).
Subgroup analysis và meta-regression để kiểm định moderators.

\textbf{Cơ sở kế thừa}: chuẩn meta-analysis trong quản trị chiến lược
(Combs et al., 2011; Aguinis et al., 2011).

\hypertarget{342-muxf4-huxecnh-empirical--tuyux1ebfn-tuxednh-cux1a1-bux1ea3n}{%
\subsubsection{3.4.2 Mô hình empirical --- Tuyến tính cơ
bản}\label{342-muxf4-huxecnh-empirical--tuyux1ebfn-tuxednh-cux1a1-bux1ea3n}}

\[\ln(LP)_i = \beta_0 + \beta_1 FSTS_i + \boldsymbol{\gamma} \mathbf{X}_i + \varepsilon_i\]

với \(\mathbf{X}_i\) là vector control variables.

\hypertarget{343-muxf4-huxecnh-phi-tuyux1ebfn-h1}{%
\subsubsection{3.4.3 Mô hình phi tuyến
(H1)}\label{343-muxf4-huxecnh-phi-tuyux1ebfn-h1}}

\[\ln(LP)_i = \beta_0 + \beta_1 FSTS_i + \beta_2 FSTS_i^2 + \boldsymbol{\gamma} \mathbf{X}_i + \varepsilon_i\]

Nếu \(\beta_1 > 0\) và \(\beta_2 < 0\) --- inverted-U được xác nhận.
Turning point tính theo \(-\beta_1 / (2\beta_2)\).

Đối với Trung Quốc, mô hình mở rộng cubic theo specification đã được
công bố:

\[\ln(LP)_i = \beta_0 + \beta_1 FSTS_i + \beta_2 FSTS_i^2 + \beta_3 FSTS_i^3 + \boldsymbol{\gamma} \mathbf{X}_i + \varepsilon_i\]

\textbf{Cơ sở kế thừa}: Lu và Beamish (2004), Hitt et al. (1997),
Contractor et al. (2003); Do và Phan (2026g) cho cubic specification ở
China với turning point \textasciitilde47,8\% FSTS.

\textbf{Lind--Mehlum U-test} (Lind \& Mehlum, 2010; Haans et al., 2016)
để xác nhận chặt inverted-U: kiểm định slope dương ở bên trái và slope
âm ở bên phải của domain \(FSTS\).

\hypertarget{344-muxf4-huxecnh-moderation-h2h6}{%
\subsubsection{3.4.4 Mô hình moderation
(H2--H6)}\label{344-muxf4-huxecnh-moderation-h2h6}}

\hypertarget{two-way-interaction-h2-h3}{%
\paragraph{Two-way interaction (H2,
H3)}\label{two-way-interaction-h2-h3}}

\[\ln(LP)_i = \beta_0 + \beta_1 FSTS_i + \beta_2 FSTS_i^2 + \beta_3 M_i + \beta_4 (FSTS_i \times M_i) + \beta_5 (FSTS_i^2 \times M_i) + \boldsymbol{\gamma} \mathbf{X}_i + \varepsilon_i\]

với \(M_i\) là moderator (TCI hoặc DAI).

\hypertarget{three-way-interaction-synthesis}{%
\paragraph{Three-way interaction
(synthesis)}\label{three-way-interaction-synthesis}}

\[\ln(LP)_i = \beta_0 + \beta_1 FSTS_i + \beta_2 FSTS_i^2 + \beta_3 D_i + \beta_4 I_i + \beta_5 M_i + \sum \text{interactions} + \boldsymbol{\gamma} \mathbf{X}_i + \varepsilon_i\]

với \(D\) = digital, \(I\) = institutional, \(M\) = manager.

\textbf{Cơ sở kế thừa}: chuẩn moderation trong quản trị chiến lược
(Aiken \& West, 1991; Dawson, 2014).

\textbf{Đóng góp mới}: three-way moderation digital × institutional ×
manager trong bối cảnh châu Á và Pacific với 47 nước chưa được kiểm định
trước đây.

\hypertarget{345-muxf4-huxecnh-temporal-heterogeneity-h6}{%
\subsubsection{3.4.5 Mô hình temporal heterogeneity
(H6)}\label{345-muxf4-huxecnh-temporal-heterogeneity-h6}}

\[\ln(LP)_i = \beta_0 + \beta_1 FSTS_i + \beta_2 FSTS_i^2 + \beta_3 Y_{2024,i} + \beta_4 (FSTS_i \times Y_{2024,i}) + \beta_5 (FSTS_i^2 \times Y_{2024,i}) + \boldsymbol{\gamma} \mathbf{X}_i + \varepsilon_i\]

Nếu \(\beta_4\) hoặc \(\beta_5\) có ý nghĩa, hình dạng đường quan hệ đã
dịch chuyển giữa hai mốc thời gian.

\textbf{Cơ sở kế thừa}: Wu et al. (2022) với twenty-year evolution; Xiao
et al. (2013) với China-specific contingencies; Do và Phan (2026g) làm
baseline cubic 2012.

\textbf{Đóng góp mới}: áp dụng cho China 2012 và 2024 để kiểm định
stability của cubic turning point (\textasciitilde47,8\% FSTS) sau hơn
một thập kỷ chuyển đổi số.

\hypertarget{3451-muxf4-huxecnh-cux1ee5-thux1ec3--nghiuxean-cux1ee9u-3-viux1ec7t-nam-wbes-200920152023}{%
\subsubsection{3.4.5.1 Mô hình cụ thể --- Nghiên cứu 3 (Việt Nam, WBES
2009/2015/2023)}\label{3451-muxf4-huxecnh-cux1ee5-thux1ec3--nghiuxean-cux1ee9u-3-viux1ec7t-nam-wbes-200920152023}}

Nghiên cứu 3 sử dụng chuỗi mô hình lồng nhau M0--M8 ước lượng riêng theo
từng sóng khảo sát (2009, 2015, 2023) và gộp pooled. Ký hiệu tiếng Việt:

\begin{itemize}
\tightlist
\item
  \textbf{lnNSLD\_it} = lnLP\_it: log năng suất lao động (ln doanh thu
  PPP / lao động thường xuyên)
\item
  \textbf{CDDXK\_c\_it} = FSTS\_c\_it: cường độ xuất khẩu điều chỉnh
  trung bình trong sóng (mean-centred)
\item
  \textbf{CDDXK\_c²\_it} = FSTS\_c²\_it: bình phương cường độ xuất khẩu
  điều chỉnh
\item
  \textbf{NLCN\_z\_it} = TCI\_z\_it: năng lực công nghệ (chuẩn hoá z
  trong sóng)
\item
  \textbf{CSS\_z\_it} = DAI\_z\_it: chỉ số số hoá cơ sở --- Tier-1
  (chuẩn hoá z trong sóng)
\item
  \textbf{lnLD, TuoiDN, SoHuuNuocNgoai}: kiểm soát quy mô, tuổi, sở hữu
\item
  \textbf{δ\_s}: hiệu ứng cố định ngành (1-digit ISIC); \textbf{λ\_t}:
  hiệu ứng cố định sóng (chỉ pooled)
\end{itemize}

\textbf{M0 --- Mô hình cơ sở (biến kiểm soát):}
\[\ln NSLD_{it} = \alpha + \gamma_1 \ln LD_{it} + \gamma_2 TuoiDN_{it} + \gamma_3 SoHuuNN_{it} + \delta_s + [\lambda_t] + \varepsilon_{it}\]

\textbf{M1 --- Quốc tế hoá tuyến tính:}
\[\ln NSLD_{it} = \alpha + \beta_1 CDDXK\_c_{it} + \boldsymbol{\gamma}\mathbf{X}_{it} + \delta_s + [\lambda_t] + \varepsilon_{it}\]

\textbf{M2 --- Quốc tế hoá phi tuyến (kiểm định H1 --- inverted-U):}
\[\ln NSLD_{it} = \alpha + \beta_1 CDDXK\_c_{it} + \beta_2 CDDXK\_c^2_{it} + \boldsymbol{\gamma}\mathbf{X}_{it} + \delta_s + [\lambda_t] + \varepsilon_{it}\]

H1: \(\beta_1 > 0\) và \(\beta_2 < 0\); điểm quay
\(TP^* = -\beta_1 / (2\beta_2)\), xác nhận bởi kiểm định Lind--Mehlum
(2010).

\textbf{M3, Điều tiết NLCN (kiểm định H2):}
\[\ln NSLD_{it} = \alpha + \beta_1 CDDXK\_c + \beta_2 CDDXK\_c^2 + \beta_3 NLCN\_z + \beta_4(CDDXK\_c \times NLCN\_z) + \beta_5(CDDXK\_c^2 \times NLCN\_z) + \boldsymbol{\gamma}\mathbf{X} + \delta_s + [\lambda_t] + \varepsilon\]

\textbf{M4 --- Điều tiết CSS/DAI (H4 thăm dò):}
\[\ln NSLD_{it} = \alpha + \beta_1 CDDXK\_c + \beta_2 CDDXK\_c^2 + \beta_3 CSS\_z + \beta_4(CDDXK\_c \times CSS\_z) + \beta_5(CDDXK\_c^2 \times CSS\_z) + \boldsymbol{\gamma}\mathbf{X} + \delta_s + [\lambda_t] + \varepsilon\]

\textbf{M5 --- NLCN trực tiếp (không tương tác):}
\[\ln NSLD_{it} = \alpha + \beta_1 CDDXK\_c + \beta_2 CDDXK\_c^2 + \beta_3 NLCN\_z + \boldsymbol{\gamma}\mathbf{X} + \delta_s + [\lambda_t] + \varepsilon\]

\textbf{M6 --- CSS/DAI trực tiếp (không tương tác):}
\[\ln NSLD_{it} = \alpha + \beta_1 CDDXK\_c + \beta_2 CDDXK\_c^2 + \beta_3 CSS\_z + \boldsymbol{\gamma}\mathbf{X} + \delta_s + [\lambda_t] + \varepsilon\]

\textbf{M7 --- Cả hai trực tiếp, không tương tác:}
\[\ln NSLD_{it} = \alpha + \beta_1 CDDXK\_c + \beta_2 CDDXK\_c^2 + \beta_3 NLCN\_z + \beta_4 CSS\_z + \boldsymbol{\gamma}\mathbf{X} + \delta_s + [\lambda_t] + \varepsilon\]

\textbf{M8 --- Mô hình đầy đủ (trực tiếp + điều tiết CSS):}
\[\ln NSLD_{it} = \alpha + \beta_1 CDDXK\_c + \beta_2 CDDXK\_c^2 + \beta_3 NLCN\_z + \beta_4 CSS\_z + \beta_5(CDDXK\_c \times CSS\_z) + \beta_6(CDDXK\_c^2 \times CSS\_z) + \boldsymbol{\gamma}\mathbf{X} + \delta_s + [\lambda_t] + \varepsilon\]

\emph{Tất cả mô hình đều ước lượng bằng OLS với sai số chuẩn vững HC1.
Hiệu ứng cố định sóng \(\lambda_t\) chỉ áp dụng trong mô hình pooled. Mô
hình được ước lượng riêng cho 3 sóng và gộp (N = 989 / 956 / 1.013 /
2.958).}

\textbf{Bảng định nghĩa biến --- Nghiên cứu 3 (Việt Nam):}

\begin{longtable}[]{@{}llll@{}}
\toprule\noalign{}
Ký hiệu & Biến WBES & Cách tính & Vai trò \\
\midrule\noalign{}
\endhead
\bottomrule\noalign{}
\endlastfoot
lnNSLD & d2, l1 & ln(d2 / l1): ln(doanh thu PPP / lao động thường xuyên)
& Biến phụ thuộc \\
CDDXK & d3c & d3c / 100: cường độ xuất khẩu trực tiếp (0--1) & Biến độc
lập \\
CDDXK\_c & d3c & CDDXK − trung bình sóng: chuẩn hoá điều chỉnh & Biến
độc lập (centred) \\
CDDXK\_c² & d3c & CDDXK\_c bình phương: hạng phi tuyến & Kiểm định
inverted-U (H1) \\
NLCN\_z & b8, e6 & z-std trong sóng của TB(b8₀₁, e6₀₁): chứng chỉ chất
lượng + công nghệ ngoại & Năng lực công nghệ (H2) \\
CSS\_z & c22b & z-std trong sóng của c22b₀₁: hiện diện website & Số hoá
Tier-1 (H4 thăm dò) \\
lnLD & l1 & ln(l1): log số lao động thường xuyên & Kiểm soát: quy mô
doanh nghiệp \\
TuoiDN & b5 & năm khảo sát − b5: số năm hoạt động & Kiểm soát: tuổi
doanh nghiệp \\
SoHuuNN & b2b & 1 nếu b2b \textgreater{} 0: có vốn nước ngoài & Kiểm
soát: hình thức sở hữu \\
δ\_s & a4b / a4a & ngành 1-chữ số ISIC (a4b cho 2009/2015; a4a cho 2023)
& Hiệu ứng cố định ngành \\
λ\_t & wave & chỉ báo sóng khảo sát (chỉ pooled) & Hiệu ứng cố định thời
kỳ \\
\end{longtable}

\hypertarget{3452-muxf4-huxecnh-cux1ee5-thux1ec3--nghiuxean-cux1ee9u-4-singapore-wbes-2023}{%
\subsubsection{3.4.5.2 Mô hình cụ thể --- Nghiên cứu 4 (Singapore, WBES
2023)}\label{3452-muxf4-huxecnh-cux1ee5-thux1ec3--nghiuxean-cux1ee9u-4-singapore-wbes-2023}}

Nghiên cứu 4 sử dụng dữ liệu mặt cắt ngang WBES Singapore 2023 (N =
623). Thiết kế đơn sóng, không có thành phần λ\_t. Ký hiệu tiếng Việt:

\begin{itemize}
\tightlist
\item
  \textbf{lnNSLD\_i} = lnLP\_i: log năng suất lao động (ln doanh thu PPP
  / lao động thường xuyên)
\item
  \textbf{CDDXK\_i} = FSTS\_i: cường độ xuất khẩu trực tiếp (d3c / 100)
\item
  \textbf{CDDXK\_c\_i} = FSTS\_c\_i: CDDXK trung bình mẫu; CDDXK\_c² là
  bình phương
\item
  \textbf{NLCN\_z\_i} = TCI\_z\_i: năng lực công nghệ, z-std của trung
  bình(b8₀₁, e6₀₁)
\item
  \textbf{CSS\_z\_i} = DAI\_z\_i: chỉ số số hoá \textbf{Tier-1+2}, z-std
  của trung bình(c22b₀₁, k33₀₁, k38₀₁); khác P3 ở chỗ bao gồm e-payment
  hai chiều (Tier-2)
\item
  \textbf{lnLD\_i}, \textbf{TuoiDN\_i}, \textbf{SoHuuNN\_i}: biến kiểm
  soát tương tự P3
\item
  \textbf{δ\_s}: sector fixed effects (ISIC 1-chữ số)
\item
  \textbf{IMR\_i} = nghịch đảo tỷ lệ Mills từ mô hình probit tham gia
  xuất khẩu (kiểm tra độ nhạy Heckman)
\end{itemize}

\textbf{Chuỗi mô hình lồng nhau M0--M5:}

\begin{quote}
M0 (Baseline): lnNSLD\_i = α + γ₁ lnLD\_i + γ₂ TuoiDN\_i + γ₃ SoHuuNN\_i
+ δ\_s + ε\_i
\end{quote}

\begin{quote}
M1 (Tuyến tính FSTS): lnNSLD\_i = α + β₁ CDDXK\_c\_i + γ·X\_i + δ\_s +
ε\_i
\end{quote}

\begin{quote}
M2 (Bậc hai FSTS, kiểm định H1): lnNSLD\_i = α + β₁ CDDXK\_c\_i + β₂
CDDXK\_c²\_i + γ·X\_i + δ\_s + ε\_i H1: β₁ \textgreater{} 0, β₂
\textless{} 0; TP* = −β₁/(2β₂) ≈ 88,6\% FSTS {[}CI bootstrap {[}53\%,
253\%{]}{]} Lưu ý: Lind--Mehlum p = 0,303, không bác bỏ tuyến tính trong
phạm vi dữ liệu quan sát; đây là kết quả thông tin tích cực theo khung
bão hòa (saturation framework)
\end{quote}

\begin{quote}
M3 (+ NLCN --- H1-TCI trực tiếp): lnNSLD\_i = α + β₁ CDDXK\_c + β₂
CDDXK\_c² + β₃ NLCN\_z + γ·X + δ\_s + ε
\end{quote}

\begin{quote}
M4 (+ CSS trực tiếp): lnNSLD\_i = α + β₁ CDDXK\_c + β₂ CDDXK\_c² + β₃
NLCN\_z + β₄ CSS\_z + γ·X + δ\_s + ε
\end{quote}

\begin{quote}
M5 (Mô hình đầy đủ, kiểm định H3): lnNSLD\_i = α + β₁ CDDXK\_c + β₂
CDDXK\_c² + β₃ NLCN\_z + β₄ CSS\_z + β₅(CDDXK\_c × CSS\_z) +
β₆(CDDXK\_c² × CSS\_z) + γ·X + δ\_s + ε H3: β₆ \textgreater{} 0 --- DAI
khuếch đại lợi nhuận I→P ở cường độ xuất khẩu cao (coordination platform
mechanism; Stallkamp \& Schotter, 2021) Kết quả: β₆ = +3,119 (p = 0,005)
trong mẫu đầy đủ; β₆ = +2,821 (p = 0,003, F-test) trong mẫu chỉ xuất
khẩu (N = 84, lưu ý: công suất thống kê ≈ 16\%)
\end{quote}

\textbf{Bảng định nghĩa biến --- Nghiên cứu 4 (Singapore):}

\begin{longtable}[]{@{}llll@{}}
\toprule\noalign{}
Ký hiệu VN & Mã WBES & Cách tính & Vai trò \\
\midrule\noalign{}
\endhead
\bottomrule\noalign{}
\endlastfoot
lnNSLD & d2, l1 & ln(d2 / l1) & Biến phụ thuộc \\
CDDXK\_c & d3c & FSTS − mean(FSTS): centred & Biến độc lập \\
CDDXK\_c² & d3c & CDDXK\_c bình phương & Phi tuyến (H1) \\
NLCN\_z & b8, e6 & z-std của TB(b8₀₁, e6₀₁) & Năng lực công nghệ (H2) \\
CSS\_z & c22b, k33, k38 & z-std của TB(c22b₀₁, k33₀₁, k38₀₁) ---
\textbf{Tier-1+2} & Số hoá Tier-1+2 (H3) \\
lnLD & l1 & ln(l1) & Quy mô doanh nghiệp \\
TuoiDN & b5 & năm khảo sát − b5 & Tuổi doanh nghiệp \\
SoHuuNN & b2b & 1 nếu b2b \textgreater{} 0 & Hình thức sở hữu \\
IMR & probit selection & Nghịch đảo tỷ lệ Mills (kiểm tra Heckman) &
Kiểm soát chọn lựa tham gia xuất khẩu \\
δ\_s & a4b & Sector FE (ISIC 1-chữ số) & Hiệu ứng cố định ngành \\
\end{longtable}

\begin{center}\rule{0.5\linewidth}{0.5pt}\end{center}

\hypertarget{3453-muxf4-huxecnh-cux1ee5-thux1ec3--nghiuxean-cux1ee9u-5-trung-quux1ed1c-wbes-2012-vuxe0-2024}{%
\subsubsection{3.4.5.3 Mô hình cụ thể --- Nghiên cứu 5 (Trung Quốc, WBES
2012 và
2024)}\label{3453-muxf4-huxecnh-cux1ee5-thux1ec3--nghiuxean-cux1ee9u-5-trung-quux1ed1c-wbes-2012-vuxe0-2024}}

Nghiên cứu 5 sử dụng hai sóng WBES Trung Quốc (2012: N = 2.610; 2024: N
= 1.934; pooled N = 4.544). Thiết kế đa sóng không panel (217 doanh
nghiệp xuất hiện cả hai sóng tạo thành "panel core" được sử dụng cho
cluster-robust SE). Ký hiệu tiếng Việt:

\begin{itemize}
\tightlist
\item
  \textbf{lnNSLD\_it} = lnLP\_it: log năng suất lao động (biến phụ
  thuộc)
\item
  \textbf{CDDXK\_c\_it} = FSTS\_c\_it: cường độ xuất khẩu centred theo
  trung bình sóng
\item
  \textbf{NLCN\_z\_it} = TCI\_full\_z\_it: năng lực công nghệ toàn diện,
  z-std của TB(b8₀₁, e6₀₁, b4₀₁, b7a₀₁); mở rộng hơn P3
\item
  \textbf{CSS\_z\_it} = DAI\_core\_it: chỉ số số hoá cơ bản, chỉ c22b₀₁
  (Tier-1 mỏng, single-item binary)
\item
  \textbf{lnLD\_it}, \textbf{TuoiDN\_it}, \textbf{SoHuuNN\_it}: biến
  kiểm soát
\item
  \textbf{δ\_s}: sector fixed effects; \textbf{λ\_t}: wave fixed effects
  (pooled)
\end{itemize}

\textbf{Chuỗi mô hình lồng nhau M0--M6:}

\begin{quote}
M0 (Baseline): lnNSLD\_it = α + γ₁ lnLD\_it + γ₂ TuoiDN\_it + γ₃
SoHuuNN\_it + δ\_s + λ\_t + ε\_it
\end{quote}

\begin{quote}
M1 (Tuyến tính FSTS): lnNSLD\_it = α + β₁ CDDXK\_c\_it + γ·X\_it + δ\_s
+ λ\_t + ε\_it
\end{quote}

\begin{quote}
M2 (Bậc hai FSTS, kiểm định H1): lnNSLD\_it = α + β₁ CDDXK\_c\_it + β₂
CDDXK\_c²\_it + γ·X\_it + δ\_s + λ\_t + ε\_it H1: β₁ \textgreater{} 0,
β₂ \textless{} 0; TP* = −β₁/(2β₂), điểm uốn 49,4\% (2012), 47,2\%
(2024), và 48,8\% (pooled) Kiểm định ổn định cấu trúc: Paternoster
(1998) z-test: z(FSTS) = +0,82 (p = 0,412); z(FSTS²) = −0,61 (p = 0,545)
→ không bác bỏ bình đẳng hệ số giữa hai sóng
\end{quote}

\begin{quote}
M3 (+ NLCN trực tiếp --- H2 level-shift): lnNSLD\_it = α + β₁ CDDXK\_c +
β₂ CDDXK\_c² + β₃ NLCN\_z + γ·X + δ\_s + λ\_t + ε Kết quả: β₃ = +0,28
(2012, p \textless{} 0,001), +0,43 (2024, p \textless{} 0,001) --- TCI
là "bộ tăng mức" (level-shifter), không điều tiết độ cong
\end{quote}

\begin{quote}
M4 (+ CSS trực tiếp): lnNSLD\_it = α + β₁ CDDXK\_c + β₂ CDDXK\_c² + β₃
NLCN\_z + β₄ CSS\_z + γ·X + δ\_s + λ\_t + ε
\end{quote}

\begin{quote}
M5 (Kiểm định ổn định xuyên sóng --- H2b): Ước lượng M2 riêng biệt cho
2012 và 2024, sau đó áp dụng: z = (β̂₁,2012 − β̂₁,2024) / √(SE²₁,2012 +
SE²₁,2024) (Paternoster et al., 1998), tương tự cho β₂ (FSTS²)
\end{quote}

\begin{quote}
M6 (Điều tiết ba chiều, kiểm định H3/H4a/H4b): lnNSLD\_it = α + β₁
CDDXK\_c + β₂ CDDXK\_c² + β₃ NLCN\_z + β₄ CSS\_z + β₅(CDDXK\_c ×
NLCN\_z) + β₆(CDDXK\_c² × NLCN\_z) + β₇(CDDXK\_c × wave) + β₈(CDDXK\_c²
× wave) + γ·X + δ\_s + λ\_t + ε F-tests: F1 (dịch chuyển độ cong xuyên
sóng), F2 (điều tiết năng lực), F3 (dịch chuyển điều kiện) Kết quả: F2 =
3,26 (p = 0,039, không qua Bonferroni α* = 0,017) → H4b (null
capability-curvature moderation) được chấp nhận
\end{quote}

\textbf{Bảng định nghĩa biến --- Nghiên cứu 5 (Trung Quốc):}

\begin{longtable}[]{@{}llll@{}}
\toprule\noalign{}
Ký hiệu VN & Mã WBES & Cách tính & Vai trò \\
\midrule\noalign{}
\endhead
\bottomrule\noalign{}
\endlastfoot
lnNSLD & d2, l1 & ln(d2 / l1) & Biến phụ thuộc \\
CDDXK\_c & d3c & FSTS − mean\_wave(FSTS) & Biến độc lập \\
CDDXK\_c² & d3c & CDDXK\_c bình phương & Phi tuyến (H1) \\
NLCN\_z & b8, e6, b4, b7a & z-std của TB(b8₀₁, e6₀₁, b4₀₁, b7a₀₁) ---
TCI toàn diện & Năng lực công nghệ (H2) \\
CSS\_z & c22b & c22b₀₁ --- \textbf{Tier-1 đơn biến} (binary) & Số hoá
Tier-1 mỏng (kiểm soát) \\
lnLD & l1 & ln(l1) & Quy mô doanh nghiệp \\
TuoiDN & b5 & năm khảo sát − b5 & Tuổi doanh nghiệp \\
SoHuuNN & b2b & 1 nếu b2b \textgreater{} 0 & Hình thức sở hữu \\
δ\_s & ISIC sector & Sector FE & Hiệu ứng cố định ngành \\
λ\_t & wave (2012/2024) & Wave FE (chỉ pooled) & Hiệu ứng cố định thời
kỳ \\
\end{longtable}

\emph{Ghi chú: 217 doanh nghiệp xuất hiện cả hai sóng ("panel core") tạo
thêm nhận dạng biến thiên trong mẫu trong khi clustered SE theo idstd xử
lý tương quan nội nhóm. DAI Trung Quốc chỉ là Tier-1 đơn biến, không đủ
để kiểm định CDCM; hàm ý chính sách và lý thuyết khác với P4 Singapore
(Tier-1+2).}

\begin{center}\rule{0.5\linewidth}{0.5pt}\end{center}

\hypertarget{3454-muxf4-huxecnh-cux1ee5-thux1ec3--nghiuxean-cux1ee9u-7-ux111a-quux1ed1c-gia-chuxe2u-uxe1-wbes-20032025}{%
\subsubsection{3.4.5.4 Mô hình cụ thể --- Nghiên cứu 7 (Đa quốc gia châu
Á, WBES
2003--2025)}\label{3454-muxf4-huxecnh-cux1ee5-thux1ec3--nghiuxean-cux1ee9u-7-ux111a-quux1ed1c-gia-chuxe2u-uxe1-wbes-20032025}}

Nghiên cứu 7 là kiểm định quy mô lớn nhất của luận án, sử dụng dữ liệu
WBES từ \textbf{49 nền kinh tế châu Á và Thái Bình Dương} (102 cặp quốc
gia × năm, 2003--2025). Mẫu phân tích toàn phần đạt N = 84.910 (M0--M2)
đến N = 29.840 (M11 full). Thiết kế hierarchical OLS với HC1 robust
standard errors, clustered theo country-year. Ký hiệu tiếng Việt:

\begin{itemize}
\tightlist
\item
  \textbf{lnNSLD\_it}: ln(năng suất lao động) = ln(doanh thu / lao động
  thường trực), biến phụ thuộc
\item
  \textbf{CDDXK\_c\_it}: cường độ xuất khẩu mean-centred --- FSTS\_c =
  (d3c/100) − mean\_pool(FSTS)
\item
  \textbf{CDDXK\_c²\_it}: CDDXK\_c bình phương, kiểm định phi tuyến
\item
  \textbf{NLCN\_z\_it}: năng lực công nghệ z-standardised --- TCI\_z =
  z-std(b8, e6)
\item
  \textbf{CSS\_z\_it}: chỉ số số hoá z-standardised --- DAI\_z =
  z-std(website + e-pay, Tier 1+2: c22b/e1, k33)
\item
  \textbf{KNQLy\_it}: kinh nghiệm nhà quản lý trong ngành (năm, b7)
\item
  \textbf{NQL\_nu\_it}: top manager là nữ (binary, b7a/b6a)
\item
  \textbf{ICRV\_j}: phân loại chế độ thể chế ICRV (integer 1--6,
  Advanced→SIDS)
\item
  \textbf{lnLD\_it}: ln(lao động thường trực), kiểm soát quy mô doanh
  nghiệp
\item
  \textbf{TuoiDN\_it}: tuổi doanh nghiệp (năm\_khảo\_sát − b5)
\item
  \textbf{SoHuuNN\_it}: tỷ lệ sở hữu nước ngoài (b6a/100)
\item
  \textbf{δ\_ct}: country × year fixed effects (M5)
\end{itemize}

\textbf{Chuỗi mô hình lồng nhau M0--M11:}

\textbf{M0} (tuyến tính baseline):
\[\ln NSLD_{it} = \alpha + \beta_1 CDDXK\_c_{it} + \varepsilon_{it}\]

\textbf{M2} (phi tuyến --- kiểm định H1):
\[\ln NSLD_{it} = \alpha + \beta_1 CDDXK\_c_{it} + \beta_2 CDDXK\_c^2_{it} + \varepsilon_{it}\]
\[TP^* = -\beta_1 / (2\beta_2); \quad \text{Lind–Mehlum } p < {,}001 \Rightarrow TP \approx 36\%\]

\textbf{M3} (M2 + kiểm soát cơ bản):
\[\ln NSLD_{it} = \alpha + \beta_1 CDDXK\_c + \beta_2 CDDXK\_c^2 + \gamma \cdot X_{it} + \varepsilon_{it}\]

\textbf{M5} (M3 + country-year FE, kiểm định vững mạnh nhất):
\[\ln NSLD_{it} = \alpha + \beta_1 CDDXK\_c + \beta_2 CDDXK\_c^2 + \gamma \cdot X_{it} + \delta_{ct} + \varepsilon_{it}\]
\[\text{AdjR}^2 = {,}677; \quad TP \approx 40\%\]

\textbf{M7} (M3 + điều tiết NLCN, kiểm định H2):
\[\ln NSLD_{it} = \alpha + \beta_1 CDDXK\_c + \beta_2 CDDXK\_c^2 + \beta_3 NLCN\_z\]
\[+ \beta_4 (CDDXK\_c \times NLCN\_z) + \beta_5 (CDDXK\_c^2 \times NLCN\_z) + \gamma \cdot X + \varepsilon\]

\textbf{M8} (M7 + điều tiết CSS, kiểm định H3):
\[\ln NSLD_{it} = \alpha + \beta_1 CDDXK\_c + \beta_2 CDDXK\_c^2 + \beta_3 NLCN\_z + \beta_6 CSS\_z\]
\[+ \beta_7 (CDDXK\_c \times CSS\_z) + \beta_8 (CDDXK\_c^2 \times CSS\_z) + \gamma \cdot X + \varepsilon\]
\[\text{H3: } \hat\beta_7 < 0\ (p < {,}001);\ \hat\beta_8 > 0\ (p < {,}01) \rightarrow \text{DAI nén sườn tăng, giảm nhẹ sườn giảm}\]

\textbf{M9} (M8 + đặc điểm nhà quản lý --- kiểm định H4):
\[+ \beta_9 KNQLy + \beta_{10} NQL\_nu + \beta_{11}(CDDXK\_c \times KNQLy) + \gamma \cdot X + \varepsilon\]
\[\hat\beta_{10} = +0{,}185^{***} \text{ (top manager nữ: lợi thế hiệu quả ~17–20\% cross-context)}\]

\textbf{M10} (M3 + điều tiết ICRV --- kiểm định H5):
\[\ln NSLD_{it} = \alpha + \beta_1 CDDXK\_c + \beta_2 CDDXK\_c^2 + \beta_{11} ICRV_j\]
\[+ \beta_{12}(CDDXK\_c \times ICRV) + \beta_{13}(CDDXK\_c^2 \times ICRV) + \gamma \cdot X + \varepsilon\]
\[\text{H5: } TP \text{ tăng khi ICRV tăng (thể chế yếu hơn): Group I} \approx 28\%,\ \text{Group V–VI} \approx 55\%\]

\textbf{M11} (full three-way, kiểm định H7-P7):
\[\ln NSLD_{it} = \alpha + \beta_1 CDDXK\_c + \beta_2 CDDXK\_c^2 + \beta_3 NLCN\_z + \beta_6 CSS\_z + \beta_{11} ICRV\]
\[+ \beta_{7}(CDDXK\_c \times CSS\_z) + \beta_{8}(CDDXK\_c^2 \times CSS\_z)\]
\[+ \beta_{12}(CDDXK\_c \times ICRV) + \beta_{13}(CDDXK\_c^2 \times ICRV)\]
\[+ \beta_{14}(CDDXK\_c \times CSS\_z \times ICRV) + \beta_9 KNQLy + \beta_{10} NQL\_nu + \gamma \cdot X + \delta_{ct} + \varepsilon\]
\[\text{DAI} \times \text{ICRV: } p = {,}012^*;\ \text{three-way NS};\ TP = 34{,}6\%,\ \text{LM } p = {,}002\]

\textbf{Bảng định nghĩa biến --- Nghiên cứu 7 (đa quốc gia)}

\begin{longtable}[]{@{}llll@{}}
\toprule\noalign{}
Ký hiệu & WBES item & Cách tính & Vai trò \\
\midrule\noalign{}
\endhead
\bottomrule\noalign{}
\endlastfoot
lnNSLD & d2/n3, l1 & ln(doanh thu/lao động) & Biến phụ thuộc \\
CDDXK\_c & d3c & (d3c/100) − mean\_pool(FSTS): centred & Biến độc lập
(I) \\
CDDXK\_c² & d3c & CDDXK\_c bình phương & I --- phi tuyến \\
NLCN\_z & b8, e6 & z-std(cert chất lượng + CN nước ngoài) & Năng lực
công nghệ \\
CSS\_z & c22b/e1, k33 & z-std(website + e-payment): Tier 1+2 & Số hoá
(DAI) \\
KNQLy & b7 & năm kinh nghiệm nhà quản lý trong ngành & Quản trị \\
NQL\_nu & b7a/b6a & 1 nếu top manager là nữ & Quản trị \\
ICRV & phân loại & integer 1--6 (Advanced Innovation→SIDS) & Thể chế \\
lnLD & l1 & ln(lao động thường trực) & Control (quy mô) \\
TuoiDN & b5 & năm\_khảo\_sát − b5 & Control (tuổi DN) \\
SoHuuNN & b6a & tỷ lệ sở hữu nước ngoài (0--1) & Control (sở hữu) \\
δ\_ct & --- & country × year FE & Fixed effects \\
\end{longtable}

\emph{Ghi chú: Mẫu phân tích giảm dần từ N = 84.910 (M2) → 38.342 (M3,
thêm kiểm soát) → 29.840 (M11, thêm manager + ICRV + three-way). DAI P7
là Tier 1+2 khi k33 có mặt, Tier-1 only khi thiếu --- khác với P3
Vietnam (Tier-1 only) và đồng nhất với P4 Singapore (Tier-1+2). ICRV là
biến điều tiết liên tục (integer), không phải dummies, nhằm giữ N đủ lớn
trong M10--M11.}

\begin{center}\rule{0.5\linewidth}{0.5pt}\end{center}

\hypertarget{3455-muxf4-huxecnh-cux1ee5-thux1ec3--nghiuxean-cux1ee9u-8-pacific-sids-n--1469-nhuxf3m-vi}{%
\paragraph{3.4.5.5 Mô hình cụ thể --- Nghiên cứu 8: Pacific SIDS (N =
1.469, Nhóm
VI)}\label{3455-muxf4-huxecnh-cux1ee5-thux1ec3--nghiuxean-cux1ee9u-8-pacific-sids-n--1469-nhuxf3m-vi}}

Nghiên cứu 8 phân tích mẫu gồm \textbf{N = 1.469 doanh nghiệp} tại
\textbf{9 nền kinh tế Pacific SIDS} (Fiji, Kiribati, Papua New Guinea,
Samoa, Solomon Islands, Timor-Leste, Tonga, Vanuatu, Comoros) từ các
sóng WBES 2009--2025. Đây là nhóm thể chế ICRV Nhóm VI --- cấp thấp nhất
trong phân loại 6 chế độ --- đặc trưng bởi thị trường nội địa cực nhỏ,
chi phí thương mại cao và hỗ trợ thể chế yếu. Trong tổng mẫu, chỉ
\textbf{187 doanh nghiệp có hoạt động xuất khẩu} (12,7\%), phản ánh cấu
trúc thị trường đảo nhỏ.

Phát hiện trọng tâm: \textbf{gánh nặng quốc tế hóa bắt buộc (Forced
Internationalization Penalty --- FIP)} --- mối quan hệ quốc tế hóa--hiệu
quả là \textbf{đơn điệu âm} tại Nhóm VI, trái ngược với hình chữ U ngược
được quan sát ở P3--P7. Lind--Mehlum U-test không bác bỏ đơn điệu (p
\textgreater{} ,10), xác nhận không có điểm quay.

\textbf{Ký hiệu biến (tiếng Việt ↔ mã WBES):}

\begin{longtable}[]{@{}llll@{}}
\toprule\noalign{}
Ký hiệu & Tên tiếng Việt & Mã WBES & Cách tính \\
\midrule\noalign{}
\endhead
\bottomrule\noalign{}
\endlastfoot
lnNSLD & Log năng suất lao động & d2, l1 & ln(doanh thu PPP / lao động
thường trực) \\
CDDXK\_c & Cường độ xuất khẩu (điều chỉnh) & d3c & (d3c/100) −
mean\_wave(FSTS): centred theo sóng \\
CDDXK\_c² & CDDXK\_c bình phương & d3c & CDDXK\_c × CDDXK\_c \\
NLCN\_z & Năng lực công nghệ & b8, e6 & z-std(trung bình cert chất lượng
+ CN nước ngoài) \\
CSS\_z & Chỉ số số hoá (Tier-1) & c22b & z-std(website binary):
\textbf{chỉ Tier-1, không dynamic} \\
lnLD & Log quy mô lao động & l1 & ln(lao động thường trực) \\
TuoiDN & Tuổi doanh nghiệp & b5 & năm\_khảo\_sát − b5 \\
SoHuuNN & Tỷ lệ sở hữu nước ngoài & b6a & tỷ lệ 0--1 \\
δ\_c & Country fixed effects & a3b & dummy theo quốc gia \\
λ\_t & Wave fixed effects & wave & dummy theo sóng điều tra \\
\end{longtable}

\emph{Ghi chú: CSS\_z (DAI) là proxy tĩnh Tier-1 (website binary c22b).
WBES không thu thập đủ thang đo cho số hoá động tại các đảo nhỏ ---
KHÔNG gọi là "năng lực số hoá động".}

\textbf{Chuỗi mô hình M0--M3 (P8):}

M0 --- Baseline kiểm soát:

\[lnNSLD_{it} = \alpha + \gamma_1 lnLD_{it} + \gamma_2 TuoiDN_{it} + \gamma_3 SoHuuNN_{it} + \delta_c + \lambda_t + \varepsilon_{it}\]

M1 --- Kiểm định FIP (H1b): thêm cường độ xuất khẩu tuyến tính

\[lnNSLD_{it} = \alpha + \beta_1 CDDXK\_c_{it} + \gamma \cdot X_{it} + \delta_c + \lambda_t + \varepsilon_{it}\]

\[\beta_1 = -0{,}404,\ p = {,}032\ (\text{country+year FE, N} = 1{.}469) \Rightarrow \textbf{FIP xác nhận}\]

M2 --- Kiểm định phi tuyến: thêm CDDXK\_c²

\[lnNSLD_{it} = \alpha + \beta_1 CDDXK\_c_{it} + \beta_2 CDDXK\_c^2_{it} + \gamma \cdot X_{it} + \delta_c + \lambda_t + \varepsilon_{it}\]

\[\beta_1\ \text{NS},\ \beta_2\ \text{NS};\ \text{Lind–Mehlum U-test NS} \Rightarrow \text{không có điểm quay, đơn điệu âm}\]

M3 --- Kiểm định năng lực điều tiết (H2 null cho SIDS):

\[lnNSLD_{it} = \alpha + \beta_1 CDDXK\_c_{it} + \beta_2 CDDXK\_c^2_{it} + \beta_3 NLCN\_z_{it} + \beta_4 CSS\_z_{it} + \gamma \cdot X_{it} + \delta_c + \lambda_t + \varepsilon_{it}\]

\[\beta_3\ (NLCN\_z):\ p = {,}495\ \text{NS};\ \beta_4\ (CSS\_z):\ p = {,}402\ \text{NS} \Rightarrow \text{H2 bị bác bỏ tại Nhóm VI}\]

\textbf{Phân tích độ vững (robustness checks):}

\begin{longtable}[]{@{}llll@{}}
\toprule\noalign{}
Specification & β(CDDXK\_c) & SE & p \\
\midrule\noalign{}
\endhead
\bottomrule\noalign{}
\endlastfoot
M1 country+year FE (N = 1.469) & −0,404 & --- & ,032 \\
Year FE only (không country FE) & −1,236 & --- & \textless{} ,001 \\
Bivariate (không controls) & −1,596 & --- & \textless{} ,001 \\
Exporters-only (N = 187) & −0,901 & --- & ,027 \\
\end{longtable}

\emph{Hệ số âm nhất quán qua 4 specification xác nhận FIP không phải do
selection bias hay outlier.}

\textbf{Bảng định nghĩa biến --- Nghiên cứu 8 (Pacific SIDS)}

\begin{longtable}[]{@{}llll@{}}
\toprule\noalign{}
Ký hiệu & WBES item & Cách tính & Vai trò \\
\midrule\noalign{}
\endhead
\bottomrule\noalign{}
\endlastfoot
lnNSLD & d2, l1 & ln(doanh thu PPP / lao động) & Biến phụ thuộc \\
CDDXK\_c & d3c & (d3c/100) − mean\_wave: centred & Biến độc lập (I) \\
CDDXK\_c² & d3c & CDDXK\_c bình phương & I --- kiểm định phi tuyến \\
NLCN\_z & b8, e6 & z-std(cert + CN nước ngoài) & Năng lực công nghệ \\
CSS\_z & c22b & z-std(website binary): Tier-1 only & Số hoá --- proxy
tĩnh \\
lnLD & l1 & ln(lao động thường trực) & Control (quy mô) \\
TuoiDN & b5 & năm\_KS − năm\_thành\_lập & Control (tuổi DN) \\
SoHuuNN & b6a & tỷ lệ sở hữu nước ngoài & Control (sở hữu) \\
δ\_c & a3b & country FE & Fixed effects \\
λ\_t & wave & wave FE & Fixed effects \\
\end{longtable}

\emph{Ghi chú: FIP là điều kiện biên (boundary condition) của mô hình
chữ U ngược, áp dụng đặc thù tại ICRV Nhóm VI nơi 3 điều kiện cấu trúc
bị vi phạm đồng thời: (1) không có thị trường nội địa đủ lớn, (2) chi
phí thương mại không kiểm soát được, (3) hỗ trợ thể chế không đủ chức
năng. DAI (CSS\_z) Tier-1 only --- WBES không thu thập thang đo đầy đủ
cho số hoá động tại SIDS.}

\begin{center}\rule{0.5\linewidth}{0.5pt}\end{center}

\hypertarget{346-sai-sux1ed1-chuux1ea9n}{%
\subsubsection{3.4.6 Sai số chuẩn}\label{346-sai-sux1ed1-chuux1ea9n}}

Tất cả các mô hình đều sử dụng HC1/HC3 robust standard errors (Long \&
Ervin, 2000; White, 1980) để xử lý heteroscedasticity. Đối với
multi-country, bổ sung clustered SE theo country.

\hypertarget{35-kiux1ec3m-ux111ux1ecbnh-ux111ux1ed9-vux1eefng}{%
\subsection{3.5 Kiểm định độ
vững}\label{35-kiux1ec3m-ux111ux1ecbnh-ux111ux1ed9-vux1eefng}}

\hypertarget{351-thay-ux111ux1ed5i-thux1b0ux1edbc-ux111o}{%
\subsubsection{3.5.1 Thay đổi thước
đo}\label{351-thay-ux111ux1ed5i-thux1b0ux1edbc-ux111o}}

\begin{itemize}
\tightlist
\item
  Biến phụ thuộc: thay \(\ln(LP)\) bằng ROS, sales growth, employment
  growth.
\item
  Biến độc lập: thay FSTS bằng export status (binary), foreign sales
  count.
\item
  Biến moderator: thay DAI\_thin bằng DAI\_rich (cho subset 2024+).
\end{itemize}

\hypertarget{352-mux1eabu-phux1ee5}{%
\subsubsection{3.5.2 Mẫu phụ}\label{352-mux1eabu-phux1ee5}}

Loại trừ micro firms (\textless5 lao động); chỉ giữ SMEs; chỉ giữ
manufacturing; chỉ giữ exporters (FSTS \textgreater{} 0); subgroup theo
6 sub-regime ICRV.

\hypertarget{353-outliers}{%
\subsubsection{3.5.3 Outliers}\label{353-outliers}}

Winsorization 1\% và 5\% đối với labor productivity và employment, áp
dụng trong từng country-year.

\hypertarget{354-multicollinearity-vuxe0-heteroscedasticity}{%
\subsubsection{3.5.4 Multicollinearity và
heteroscedasticity}\label{354-multicollinearity-vuxe0-heteroscedasticity}}

\begin{itemize}
\tightlist
\item
  VIF cho tất cả mô hình; ngưỡng VIF \textless{} 5 (Hair et al., 2019).
\item
  Breusch--Pagan test (Breusch \& Pagan, 1979).
\end{itemize}

\hypertarget{355-sensitivity-analysis-cho-meta-analysis}{%
\subsubsection{3.5.5 Sensitivity analysis cho
meta-analysis}\label{355-sensitivity-analysis-cho-meta-analysis}}

\begin{itemize}
\tightlist
\item
  Leave-one-out: loại từng nghiên cứu và tái tính pooled effect.
\item
  Trim-and-fill (Duval \& Tweedie, 2000) cho publication bias.
\end{itemize}

\hypertarget{36-tuxf3m-tux1eaft-phux1b0ux1a1ng-phuxe1p-kux1ebf-thux1eeba-vuxe0-ux111uxf3ng-guxf3p-mux1edbi}{%
\subsection{3.6 Tóm tắt phương pháp: Kế thừa và đóng góp
mới}\label{36-tuxf3m-tux1eaft-phux1b0ux1a1ng-phuxe1p-kux1ebf-thux1eeba-vuxe0-ux111uxf3ng-guxf3p-mux1edbi}}

\begin{longtable}[]{@{}lll@{}}
\toprule\noalign{}
Thành phần & Cơ sở kế thừa & Đóng góp mới \\
\midrule\noalign{}
\endhead
\bottomrule\noalign{}
\endlastfoot
Mixed design & Borenstein et al. (2009); Hunter \& Schmidt (2004) & Áp
dụng cho \textbf{47 châu Á + Pacific economies} trong một luận án \\
Meta-analysis 1982--2026 & Bausch \& Krist (2007); Kirca et al. (2012);
Marano et al. (2016) & Mở rộng coverage, bổ sung digital và 6 sub-regime
ICRV moderators \\
Pool empirical & Do \& Phan (2026a) --- 17 nước châu Á mới nổi & Mở rộng
từ 17 → \textbf{47 nước} với 101.185 doanh nghiệp · 108 cặp quốc gia ×
năm \\
Labor productivity DV & Bloom et al. (2012); Hsieh \& Klenow (2009); Do
\& Phan (2026a) & Lập luận rõ về firm performance đa chiều trong bối
cảnh WBES \\
FSTS + FSTS² + FSTS³ & Lu \& Beamish (2004); Hitt et al. (1997); Do \&
Phan (2026g) & Đưa vào cùng three-way moderation framework \\
TCI vs DAI & Bhandari et al. (2023); Verhoef et al. (2021) &
Non-overlapping construct purity protocol cho WBES; digital shield mở
rộng (Do \& Phan, 2026a) \\
Lind--Mehlum U-test & Lind \& Mehlum (2010); Haans et al. (2016) & Áp
dụng cho từng country sample và multi-country \\
Three-way moderation & Aiken \& West (1991); Dawson (2014) & digital ×
institutional × manager mới cho châu Á và Pacific \\
Temporal heterogeneity & Wu et al. (2022); Xiao et al. (2013); Do \&
Phan (2026g) & China 2012--2024 stability test (turning point
\textasciitilde47,8\% FSTS) \\
Robust SE & Long \& Ervin (2000); White (1980) & Chuẩn \\
Publication bias & Egger et al. (1997); Begg \& Mazumdar (1994) &
Chuẩn \\
\end{longtable}

\hypertarget{37-ux111ux1ea1o-ux111ux1ee9c-nghiuxean-cux1ee9u}{%
\subsection{3.7 Đạo đức nghiên
cứu}\label{37-ux111ux1ea1o-ux111ux1ee9c-nghiuxean-cux1ee9u}}

Luận án sử dụng dữ liệu thứ cấp từ WBES, bộ dữ liệu công khai được World
Bank cung cấp miễn phí cho mục đích nghiên cứu (World Bank, n.d.). Tất
cả các citation theo APA 7th để ghi nhận đầy đủ công sức của các tác giả
trước, tránh đạo văn. Code và dữ liệu phân tích được lưu trữ trên các
nhánh chuyên biệt của repository để bảo đảm khả năng tái lập kết quả.

\begin{center}\rule{0.5\linewidth}{0.5pt}\end{center}

\hypertarget{chux1b0ux1a1ng-4-kux1ebft-quux1ea3-nghiuxean-cux1ee9u}{%
\section{CHƯƠNG 4: KẾT QUẢ NGHIÊN
CỨU}\label{chux1b0ux1a1ng-4-kux1ebft-quux1ea3-nghiuxean-cux1ee9u}}

\begin{center}\rule{0.5\linewidth}{0.5pt}\end{center}

\hypertarget{giux1edbi-thiux1ec7u-chux1b0ux1a1ng}{%
\subsection{Giới thiệu
chương}\label{giux1edbi-thiux1ec7u-chux1b0ux1a1ng}}

Chương 4 trình bày toàn bộ kết quả nghiên cứu của luận án theo logic đi
từ bức tranh thực trạng mô tả (Mục 4.1) đến bằng chứng tổng hợp ở cấp
meta-analytic (Mục 4.2), sau đó đến bằng chứng quốc gia cụ thể và cuối
cùng là kiểm định toàn mẫu đa quốc gia. Cấu trúc trình bày gồm tám phần
chính, bắt đầu bằng thực trạng quốc tế hóa và hiệu quả hoạt động kinh
doanh tại châu Á theo các phân nhóm con thể chế ICRV, tiếp theo là: (ii)
kết quả meta-analysis ba tầng (P6); (iii) bối cảnh thể chế tiên
tiến---đổi mới tại Singapore (P4); (iv) bối cảnh thể chế trung bình cao
tại Trung Quốc (P5); (v) bối cảnh chuyển đổi bậc thấp---trung tại Việt
Nam (P3); (vi) kiểm định toàn mẫu châu Á đa bối cảnh (P7); (vii) trường
hợp biên là các nền kinh tế đảo nhỏ đang phát triển ở Thái Bình Dương
(P8); và (viii) tổng hợp kết quả theo hệ giả thuyết.

\begin{center}\rule{0.5\linewidth}{0.5pt}\end{center}

\hypertarget{41-thux1ef1c-trux1ea1ng-quux1ed1c-tux1ebf-huxf3a-vuxe0-hiux1ec7u-quux1ea3-houx1ea1t-ux111ux1ed9ng-kinh-doanh-tux1ea1i-chuxe2u-uxe1}{%
\subsection{4.1 Thực trạng quốc tế hóa và hiệu quả hoạt động kinh doanh
tại châu
Á}\label{41-thux1ef1c-trux1ea1ng-quux1ed1c-tux1ebf-huxf3a-vuxe0-hiux1ec7u-quux1ea3-houx1ea1t-ux111ux1ed9ng-kinh-doanh-tux1ea1i-chuxe2u-uxe1}}

\hypertarget{411-bux1ee9c-tranh-tux1ed5ng-thux1ec3-vux1ec1-hiux1ec7u-quux1ea3-doanh-nghiux1ec7p-chuxe2u-uxe1}{%
\subsubsection{4.1.1 Bức tranh tổng thể về hiệu quả doanh nghiệp châu
Á}\label{411-bux1ee9c-tranh-tux1ed5ng-thux1ec3-vux1ec1-hiux1ec7u-quux1ea3-doanh-nghiux1ec7p-chuxe2u-uxe1}}

Phân tích mô tả trên pool 101.185 doanh nghiệp từ 47 nền kinh tế châu Á
và Thái Bình Dương trong giai đoạn 2009--2025 cho thấy bức tranh hiệu
quả hoạt động kinh doanh có sự phân tán lớn và không hội tụ. Năng suất
lao động, đo bằng ln(doanh thu/lao động), có độ lệch chuẩn dao động từ
0,49 (Advanced tài nguyên dẫn dắt) đến 1,36 (Frontier), với tỷ số
P90/P10 tăng đơn điệu theo chiều từ nhóm thể chế cao đến nhóm thể chế
thấp.

\textbf{Bảng 4.1.} \emph{Phân tán năng suất lao động và đặc điểm quốc tế
hóa theo phân nhóm con ICRV (n = 101.185 doanh nghiệp, 2009--2025).}

\begin{longtable}[]{@{}llllll@{}}
\toprule\noalign{}
Phân nhóm con ICRV & n (ước tính) & sd log LP & FSTS (\%) & Xuất khẩu
(\%) & FDI≥10\% (\%) \\
\midrule\noalign{}
\endhead
\bottomrule\noalign{}
\endlastfoot
Nhóm I --- Advanced đổi mới (SG, HK, KOR, TWN, ISR) &
\textasciitilde4.220 & \textbf{1,03} & 10,2 & 23,0 & 11,1 \\
Nhóm II --- Advanced tài nguyên (GCC) & \textasciitilde1.932 &
\textbf{0,49} & \textasciitilde3 & \textasciitilde8 & --- \\
Nhóm III --- Upper-middle (CHN, MYS, THA, KAZ...) &
\textasciitilde16.693 & 1,29 & 10,3 & 21,7 & 8,4 \\
Nhóm IV --- Emerging (VNM, IDN, PHL, IND...) & \textasciitilde47.803 &
1,24 & 8,6 & 15,5 & 4,7 \\
Nhóm V --- Frontier (BGD, PAK, LAO, KHM...) & \textasciitilde28.678 &
1,36 & 10,1 & 16,6 & 5,9 \\
Nhóm VI --- SIDS Thái Bình Dương (FJI, PNG, SLB...) &
\textasciitilde1.371 & 1,32 & 6,3 & 16,3 & 23,5 \\
\textbf{Tổng} & \textbf{101.185} & --- & --- & --- & --- \\
\end{longtable}

\emph{Nguồn: Phân tích của tác giả từ World Bank Enterprise Surveys
(WBES), 2009--2025.}

Phát hiện quan trọng nhất từ thực trạng mô tả là sự phân biệt nội bộ
trong nhóm Advanced: tỷ số phân tán giữa nhóm đổi mới sáng tạo dẫn dắt
(Singapore, Hàn Quốc, Đài Loan, sd ≈ 1,03) và nhóm tài nguyên dẫn dắt
(Saudi Arabia, Qatar, Kuwait, sd ≈ 0,49) xấp xỉ 2,1 lần, phản ánh hai cơ
chế tăng trưởng và phân phối hoàn toàn khác nhau. Phân nhóm con tài
nguyên dẫn dắt có độ phân tán năng suất thấp nhất toàn bộ pool, biểu
hiện của kinh tế nhà nước tô (rentier economy) nơi phân phối lại từ xuất
khẩu dầu mỏ thu hẹp khoảng cách năng suất giữa doanh nghiệp, nhưng không
nhất thiết phản ánh hiệu quả doanh nghiệp đáng kể. Điều này gợi ý rằng
việc gộp hai loại Advanced vào một nhóm duy nhất (như các nghiên cứu
trước thường làm) sẽ che khuất cơ chế quan trọng này.

\hypertarget{412-thux1ef1c-trux1ea1ng-quux1ed1c-tux1ebf-huxf3a-phuxe2n-cux1ef1c-tham-gia-xuux1ea5t-khux1ea9u}{%
\subsubsection{4.1.2 Thực trạng quốc tế hóa, phân cực tham gia xuất
khẩu}\label{412-thux1ef1c-trux1ea1ng-quux1ed1c-tux1ebf-huxf3a-phuxe2n-cux1ef1c-tham-gia-xuux1ea5t-khux1ea9u}}

Trung vị FSTS bằng 0\% trong tất cả các phân nhóm, chỉ 15--23\% doanh
nghiệp tham gia xuất khẩu tùy theo bối cảnh. Phân phối FSTS có dạng phân
cực mạnh: đại đa số không xuất khẩu, trong khi một thiểu số nhỏ có tỷ lệ
xuất khẩu rất cao. Phân cực này đặt ra vấn đề phương pháp quan trọng cho
nghiên cứu I→P: kết quả hồi quy trên toàn mẫu bị kéo bởi hành vi của
thiểu số xuất khẩu tích cực.

Cường độ quốc tế hóa (FSTS trung bình, tính trên toàn mẫu gồm cả doanh
nghiệp không xuất khẩu) dao động từ 6,3\% (SIDS) đến 10,3\%
(Upper-middle), với mức trung bình khu vực khoảng 8,8\%. Điều đáng chú ý
là Upper-middle và Advanced có cùng FSTS trung bình
(\textasciitilde10\%) nhưng cơ chế xuất khẩu hoàn toàn khác: Advanced
chủ yếu dịch vụ và hàng hóa hàm lượng tri thức cao; Upper-middle chủ yếu
sản xuất theo chuỗi giá trị toàn cầu với biên lợi nhuận khác nhau.

\hypertarget{413-thux1ef1c-trux1ea1ng-nux103ng-lux1ef1c-sux1ed1-vuxe0-cuxf4ng-nghux1ec7}{%
\subsubsection{4.1.3 Thực trạng năng lực số và công
nghệ}\label{413-thux1ef1c-trux1ea1ng-nux103ng-lux1ef1c-sux1ed1-vuxe0-cuxf4ng-nghux1ec7}}

Phân tích năm thành phần năng lực (đổi mới sản phẩm, đổi mới quy trình,
R\&D, ISO, website) theo phân nhóm con ICRV cho thấy sự đa dạng đáng kể.

\textbf{Bảng 4.2.} \emph{Đổi mới sáng tạo và chấp nhận số theo phân nhóm
con ICRV (\%).}

\begin{longtable}[]{@{}llllll@{}}
\toprule\noalign{}
Phân nhóm con & Sản phẩm mới & Quy trình mới & R\&D & ISO & Website (DAI
Tier-1) \\
\midrule\noalign{}
\endhead
\bottomrule\noalign{}
\endlastfoot
Advanced & 22,3 & 52,3 & 16,7 & 29,9 & 59,3 \\
Upper-middle & 26,7 & 71,7 & 21,0 & 31,4 & 56,9 \\
Emerging & 17,5 & 65,2 & 16,4 & 24,9 & 49,2 \\
Frontier & 23,1 & 68,9 & 14,2 & 20,7 & 38,0 \\
SIDS Thái Bình Dương & \textbf{41,5} & 65,1 & 11,8 & 16,5 &
\textbf{58,9} \\
\end{longtable}

\emph{Nguồn: Phân tích của tác giả từ WBES, 2009--2025.}

Phát hiện đáng chú ý là pattern "nhảy vọt số" (digital leapfrogging) ở
SIDS: tỷ lệ website 58,9\% vượt cả nhóm Emerging (49,2\%) dù GNI thấp
hơn đáng kể, phản ánh thực tế rằng số hóa bề mặt (website) ở SIDS là
công cụ sinh tồn chứ không phải công cụ tối ưu hóa hiệu quả. Kết quả này
cũng hỗ trợ lập luận lý thuyết về sự phân biệt giữa DAI (chấp nhận số bề
mặt) và TCI (năng lực công nghệ chiều sâu): SIDS có DAI cao nhưng TCI
thấp nhất, và đây là nhóm duy nhất có quan hệ I→P âm đơn điệu.

Hiện tượng tương quan âm giữa DAI Tier-1 và năng suất trong nhóm
Advanced (-0,129) là ảo ảnh thống kê do bão hòa Tier-1: hầu hết doanh
nghiệp Advanced đã có website từ lâu, khiến biến nhị phân này không còn
phân biệt được doanh nghiệp hiệu quả và kém hiệu quả trong nhóm đó. Khi
loại trừ doanh nghiệp ICT, hệ số âm biến mất, xác nhận đây là vấn đề đo
lường chứ không phải quan hệ nhân quả thực.

\hypertarget{414-bux1ed1n-ruxe0o-cux1ea3n-cux1ea5u-truxfac-ux1ea3nh-hux1b0ux1edfng-ux111ux1ebfn-quux1ed1c-tux1ebf-huxf3a-vuxe0-hiux1ec7u-quux1ea3}{%
\subsubsection{4.1.4 Bốn rào cản cấu trúc ảnh hưởng đến quốc tế hóa và
hiệu
quả}\label{414-bux1ed1n-ruxe0o-cux1ea3n-cux1ea5u-truxfac-ux1ea3nh-hux1b0ux1edfng-ux111ux1ebfn-quux1ed1c-tux1ebf-huxf3a-vuxe0-hiux1ec7u-quux1ea3}}

Phân tích thực trạng xác định bốn rào cản cấu trúc phổ biến nhất ảnh
hưởng đến khả năng quốc tế hóa và hiệu quả doanh nghiệp tại châu Á: (i)
tiếp cận tài chính hạn chế, đặc biệt nghiêm trọng ở Frontier và
Emerging; (ii) thiếu hụt lao động có kỹ năng; (iii) cạnh tranh từ doanh
nghiệp phi chính thức; và (iv) nguồn cung điện không đáng tin cậy. Bốn
rào cản này hình thành "institutional voids" (Khanna \& Palepu, 2010)
tạo ra chi phí giao dịch bổ sung cho doanh nghiệp khi mở rộng hoạt động
ra ngoài biên giới. Sự hiện diện đồng thời của nhiều voids tại Frontier
và SIDS giải thích tại sao quan hệ I→P ở các nhóm này yếu hoặc âm.

\begin{center}\rule{0.5\linewidth}{0.5pt}\end{center}

\hypertarget{42-kux1ebft-quux1ea3-tux1ed5ng-hux1ee3p-ux111ux1ecbnh-lux1b0ux1ee3ng--meta-analysis-ba-tux1ea7ng-p6}{%
\subsection{4.2 Kết quả tổng hợp định lượng --- Meta-analysis ba tầng
(P6)}\label{42-kux1ebft-quux1ea3-tux1ed5ng-hux1ee3p-ux111ux1ecbnh-lux1b0ux1ee3ng--meta-analysis-ba-tux1ea7ng-p6}}

\hypertarget{421-chux1ec9-sux1ed1-hiux1ec7u-ux1ee9ng-tux1ed5ng-hux1ee3p}{%
\subsubsection{4.2.1 Chỉ số hiệu ứng tổng
hợp}\label{421-chux1ec9-sux1ed1-hiux1ec7u-ux1ee9ng-tux1ed5ng-hux1ee3p}}

Phân tích meta-analytic tổng hợp ba tầng (three-level multilevel
meta-analysis --- MARA) được tiến hành trên tập dữ liệu gồm \(k = 238\)
nghiên cứu với tổng cộng hàng trăm cặp effect-size đa dạng theo vùng địa
lý, giai đoạn và phương pháp. Kết quả cho thấy hiệu ứng tổng hợp (pooled
correlation) là \(r = 0{,}074\), với khoảng tin cậy 95\% dương và có ý
nghĩa thống kê. Giá trị \(r\) này khẳng định bức tranh tổng quát từ các
meta-analyses trước đây của Bausch và Krist (2007), Kirca et al. (2012),
Marano et al. (2016) và Arte và Larimo (2022): mối quan hệ tổng hợp giữa
quốc tế hóa và hiệu quả hoạt động là dương nhưng có biên độ khiêm tốn.

Đặc biệt đáng lưu ý, kết quả P6 hoàn toàn nhất quán với baseline
meta-analysis đã công bố tại hội nghị ICBEF 2024 của Do và Phan (2024),
với \(k = 113\) nghiên cứu và \(r = 0{,}07\), sự hội tụ này tạo ra bằng
chứng chéo vững chắc khi mở rộng pool lên gần gấp đôi số nghiên cứu.

\hypertarget{422-phuxe2n-tuxedch-dux1ecb-biux1ec7t-ba-tux1ea7ng}{%
\subsubsection{4.2.2 Phân tích dị biệt ba
tầng}\label{422-phuxe2n-tuxedch-dux1ecb-biux1ec7t-ba-tux1ea7ng}}

Kết quả quan trọng nhất của phân tích MARA không phải là giá trị điểm
của \(r\) mà là cấu trúc của dị biệt (\(I^2\)). Tổng \(I^2 = 62{,}4\%\),
cho thấy mức dị biệt đáng kể trong tập hợp 238 nghiên cứu không phải do
sai số lấy mẫu. Quan trọng hơn, khi phân tách theo cấu trúc ba tầng:

\begin{itemize}
\tightlist
\item
  \textbf{Tầng 2} (between-effect within-study, tức dị biệt giữa các
  effect-size trong cùng một nghiên cứu): chiếm \(54{,}1\%\) tổng phương
  sai, đây là nguồn dị biệt chủ đạo.
\item
  \textbf{Tầng 3} (between-study, tức dị biệt giữa các nghiên cứu): chỉ
  chiếm \(8{,}4\%\).
\end{itemize}

Phát hiện này có hàm ý lý thuyết sâu sắc: heterogeneity trong I→P
literature xuất phát chủ yếu từ sự biến thiên \textbf{bên trong các
nghiên cứu} (ví dụ: cùng một mẫu doanh nghiệp nhưng dùng các thước đo
khác nhau cho cùng cấu trúc cho ra kết quả khác nhau), chứ không phải
giữa các nghiên cứu với nhau. Điều này gợi ý rằng sự không nhất quán
trong literature I→P phần lớn là vấn đề \textbf{đo lường và bối cảnh},
không phải là bằng chứng về sự thiếu nhất quán thực sự trong quan hệ
nhân quả giữa quốc tế hóa và hiệu quả.

\hypertarget{423-ux111iux1ec1u-tiux1ebft-bux1edfi-chux1ebf-ux111ux1ed9-thux1ec3-chux1ebf--icrv-moderation}{%
\subsubsection{4.2.3 Điều tiết bởi chế độ thể chế --- ICRV
moderation}\label{423-ux111iux1ec1u-tiux1ebft-bux1edfi-chux1ebf-ux111ux1ed9-thux1ec3-chux1ebf--icrv-moderation}}

Kiểm định điều tiết bằng nhóm chế độ thể chế ICRV (Institutional Context
Regime Variation) cho kết quả có ý nghĩa thống kê: \(Q_M = 17{,}35\),
\(df = 4\), \(p = {,}002\). Điều này xác nhận rằng chế độ thể chế, được
phân loại thành sáu nhóm từ Advanced Innovation-Driven (Nhóm I:
Singapore, Hong Kong, Đài Loan, Hàn Quốc) đến SIDS Thái Bình Dương (Nhóm
VI), là moderator có ý nghĩa thực sự cho mối quan hệ tổng hợp I→P trong
văn liệu. Kiểm định \(Q_M\) vượt ngưỡng tới hạn (\(p = {,}002\)), hàm ý
rằng sự khác biệt giữa các nhóm ICRV không thể giải thích bằng sai số
lấy mẫu.

Gradient theo ICRV hiện diện rõ ràng trong kết quả: các nền kinh tế
Advanced Innovation-Driven có hiệu ứng I→P lớn hơn so với các nền kinh
tế Frontier hay SIDS. Kết quả này nhất quán với cơ chế lý thuyết từ
Institutional Theory, môi trường thể chế chất lượng cao tạo điều kiện để
doanh nghiệp chuyển hóa quốc tế hóa thành hiệu quả dễ dàng hơn (North,
1990; Marano et al., 2016).

\hypertarget{424-kiux1ec3m-ux111ux1ecbnh-publication-bias}{%
\subsubsection{4.2.4 Kiểm định publication
bias}\label{424-kiux1ec3m-ux111ux1ecbnh-publication-bias}}

Phân tích publication bias gồm ba kiểm định bổ sung nhau. \textbf{Kiểm
định Egger} cho kết quả \(b = 0{,}475\) (\(p = {,}057\)), sát ngưỡng
\(\alpha = 0{,}05\) nhưng không đạt ý nghĩa thống kê theo tiêu chuẩn
thông thường. \textbf{Kiểm định Begg--Mazumdar} (Begg \& Mazumdar, 1994)
cho kết quả \(\tau = -0{,}134\) (\(p = {,}0007\)), có ý nghĩa thống kê
rõ ràng, cho thấy các nghiên cứu với sai số chuẩn lớn hơn (mẫu nhỏ hơn)
có xu hướng báo cáo effect size thấp hơn, phù hợp với lệch lạc công bố
ưu tiên kết quả dương tính. \textbf{Phân tích trim-and-fill} điều chỉnh
ước lượng từ \(r = 0{,}074\) xuống \(r = 0{,}035\) (với \(k = 58\)
nghiên cứu được imputed), mức điều chỉnh đáng kể nhưng không đổi chiều
hướng tích cực của quan hệ tổng hợp. \textbf{Fail-safe \(N\) = 45.848}
vượt xa ngưỡng tới hạn \(5k + 10 = 1{,}200\), khẳng định publication
bias không thể bác bỏ hoàn toàn kết quả tổng hợp. Tổng hợp lại, bằng
chứng publication bias được xác nhận qua kiểm định Begg--Mazumdar có ý
nghĩa (\(p = {,}0007\)): pooled effect thực \(r = 0{,}074\) có thể bị
thổi phồng so với hiệu ứng dân số thực, với ước lượng bảo thủ hơn là
\(r = 0{,}035\).

\begin{center}\rule{0.5\linewidth}{0.5pt}\end{center}

\hypertarget{43-kux1ebft-quux1ea3-bux1ed1i-cux1ea3nh-thux1ec3-chux1ebf-tiuxean-tiux1ebfn-ux111ux1ed5i-mux1edbi-singapore-p4}{%
\subsection{4.3 Kết quả bối cảnh thể chế tiên tiến, Đổi mới: Singapore
(P4)}\label{43-kux1ebft-quux1ea3-bux1ed1i-cux1ea3nh-thux1ec3-chux1ebf-tiuxean-tiux1ebfn-ux111ux1ed5i-mux1edbi-singapore-p4}}

\hypertarget{431-xuxe1c-nhux1eadn-quan-hux1ec7-phi-tuyux1ebfn-chux1eef-u-ngux1b0ux1ee3c}{%
\subsubsection{4.3.1 Xác nhận quan hệ phi tuyến chữ U
ngược}\label{431-xuxe1c-nhux1eadn-quan-hux1ec7-phi-tuyux1ebfn-chux1eef-u-ngux1b0ux1ee3c}}

Phân tích trên mẫu doanh nghiệp Singapore từ dữ liệu WBES với
\(N = 237\) doanh nghiệp xuất khẩu xác nhận mối quan hệ phi tuyến
(inverted U-shape) giữa FSTS và hiệu quả hoạt động ở mức có ý nghĩa
thống kê (\(p < {,}05\)). Mô hình bậc hai với biến FSTS và FSTS\(^2\)
cho hệ số bậc hai âm, phù hợp với dạng chữ U ngược như kỳ vọng từ H1
trong khung lý thuyết CDCM.

Điểm turning point ước lượng được ở mức \textbf{TP ≈ 88,6\% FSTS}, đây
là điểm turning point cao nhất trong tất cả các bối cảnh ICRV được phân
tích trong luận án. Khoảng tin cậy của turning point rộng (CI: {[}53\%,
253\%{]}), phản ánh hạn chế về quy mô mẫu và thống kê chính xác, nhưng
điểm ước lượng trung tâm cho thấy rằng doanh nghiệp Singapore có thể duy
trì đà tăng trưởng hiệu quả khi quốc tế hóa ở mức độ rất cao trước khi
đạt đỉnh và suy giảm.

\hypertarget{432-ux111iux1ec1u-tiux1ebft-bux1edfi-digital-adoption-index-dai}{%
\subsubsection{4.3.2 Điều tiết bởi Digital Adoption Index
(DAI)}\label{432-ux111iux1ec1u-tiux1ebft-bux1edfi-digital-adoption-index-dai}}

Mô hình M5 bổ sung tương tác \(FSTS \times DAI\) (bao gồm cả Tier-1 và
Tier-2 của chỉ số DAI) cho thấy hiệu ứng khuếch đại (amplification) tại
mức FSTS cao: hệ số tương tác \(FSTS^2 \times DAI < 0\) và có ý nghĩa
thống kê (\(p < {,}05\)). Cụ thể, ở mức độ quốc tế hóa cao, doanh nghiệp
Singapore có năng lực chấp nhận số (digital adoption) cao hơn duy trì
hiệu quả tốt hơn so với doanh nghiệp có mức DAI thấp hơn, cho thấy DAI
đóng vai trò nguồn lực bổ trợ (situational resource) làm chậm đà suy
giảm phía bên phải của đường phi tuyến.

\hypertarget{433-giux1edbi-hux1ea1n-suy-luux1eadn-vuxe0-cux1ea3nh-buxe1o-thux1ed1ng-kuxea}{%
\subsubsection{4.3.3 Giới hạn suy luận và cảnh báo thống
kê}\label{433-giux1edbi-hux1ea1n-suy-luux1eadn-vuxe0-cux1ea3nh-buxe1o-thux1ed1ng-kuxea}}

Cần nhấn mạnh rằng kết quả Singapore là \textbf{ranh giới suy luận}
(inferential bounds) do hạn chế nghiêm trọng về quy mô mẫu: \(N = 237\)
quan sát, chỉ gồm các doanh nghiệp xuất khẩu, dẫn đến power thống kê ước
tính chỉ đạt khoảng 16\%. Luận án diễn giải kết quả Singapore như
\emph{bằng chứng phù hợp nhưng chưa kết luận} (consistent but
underpowered evidence), không bác bỏ giả thuyết nhưng cũng không tuyên
bố xác nhận dứt khoát.

\begin{center}\rule{0.5\linewidth}{0.5pt}\end{center}

\hypertarget{44-kux1ebft-quux1ea3-bux1ed1i-cux1ea3nh-thux1ec3-chux1ebf-trung-buxecnh-cao-trung-quux1ed1c-20122024-p5}{%
\subsection{4.4 Kết quả bối cảnh thể chế trung bình cao: Trung Quốc,
2012--2024
(P5)}\label{44-kux1ebft-quux1ea3-bux1ed1i-cux1ea3nh-thux1ec3-chux1ebf-trung-buxecnh-cao-trung-quux1ed1c-20122024-p5}}

\hypertarget{441-xuxe1c-nhux1eadn-quan-hux1ec7-phi-tuyux1ebfn-vuxe0-turning-point}{%
\subsubsection{4.4.1 Xác nhận quan hệ phi tuyến và turning
point}\label{441-xuxe1c-nhux1eadn-quan-hux1ec7-phi-tuyux1ebfn-vuxe0-turning-point}}

Phân tích hai giai đoạn trên dữ liệu WBES Trung Quốc (Nhóm III ---
Upper-middle) giai đoạn 2012--2024 xác nhận dạng quan hệ phi tuyến chữ U
ngược giữa FSTS và hiệu quả hoạt động doanh nghiệp. Điểm turning point
được ước lượng tại:

\begin{itemize}
\tightlist
\item
  \textbf{Năm 2012}: TP = 49,37\% FSTS
\item
  \textbf{Năm 2024}: TP = 47,19\% FSTS
\item
  \textbf{Pooled (toàn mẫu)}: TP = 48,78\% FSTS
\end{itemize}

Các con số này cho thấy doanh nghiệp Trung Quốc đạt hiệu quả tối ưu khi
cường độ quốc tế hóa (FSTS) ở mức gần 49\%, một mức độ tập trung quốc tế
vừa phải trong bối cảnh nền kinh tế nội địa khổng lồ của Trung Quốc cũng
cung cấp cơ hội tăng trưởng song song.

\hypertarget{442-kiux1ec3m-ux111ux1ecbnh-tuxednh-ux1ed5n-ux111ux1ecbnh-theo-thux1eddi-gian--temporal-stability}{%
\subsubsection{4.4.2 Kiểm định tính ổn định theo thời gian --- Temporal
Stability}\label{442-kiux1ec3m-ux111ux1ecbnh-tuxednh-ux1ed5n-ux111ux1ecbnh-theo-thux1eddi-gian--temporal-stability}}

Luận án đặt câu hỏi thực nghiệm quan trọng: liệu hình dạng phi tuyến và
điểm turning point có thay đổi qua thập kỷ 2012--2024, giai đoạn Trung
Quốc chứng kiến nhiều thay đổi chính sách thương mại, căng thẳng địa
chính trị và chuyển đổi số mạnh mẽ? Kiểm định Paternoster cho kết quả
\(z = +0{,}82\), \(p = {,}412\), không có ý nghĩa thống kê. Điều này xác
nhận rằng \textbf{điểm turning point ổn định qua thời gian}: sự dịch
chuyển từ TP = 49,37\% (2012) xuống TP = 47,19\% (2024) là không đáng kể
về mặt thống kê.

Kiểm định tương tác thời gian (Temporal H6): \(F = 1{,}83\),
\(p = {,}176\), không có ý nghĩa thống kê, nghĩa là hình dạng tổng thể
của quan hệ phi tuyến không biến đổi đáng kể qua giai đoạn phân tích.
Phát hiện này cho thấy trong bối cảnh Upper-middle như Trung Quốc, cơ
chế căn bản của quan hệ I→P duy trì tính ổn định đáng kể, có thể phản
ánh sự bù trừ giữa chi phí thích ứng tăng lên do chính sách thương mại
khó đoán và năng lực tổ chức ngày càng trưởng thành của các doanh nghiệp
xuất khẩu Trung Quốc.

\hypertarget{443-vai-truxf2-cux1ee7a-tci-vuxe0-dai}{%
\subsubsection{4.4.3 Vai trò của TCI và
DAI}\label{443-vai-truxf2-cux1ee7a-tci-vuxe0-dai}}

Tương tác \(FSTS \times TCI\) không đạt ý nghĩa thống kê (null
moderation, \(p > {,}10\)), cho thấy Technological Capability Index
không phát huy vai trò điều tiết đáng kể ở bối cảnh Trung Quốc theo mô
hình P5. Tương tự, DAI cũng không cho thấy điều tiết có ý nghĩa trong
bối cảnh này. Kết quả null về moderation tại Trung Quốc có thể phản ánh
thực tế rằng năng lực công nghệ đã được phân phối tương đối đồng đều hơn
trong mẫu doanh nghiệp xuất khẩu Trung Quốc, làm giảm phương sai đủ để
phát hiện hiệu ứng điều tiết.

\begin{center}\rule{0.5\linewidth}{0.5pt}\end{center}

\hypertarget{45-kux1ebft-quux1ea3-bux1ed1i-cux1ea3nh-chuyux1ec3n-ux111ux1ed5i-bux1eadc-thux1ea5ptrung-viux1ec7t-nam-p3}{%
\subsection{4.5 Kết quả bối cảnh chuyển đổi bậc thấp---trung: Việt Nam
(P3)}\label{45-kux1ebft-quux1ea3-bux1ed1i-cux1ea3nh-chuyux1ec3n-ux111ux1ed5i-bux1eadc-thux1ea5ptrung-viux1ec7t-nam-p3}}

\hypertarget{451-xuxe1c-nhux1eadn-quan-hux1ec7-phi-tuyux1ebfn-vuxe0-hux1ec7-sux1ed1-hux1ed3i-quy}{%
\subsubsection{4.5.1 Xác nhận quan hệ phi tuyến và hệ số hồi
quy}\label{451-xuxe1c-nhux1eadn-quan-hux1ec7-phi-tuyux1ebfn-vuxe0-hux1ec7-sux1ed1-hux1ed3i-quy}}

Phân tích hồi quy bảng (panel regression) trên ba làn sóng khảo sát WBES
Việt Nam (2009, 2015, 2023) xác nhận dạng phi tuyến chữ U ngược. Mô hình
M4, mô hình chính, cho kết quả:

\[\hat{\beta}(\text{FSTS}_c) = +0{,}431, \quad \hat{\beta}(\text{FSTS}_c^2) = -0{,}553\]

Cả hai hệ số đều có ý nghĩa thống kê, với hệ số bậc nhất dương và bậc
hai âm xác nhận đường cong chữ U ngược. Kiểm định hình thức bằng phương
pháp Lind--Mehlum (U-test) cho kết quả \(p < {,}001\) --- bằng chứng
thống kê mạnh nhất trong tất cả các mẫu quốc gia của luận án.

\hypertarget{452-ux111iux1ec3m-turning-point-theo-luxe0n-suxf3ng}{%
\subsubsection{4.5.2 Điểm turning point theo làn
sóng}\label{452-ux111iux1ec3m-turning-point-theo-luxe0n-suxf3ng}}

Turning point được ước lượng riêng cho từng làn sóng và cho pooled:

\begin{longtable}[]{@{}ll@{}}
\toprule\noalign{}
Làn sóng & Turning Point \\
\midrule\noalign{}
\endhead
\bottomrule\noalign{}
\endlastfoot
2009 & 46,2\% FSTS \\
2015 & 39,3\% FSTS \\
2023 & 41,6\% FSTS \\
Pooled & 39,7\% FSTS \\
\end{longtable}

Điểm turning point trong khoảng \textbf{39--46\% FSTS} đặt Việt Nam ở
mức "thấp hơn đáng kể" so với Singapore (TP ≈ 88,6\%) và tương đương
hoặc thấp hơn Trung Quốc (TP ≈ 48,78\%). Kiểm định Paternoster so sánh
turning point giữa 2009 và 2015 cho kết quả \(z = 3{,}353\),
\(p < {,}001\) --- có ý nghĩa thống kê, cho thấy turning point đã dịch
chuyển đáng kể từ 46,2\% xuống 39,3\% giữa hai làn sóng.

\hypertarget{453-ux111iux1ec1u-tiux1ebft-bux1edfi-tci--tuxedch-cux1ef1c-vuxe0-cuxf3-uxfd-nghux129a}{%
\subsubsection{4.5.3 Điều tiết bởi TCI --- Tích cực và có ý
nghĩa}\label{453-ux111iux1ec1u-tiux1ebft-bux1edfi-tci--tuxedch-cux1ef1c-vuxe0-cuxf3-uxfd-nghux129a}}

Mô hình M6 với tương tác \(FSTS \times TCI\) cho kết quả \textbf{dương
và có ý nghĩa thống kê}: TCI làm nâng cao mặt bằng hiệu quả và khuếch
đại tác động tích cực của quốc tế hóa. Doanh nghiệp Việt Nam có năng lực
công nghệ tốt hơn (đo bằng TCI) đạt hiệu quả cao hơn từ quá trình quốc
tế hóa --- nhất quán với H2 trong khung lý thuyết CDCM.

\hypertarget{454-ux111iux1ec1u-tiux1ebft-bux1edfi-dai--khuxf4ng-cuxf3-uxfd-nghux129a}{%
\subsubsection{4.5.4 Điều tiết bởi DAI --- Không có ý
nghĩa}\label{454-ux111iux1ec1u-tiux1ebft-bux1edfi-dai--khuxf4ng-cuxf3-uxfd-nghux129a}}

Mô hình M5 với tương tác \(FSTS \times DAI\) (chỉ bao gồm Tier-1 DAI do
giới hạn đo lường WBES trong các làn sóng này) cho kết quả yếu dương
nhưng không có ý nghĩa thống kê. Kết quả này có thể phản ánh hạn chế về
đo lường: Tier-1 DAI chỉ nắm bắt mức độ chấp nhận số bề mặt (website,
email, online sales) chứ không đo lường được chiều sâu năng lực số. Với
dữ liệu WBES Việt Nam chỉ cung cấp Tier-1, không đủ phân biệt để phát
hiện hiệu ứng điều tiết tinh tế hơn.

\begin{center}\rule{0.5\linewidth}{0.5pt}\end{center}

\hypertarget{46-kux1ebft-quux1ea3-kiux1ec3m-ux111ux1ecbnh-touxe0n-mux1eabu-ux111a-bux1ed1i-cux1ea3nh-chuxe2u-uxe1-p7}{%
\subsection{4.6 Kết quả kiểm định toàn mẫu đa bối cảnh châu Á
(P7)}\label{46-kux1ebft-quux1ea3-kiux1ec3m-ux111ux1ecbnh-touxe0n-mux1eabu-ux111a-bux1ed1i-cux1ea3nh-chuxe2u-uxe1-p7}}

\hypertarget{461-mux1eabu-vuxe0-muxf4-huxecnh-chuxednh}{%
\subsubsection{4.6.1 Mẫu và mô hình
chính}\label{461-mux1eabu-vuxe0-muxf4-huxecnh-chuxednh}}

Phân tích P7 là kiểm định quy mô lớn nhất của luận án, kết hợp tất cả
các bối cảnh ICRV trong khu vực châu Á và Thái Bình Dương với
\textbf{\(N = 84.910\) quan sát} (mô hình M2 --- quadratic FSTS không có
biến kiểm soát). Khi bổ sung đầy đủ biến kiểm soát, mẫu hồi quy giảm
xuống \(N = 38.342\) quan sát (M3--M5), phản ánh missing data trong biến
sở hữu nước ngoài và tuổi doanh nghiệp ở các nền kinh tế nhỏ và sóng
WBES sớm. Mô hình M5 với country fixed effects và year fixed effects ---
cấu hình mạnh nhất về kiểm soát phương sai không quan sát
(\(\text{adjR}^2 = {,}677\)) --- ước lượng turning point tổng thể
(pooled):

\[\text{TP(M5)} = 40{,}0\% \text{ FSTS} \quad (p_{\text{LM}} < {,}001)\]

Turning point dao động nhất quán trong dải hẹp 33,8--40,0\% FSTS qua tất
cả các đặc tả M2--M9, và được xác nhận lại trong mô hình điều tiết ba
chiều đầy đủ M11 (\(\text{TP} = 34{,}6\%\), \(p_{\text{LM}} = {,}002\),
\(N = 29.840\)). Tính nhất quán này bác bỏ giả thuyết rằng chữ U ngược
chỉ là artefact của mẫu nhỏ hay thiếu biến kiểm soát.

\hypertarget{462-gradient-icrv--xuxe1c-nhux1eadn-h5}{%
\subsubsection{4.6.2 Gradient ICRV --- Xác nhận
H5}\label{462-gradient-icrv--xuxe1c-nhux1eadn-h5}}

Kết quả quan trọng nhất của P7 là xác nhận \textbf{gradient ICRV}:
turning point tăng dần khi chuyển từ nhóm thể chế mạnh hơn sang nhóm thể
chế yếu hơn:

\[\text{TP(Advanced Innovation)} \approx 28\% < \text{TP(Upper-middle)} < \text{TP(Frontier/SIDS)} \approx 55\%\]

Gradient này nhất quán với dự đoán của Institutional Theory: trong môi
trường thể chế mạnh (Nhóm I), doanh nghiệp hoàn thu chi phí quốc tế hóa
tại mức FSTS thấp hơn, tức là đạt đỉnh hiệu quả sớm hơn. Ngược lại, tại
các nền kinh tế thể chế yếu (Nhóm V--VI), doanh nghiệp cần mức độ quốc
tế hóa cao hơn để bù đắp chi phí giao dịch và bất định thể chế, nên
turning point xuất hiện muộn hơn --- ở mức FSTS cao hơn.

Kiểm định điều tiết tổng thể \(Q_M\) về ICRV có ý nghĩa thống kê cao,
nhất quán với kết quả meta-analysis trong P6 (\(Q_M = 17{,}35\),
\(p = {,}002\)). Hai nguồn bằng chứng độc lập --- meta-analytic và
primary empirical --- cùng chỉ ra gradient ICRV, tạo ra independent
confirmation có giá trị.

\hypertarget{463-ux111iux1ec1u-tiux1ebft-bux1edfi-tci-vuxe0-dai-trong-touxe0n-mux1eabu}{%
\subsubsection{4.6.3 Điều tiết bởi TCI và DAI trong toàn
mẫu}\label{463-ux111iux1ec1u-tiux1ebft-bux1edfi-tci-vuxe0-dai-trong-touxe0n-mux1eabu}}

\begin{itemize}
\tightlist
\item
  \textbf{TCI}: có ý nghĩa thống kê dương trên toàn mẫu --- một đơn vị
  độ lệch chuẩn tăng TCI nâng năng suất lao động khoảng 41\%
  (exp(0,344)−1). Tuy nhiên, các tương tác FSTS×TCI và FSTS²×TCI
  \textbf{không có ý nghĩa} trong mẫu 49 nền kinh tế đa dạng thể chế. H2
  \textbf{được xác nhận một phần}: TCI là yếu tố nâng mặt bằng năng suất
  nhất quán, nhưng vai trò uốn hình dạng đường cong I→P chỉ xuất hiện
  trong bối cảnh quốc gia cụ thể (P3 Việt Nam), không phổ quát xuyên
  toàn mẫu.
\item
  \textbf{DAI}: có ý nghĩa thống kê toàn mẫu --- hiệu ứng trực tiếp
  \(\hat{\beta}(\text{DAI}) = +0{,}155\) (\(p < {,}001\)) tương ứng lợi
  thế năng suất \textasciitilde17\% cho doanh nghiệp có website. Quan
  trọng hơn, cả hai tương tác điều tiết đều có ý nghĩa:
  \(\hat{\beta}(\text{FSTS} \times \text{DAI}) = -0{,}614\)
  (\(p < {,}001\)) và
  \(\hat{\beta}(\text{FSTS}^2 \times \text{DAI}) = +0{,}766\)
  (\(p < {,}01\)). Pattern âm rồi dương này cho thấy doanh nghiệp DAI
  cao có đường cong I→P phẳng hơn nhưng dịch lên --- tổn thất hiệu suất
  ít hơn khi quốc tế hóa cao. Tương tác DAI×ICRV (\(p = {,}012\)) xác
  nhận vai trò của DAI mạnh hơn ở các môi trường thể chế yếu hơn (nhất
  quán với CDCM).
\end{itemize}

\begin{center}\rule{0.5\linewidth}{0.5pt}\end{center}

\hypertarget{47-kux1ebft-quux1ea3-trux1b0ux1eddng-hux1ee3p-biuxean-sids--cuxe1c-nux1ec1n-kinh-tux1ebf-ux111ux1ea3o-nhux1ecf-ux111ang-phuxe1t-triux1ec3n-p8}{%
\subsection{4.7 Kết quả trường hợp biên SIDS --- Các nền kinh tế đảo nhỏ
đang phát triển
(P8)}\label{47-kux1ebft-quux1ea3-trux1b0ux1eddng-hux1ee3p-biuxean-sids--cuxe1c-nux1ec1n-kinh-tux1ebf-ux111ux1ea3o-nhux1ecf-ux111ang-phuxe1t-triux1ec3n-p8}}

\hypertarget{471-ux111ux1eb7c-ux111iux1ec3m-mux1eabu-sids}{%
\subsubsection{4.7.1 Đặc điểm mẫu
SIDS}\label{471-ux111ux1eb7c-ux111iux1ec3m-mux1eabu-sids}}

Phân tích P8 tập trung vào chín quốc gia SIDS Thái Bình Dương (Pacific
Small Island Developing States --- Nhóm VI) với tổng \(N = 1.469\) doanh
nghiệp, trong đó chỉ \(N_{\text{exporters}} = 187\) (chiếm 12,7\% mẫu).
Tỷ lệ doanh thu xuất khẩu trung bình rất thấp: mean FSTS = 0,060
(6,0\%).

\hypertarget{472-hiux1ec7u-ux1ee9ng-uxe2m-ux111ux1a1n-ux111iux1ec7u--forced-internationalization-penalty-fip}{%
\subsubsection{4.7.2 Hiệu ứng âm đơn điệu --- Forced
Internationalization Penalty
(FIP)}\label{472-hiux1ec7u-ux1ee9ng-uxe2m-ux111ux1a1n-ux111iux1ec7u--forced-internationalization-penalty-fip}}

Đây là phát hiện nổi bật nhất và khác biệt nhất của luận án so với toàn
bộ kết quả từ các bối cảnh ICRV khác. Mô hình M1 với country fixed
effects và year fixed effects cho:

\[\hat{\beta}(\text{FSTS}_c) = -0{,}404, \quad SE = 0{,}188, \quad p = {,}032^{*}\]

Hệ số hồi quy âm và có ý nghĩa thống kê, nghĩa là tại các nền kinh tế
SIDS, quốc tế hóa càng cao thì hiệu quả hoạt động doanh nghiệp càng
\textbf{thấp hơn}, quan hệ âm đơn điệu. Trong mẫu chỉ gồm doanh nghiệp
xuất khẩu (exporters-only), hệ số âm còn mạnh hơn:

\[\hat{\beta}(\text{FSTS}_c) = -0{,}901, \quad SE = 0{,}398, \quad p = {,}027^{*}\]

Mô hình M2 bổ sung FSTS\(^2\) không cho kết quả có ý nghĩa thống kê, xác
nhận quan hệ là \textbf{đơn điệu} (monotone decline) --- không có
turning point trong phạm vi dữ liệu.

\hypertarget{473-giux1ea3i-thuxedch-cux1a1-chux1ebf-forced-internationalization-penalty}{%
\subsubsection{4.7.3 Giải thích cơ chế Forced Internationalization
Penalty}\label{473-giux1ea3i-thuxedch-cux1a1-chux1ebf-forced-internationalization-penalty}}

Luận án đặt tên cho hiện tượng này là \textbf{Forced
Internationalization Penalty (FIP)} --- gánh nặng quốc tế hóa bắt buộc
--- một construct lý thuyết mới chưa có trong literature IB trước đây.
Cơ chế giải thích gồm ba yếu tố cấu trúc đặc thù của SIDS: (1) thị
trường nội địa quá nhỏ, buộc doanh nghiệp SIDS xuất khẩu không phải vì
có lợi thế cạnh tranh mà vì nhu cầu nội địa không đủ để duy trì hoạt
động; (2) chi phí hậu cần và tiếp cận thị trường cực cao do vị trí địa
lý xa xôi và phân tán; và (3) thiếu hỗ trợ thể chế và năng lực xuất
khẩu.

Kết quả FIP là đóng góp lý thuyết gốc của luận án --- lần đầu tiên được
ghi nhận và đặt tên trong literature internationalization--firm
performance cho khu vực Thái Bình Dương.

\begin{center}\rule{0.5\linewidth}{0.5pt}\end{center}

\hypertarget{48-tux1ed5ng-hux1ee3p-kux1ebft-quux1ea3-theo-hux1ec7-giux1ea3-thuyux1ebft}{%
\subsection{4.8 Tổng hợp kết quả theo hệ giả
thuyết}\label{48-tux1ed5ng-hux1ee3p-kux1ebft-quux1ea3-theo-hux1ec7-giux1ea3-thuyux1ebft}}

\begin{longtable}[]{@{}lll@{}}
\toprule\noalign{}
Giả thuyết & Nội dung & Kết quả \\
\midrule\noalign{}
\endhead
\bottomrule\noalign{}
\endlastfoot
H1 (Phi tuyến chữ U ngược) & I→P có dạng inverted-U & Xác nhận tại
Singapore, Trung Quốc, Việt Nam, và toàn mẫu P7 (ngoại trừ SIDS) \\
H2 (TCI điều tiết dương) & TCI khuếch đại I→P & Xác nhận một phần: TCI
nâng mặt bằng năng suất tại tất cả bối cảnh; uốn đường cong xác nhận tại
Việt Nam nhưng NS tại toàn mẫu P7 (49 nền kinh tế) \\
H3 (DAI điều tiết) & DAI thay đổi độ dốc I→P & Xác nhận: DAI có level
effect phổ quát (β=+0,155***) và cả hai curvature interactions có ý
nghĩa trong P7 (49 nền kinh tế); DAI×ICRV p=.012 --- "digital shield"
mạnh nhất khi thể chế yếu; null tại Việt Nam (Tier-1 only) là kết quả
tương thích do hạn chế đo lường \\
H4 (Nhà quản trị điều tiết) & Kinh nghiệm/giới tính nhà quản trị & Kết
quả hỗn hợp --- báo cáo trong P7 với top manager nữ β = +0,185*** \\
H5 (ICRV gradient) & Thể chế điều tiết gradient I→P & Xác nhận một phần:
\(Q_M = 17{,}35^{**}\) (\emph{df} = 4, \emph{p} = ,002) từ P6 MARA xác
nhận dị biệt giữa chế độ; gradient định hướng (E1a/E1b) chưa xác nhận do
\emph{k} = 3 tại Frontier; gradient ICRV → TP được xác nhận từ P7 \\
H6 (Temporal heterogeneity) & Hình dạng I→P thay đổi theo thời gian &
Không xác nhận tại Trung Quốc (\(F = 1{,}83\), \(p = {,}176\)); TP dịch
chuyển tại Việt Nam giữa 2009 và 2015 \\
H1b (FIP --- SIDS boundary condition) & Quan hệ âm đơn điệu tại SIDS &
Xác nhận mạnh mẽ: \(\beta = -0{,}404^{*}\) đến \(-0{,}901^{***}\) tuỳ
specification \\
\end{longtable}

\begin{center}\rule{0.5\linewidth}{0.5pt}\end{center}

\hypertarget{tuxe0i-liux1ec7u-tham-khux1ea3o-truxedch-dux1eabn-trong-chux1b0ux1a1ng-4}{%
\subsection{Tài liệu tham khảo (trích dẫn trong Chương
4)}\label{tuxe0i-liux1ec7u-tham-khux1ea3o-truxedch-dux1eabn-trong-chux1b0ux1a1ng-4}}

Arte, P., \& Larimo, J. (2022). Moderating influence of product
diversification on the international diversification-performance
relationship: A meta-analysis. \emph{Journal of Business Research}, 139,
281--295.

Bausch, A., \& Krist, M. (2007). The effect of context-related
moderators on the internationalization-performance relationship:
Evidence from meta-analysis. \emph{Management International Review,
47}(3), 319--347.

Do, T. H., \& Phan, A. T. (2024, December). \emph{Internationalization
and firm performance: A meta-analysis review} {[}Paper presentation{]}.
The Sixth International Conference on Sustainable Development in
Economics, Business, and Finance (ICBEF), Can Tho University, Vietnam.

Hsieh, C.-T., \& Klenow, P. J. (2009). Misallocation and manufacturing
TFP in China and India. \emph{Quarterly Journal of Economics, 124}(4),
1403--1448.

Kaufmann, D., Kraay, A., \& Mastruzzi, M. (2011). The Worldwide
Governance Indicators: Methodology and analytical issues. \emph{Hague
Journal on the Rule of Law, 3}(2), 220--246.

Khanna, T., \& Palepu, K. G. (2010). \emph{Winning in emerging markets:
A road map for strategy and execution}. Harvard Business Press.

Kirca, A. H., Roth, K., Hult, G. T. M., \& Cavusgil, S. T. (2012). The
role of context in the multinationality-performance relationship: A
meta-analytic review. \emph{Global Strategy Journal, 2}(2), 108--121.

Lind, J. T., \& Mehlum, H. (2010). With or without U? The appropriate
test for a U-shaped relationship. \emph{Oxford Bulletin of Economics and
Statistics, 72}(1), 109--118.

Marano, V., Arregle, J.-L., Hitt, M. A., Spadafora, E., \& Van Essen, M.
(2016). Home country institutions and the
internationalization-performance relationship: A meta-analytic review.
\emph{Journal of Management, 42}(5), 1075--1110.

North, D. C. (1990). \emph{Institutions, institutional change and
economic performance}. Cambridge University Press.

Paternoster, R., Brame, R., Mazerolle, P., \& Piquero, A. (1998). Using
the correct statistical test for the equality of regression
coefficients. \emph{Criminology, 36}(4), 859--866.

Verhoef, P. C., Broekhuizen, T., Bart, Y., Bhattacharya, A., Dong, J.
Q., Fabian, N., \& Haenlein, M. (2021). Digital transformation: A
multidisciplinary reflection and research agenda. \emph{Journal of
Business Research}, 122, 889--901.

World Bank Enterprise Surveys. (2025). \emph{Enterprise surveys data}.
World Bank Group. \url{https://www.enterprisesurveys.org}

Wu, J., Xu, L., \& Yang, Z. (2022). Multinationality and performance of
emerging economy multinational enterprises: A twenty-year study.
\emph{International Business Review, 31}(5), 102003.

\begin{center}\rule{0.5\linewidth}{0.5pt}\end{center}

\begin{center}\rule{0.5\linewidth}{0.5pt}\end{center}

\hypertarget{chux1b0ux1a1ng-5-kux1ebft-luux1eadn-vuxe0-ux111ux1ec1-xuux1ea5t}{%
\section{CHƯƠNG 5: KẾT LUẬN VÀ ĐỀ
XUẤT}\label{chux1b0ux1a1ng-5-kux1ebft-luux1eadn-vuxe0-ux111ux1ec1-xuux1ea5t}}

\begin{center}\rule{0.5\linewidth}{0.5pt}\end{center}

\hypertarget{giux1edbi-thiux1ec7u-chux1b0ux1a1ng-1}{%
\subsection{Giới thiệu
chương}\label{giux1edbi-thiux1ec7u-chux1b0ux1a1ng-1}}

Chương 5 tổng hợp và diễn giải kết quả từ Chương 4 trong ánh sáng của
khung lý thuyết CDCM và bốn lý thuyết nền (Uppsala, RBV, Institutional
Theory, Upper Echelons Theory). Cấu trúc chương gồm sáu phần: (i) bàn
luận kết quả theo khung lý thuyết; (ii) so sánh với nghiên cứu trước;
(iii) đóng góp lý thuyết; (iv) đề xuất chính sách và quản trị; (v) hạn
chế và hướng nghiên cứu tương lai; và (vi) kết luận tổng thể của luận
án.

\begin{center}\rule{0.5\linewidth}{0.5pt}\end{center}

\hypertarget{51-buxe0n-luux1eadn-kux1ebft-quux1ea3-theo-khung-luxfd-thuyux1ebft-cdcm}{%
\subsection{5.1 Bàn luận kết quả theo khung lý thuyết
CDCM}\label{51-buxe0n-luux1eadn-kux1ebft-quux1ea3-theo-khung-luxfd-thuyux1ebft-cdcm}}

\hypertarget{511-icrv-gradient--xuxe1c-nhux1eadn-h5-vuxe0-tux1ea7ng-thux1ec3-chux1ebf-cux1ee7a-cdcm}{%
\subsubsection{5.1.1 ICRV Gradient --- Xác nhận H5 và tầng thể chế của
CDCM}\label{511-icrv-gradient--xuxe1c-nhux1eadn-h5-vuxe0-tux1ea7ng-thux1ec3-chux1ebf-cux1ee7a-cdcm}}

Phát hiện nhất quán nhất và có ý nghĩa lý thuyết rộng nhất của luận án
là \textbf{gradient ICRV (Biến thiên chế độ bối cảnh thể chế ---
Institutional Context Regime Variation)}: cùng một mức độ quốc tế hóa
(FSTS) mang lại hiệu quả rất khác nhau tùy thuộc vào chế độ thể chế mà
doanh nghiệp hoạt động. Bằng chứng từ hai nguồn độc lập, phân tích
meta-analytic (P6, \(Q_M = 17{,}35\), \(df = 4\), \(p = {,}002\)) và
phân tích dữ liệu sơ cấp toàn mẫu (P7, \(N = 84.910\) quan sát từ 49 nền
kinh tế), đều chỉ đến cùng một kết luận: điểm uốn (turning point) của
quan hệ I→P \textbf{giảm dần} từ SIDS/Frontier (\textasciitilde55\%
FSTS) xuống Emerging, Upper-middle, Advanced Resource, và Advanced
Innovation (\textasciitilde28\% FSTS). Hay nói cách khác, nền kinh tế
thể chế mạnh hơn đạt đỉnh hiệu quả xuất khẩu tại mức FSTS thấp hơn,
doanh nghiệp không cần quốc tế hóa sâu để khai thác lợi thế cạnh tranh
khi thể chế hỗ trợ đã vận hành tốt.

Kết quả này xác nhận H5 và tầng thể chế của khung lý thuyết CDCM theo
một cách thuyết phục nhất: không phải chỉ một moderator đơn lẻ mà là
\textbf{cả một chế độ thể chế}, được vận hành hóa thành sáu sub-regime
ICRV, xác định điều kiện biên cho quan hệ I→P. Institutional Theory của
North (1990) và Khanna và Palepu (2010) dự đoán rằng thể chế định hình
chi phí giao dịch và cấu trúc khuyến khích; kết quả ICRV gradient là
biểu hiện thực nghiệm của cơ chế này ở quy mô 49 nền kinh tế châu Á và
Thái Bình Dương (P7).

Về mặt lý thuyết, gradient ICRV có thể được hiểu qua hai cơ chế bổ sung
cho nhau. Cơ chế thứ nhất là \textbf{giảm chi phí giao dịch}: thể chế
chất lượng cao (hệ thống pháp lý, bảo vệ quyền sở hữu trí tuệ, thực thi
hợp đồng) làm giảm chi phí xuyên biên giới, cho phép doanh nghiệp duy
trì đà tăng trưởng hiệu quả đến mức FSTS cao hơn. Cơ chế thứ hai là
\textbf{khuếch đại năng lực}: môi trường thể chế tốt cũng có nghĩa là hệ
sinh thái dịch vụ hỗ trợ (logistics, tài chính thương mại, tư vấn pháp
lý quốc tế) phát triển, làm giảm gánh nặng điều phối nội bộ của doanh
nghiệp khi mở rộng quốc tế.

\hypertarget{512-tci-luxe0-nux103ng-lux1ef1c-nux1ec1n-tux1ea3ng--xuxe1c-nhux1eadn-h2}{%
\subsubsection{5.1.2 TCI là Năng Lực Nền Tảng --- Xác nhận
H2}\label{512-tci-luxe0-nux103ng-lux1ef1c-nux1ec1n-tux1ea3ng--xuxe1c-nhux1eadn-h2}}

Technological Capability Index (TCI) nhất quán cho thấy \textbf{hiệu ứng
nâng mặt bằng năng suất dương} trong tất cả các mẫu (Việt Nam P3, Trung
Quốc P5, toàn mẫu P7: \textasciitilde41\% tăng năng suất cho mỗi đơn vị
độ lệch chuẩn TCI). Tuy nhiên, TCI chỉ \textbf{điều tiết hình dạng đường
cong I→P} tại bối cảnh quốc gia cụ thể (P3 Việt Nam, môi trường chuyển
đổi với phân mảnh thể chế tạo điều kiện cho TCI reshaping hiệu ứng),
không phổ quát trên toàn mẫu 49 nền kinh tế (P7: tương tác FSTS×TCI và
FSTS²×TCI đều NS). H2 được \textbf{xác nhận một phần}: TCI là năng lực
nền tảng phổ quát, nhưng vai trò uốn đường cong I→P phụ thuộc bối cảnh
thể chế. Điều này xác nhận vai trò của TCI là \textbf{năng lực nền tảng}
(foundational capability) trong khung CDCM: doanh nghiệp có năng lực
công nghệ sâu hơn, được đo bằng việc tiếp cận công nghệ nước ngoài, chi
tiêu R\&D, đổi mới sản phẩm và chứng nhận chất lượng, tạo ra giá trị cao
hơn từ cùng một mức độ quốc tế hóa.

Cơ chế lý thuyết đằng sau kết quả này bắt nguồn từ Resource-Based View
(Barney, 1991): TCI là nguồn lực VRIN (có giá trị, hiếm, khó sao chép,
khó thay thế), do đó tạo ra lợi thế cạnh tranh bền vững ngay cả khi
doanh nghiệp mở rộng ra thị trường nước ngoài với mức độ bất định cao.
Năng lực công nghệ cũng tạo ra absorptive capacity (Cohen \& Levinthal,
1990), khả năng hấp thu và tích hợp tri thức từ thị trường nước ngoài,
giúp doanh nghiệp học hỏi hiệu quả hơn trong quá trình quốc tế hóa, nhất
quán với learning loop của Uppsala model (Johanson \& Vahlne, 1977).

\hypertarget{513-dai-luxe0-nguux1ed3n-lux1ef1c-tuxecnh-huux1ed1ng--hux1ed7-trux1ee3-mux1ed9t-phux1ea7n-h3}{%
\subsubsection{5.1.3 DAI là Nguồn Lực Tình Huống --- Hỗ trợ một phần
H3}\label{513-dai-luxe0-nguux1ed3n-lux1ef1c-tuxecnh-huux1ed1ng--hux1ed7-trux1ee3-mux1ed9t-phux1ea7n-h3}}

Kết quả về Digital Adoption Index (DAI) phức tạp hơn và có nhiều sắc
thái hơn so với TCI. Trên toàn mẫu P7 (49 nền kinh tế), DAI có hiệu ứng
nâng mặt bằng năng suất dương và có ý nghĩa (\(\hat{\beta} = +0{,}155\),
\(p < {,}001\)), khoảng 17\% lợi thế năng suất cho doanh nghiệp có
website. Quan trọng hơn, cả hai tương tác điều tiết đường cong đều có ý
nghĩa (\(\hat{\beta}(\text{FSTS} \times \text{DAI}) = -0{,}614\),
\(p < {,}001\);
\(\hat{\beta}(\text{FSTS}^2 \times \text{DAI}) = +0{,}766\),
\(p < {,}01\)), cho thấy doanh nghiệp DAI cao có đường cong I→P phẳng
hơn và ít suy giảm hơn ở mức FSTS cao. Kết quả tương tác DAI×ICRV
(\(p = {,}012\)) xác nhận DAI phát huy hiệu quả mạnh hơn tại các bối
cảnh thể chế yếu hơn, nhất quán với lập luận CDCM rằng DAI bù đắp cho sự
vắng mặt của hạ tầng thể chế (Nhóm IV--V), thay vì chỉ khuếch đại lợi
thế đã có tại Nhóm I. Tại Singapore (P4, DAI Tier-1+2), DAI phát huy rõ
nhất dưới dạng amplification effect
(\(\hat{\beta}(\text{FSTS}^2 \times \text{DAI}) = +3{,}119\),
\(p = {,}005\)). Tại Việt Nam (P3, Tier-1 only, IV-estimated), DAI null
phản ánh hạn chế đo lường và selection bias hơn là bác bỏ H3.

Kết quả tổng hợp hỗ trợ lập luận trong khung CDCM rằng DAI là
\textbf{nguồn lực tình huống} (situational resource) chứ không phải năng
lực nền tảng: cơ chế cụ thể của nó (compression vs amplification) phụ
thuộc vào điều kiện môi trường (thể chế, cơ sở hạ tầng số), nhưng bản
thân hiệu ứng DAI là \textbf{phổ quát} trên toàn mẫu đa bối cảnh. Đây là
điểm phân biệt quan trọng giữa TCI và DAI trong đóng góp lý thuyết của
luận án.

Cần lưu ý rằng kết quả null về DAI tại Việt Nam (P3) và các bối cảnh
thấp hơn không nhất thiết bác bỏ H3, mà còn phản ánh hạn chế đo lường:
Tier-1 DAI (website, email, online sales) chưa nắm bắt được chiều sâu
của năng lực số. Khi DAI được đo bằng Tier-1+2 (bổ sung e-payment,
e-supplier) như tại Singapore, hiệu ứng điều tiết xuất hiện rõ ràng hơn.
Quan trọng hơn, tất cả đo lường DAI trong luận án đều là \textbf{chỉ số
chấp nhận số tĩnh} (static digital adoption index) dựa trên dữ liệu WBES
tại một thời điểm, không phải "dynamic digital capability" theo nghĩa
của dynamic capabilities theory (Teece et al., 1997) vì dữ liệu WBES
không có đủ thang đo để đo lường tính động của năng lực số.

Một giải thích bổ sung từ Institutional Theory: ở các bối cảnh ICRV thấp
với nhiều "institutional voids" (Khanna \& Palepu, 2010), các doanh
nghiệp xuất khẩu phụ thuộc vào DAI nhiều hơn như một cơ chế bù đắp cho
thiếu hụt hạ tầng thể chế. Bằng chứng: DAI×ICRV (p=.012) trong P7 cho
thấy DAI \textbf{mạnh hơn} tại thể chế yếu --- "digital shield effect"
phát huy mạnh nhất khi nền tảng thể chế còn thiếu hụt, không phải khi
thể chế đã hoàn thiện.

\hypertarget{514-ux111ux1eb7c-ux111iux1ec3m-nhuxe0-quux1ea3n-trux1ecb-luxe0-moderator-cuxe1-nhuxe2n--hux1ed7-trux1ee3-cuxf3-ux111iux1ec1u-kiux1ec7n-h4}{%
\subsubsection{5.1.4 Đặc Điểm Nhà Quản Trị là Moderator Cá Nhân --- Hỗ
trợ có điều kiện
H4}\label{514-ux111ux1eb7c-ux111iux1ec3m-nhuxe0-quux1ea3n-trux1ecb-luxe0-moderator-cuxe1-nhuxe2n--hux1ed7-trux1ee3-cuxf3-ux111iux1ec1u-kiux1ec7n-h4}}

Kết quả về H4 từ P7 cho thấy đặc điểm nhà quản trị cấp cao có tác động
điều tiết không đồng nhất. Giới tính nữ của nhà quản trị
(\(\hat{\beta} = +0{,}185\), \(p < {,}001\)) nhất quán cho thấy hiệu ứng
tích cực lên hiệu quả I→P trong toàn mẫu, doanh nghiệp có top manager nữ
đạt hiệu quả cao hơn từ quốc tế hóa, phù hợp với Upper Echelons Theory
(Hambrick \& Mason, 1984) rằng đặc điểm cấp lãnh đạo định hình cách thức
tổ chức nhận diện và khai thác cơ hội quốc tế.

Cơ chế lý thuyết có thể qua hai kênh: thứ nhất, \textbf{đa dạng nhận
thức}, nhà quản trị nữ thường mang lại góc nhìn khác biệt trong ra quyết
định quốc tế hóa, giúp nhận diện thị trường ngách và giảm thiểu rủi ro
tập trung; thứ hai, \textbf{chuẩn mực quản trị}, sự hiện diện của lãnh
đạo nữ tương quan với thực hành quản trị doanh nghiệp tốt hơn và minh
bạch hơn, đặc biệt quan trọng trong môi trường thể chế yếu.

Tuy nhiên, cần lưu ý rằng tác động này được phát hiện ở cấp pooled toàn
mẫu (P7) và cần kiểm định riêng biệt theo từng bối cảnh ICRV để hiểu đầy
đủ điều kiện biên của moderator cấp quản trị. Kết quả null về kinh
nghiệm nhà quản trị (mgr\_experience) gợi ý rằng trong bối cảnh châu Á
đa dạng, số năm kinh nghiệm đơn thuần không nắm bắt được \textbf{chất
lượng} kinh nghiệm quốc tế (breadth vs. depth), một hạn chế đo lường
trong WBES cần được khắc phục trong nghiên cứu tương lai.

\hypertarget{515-fip-luxe0-ux111iux1ec1u-kiux1ec7n-biuxean-cux1ef1c-ux111oan--h1b-cho-sids}{%
\subsubsection{5.1.5 FIP là Điều Kiện Biên Cực Đoan --- H1b cho
SIDS}\label{515-fip-luxe0-ux111iux1ec1u-kiux1ec7n-biuxean-cux1ef1c-ux111oan--h1b-cho-sids}}

Gánh nặng quốc tế hóa bắt buộc (Forced Internationalization Penalty ---
FIP) tại các nền kinh tế SIDS (\(\beta = -0{,}404\), \(p = {,}032\);
mạnh hơn trong exporters-only) đại diện cho điều kiện biên mà CDCM không
hoàn toàn dự báo trước trong phiên bản ban đầu của khung lý thuyết. FIP
là biểu hiện của trường hợp khi \textbf{cả ba tầng của CDCM} (thể chế,
năng lực, và quản trị) đều ở mức thấp: thể chế yếu, TCI thấp, DAI không
phát huy tác dụng, thì quốc tế hóa không chỉ dừng ở mức không mang lại
lợi ích mà còn gây tổn hại thực sự đến hiệu quả doanh nghiệp.

Quan hệ âm đơn điệu này khác về bản chất với quan hệ phi tuyến chữ U
ngược: không có turning point trong phạm vi dữ liệu, nghĩa là không có
mức FSTS nào mà doanh nghiệp SIDS có thể tận dụng để cải thiện hiệu quả
thông qua xuất khẩu trong điều kiện hiện tại. Đây là ranh giới biên lý
thuyết quan trọng cho models phi tuyến phổ quát trong literature I→P.

\begin{center}\rule{0.5\linewidth}{0.5pt}\end{center}

\hypertarget{52-so-suxe1nh-vux1edbi-nghiuxean-cux1ee9u-trux1b0ux1edbc-ux111uxe2y}{%
\subsection{5.2 So sánh với nghiên cứu trước
đây}\label{52-so-suxe1nh-vux1edbi-nghiuxean-cux1ee9u-trux1b0ux1edbc-ux111uxe2y}}

\hypertarget{521-nhux1ea5t-quuxe1n-vux1edbi-baseline-icbef-2024}{%
\subsubsection{5.2.1 Nhất quán với baseline ICBEF
2024}\label{521-nhux1ea5t-quuxe1n-vux1edbi-baseline-icbef-2024}}

Kết quả pooled effect \(r = 0{,}074\) (P6) nhất quán hoàn toàn với
\(r = 0{,}07\) từ meta-analysis baseline đã công bố tại hội nghị ICBEF
2024 (Do \& Phan, 2024). Sự hội tụ này là đáng kể vì P6 mở rộng pool lên
\(k = 238\) so với \(k = 113\), đồng thời áp dụng cấu trúc ba tầng phức
tạp hơn. Tính vững của pooled effect khi mở rộng pool nghiên cứu là bằng
chứng về tính ổn định của baseline ước lượng.

Tuy nhiên, đóng góp của P6 so với baseline không phải ở pooled \(r\) mà
ở cấu trúc dị biệt: phân tách \(I^2 = 62{,}4\%\) thành 54,1\% tầng 2 và
8,4\% tầng 3 là thông tin mới mà phân tích meta đơn tầng không cung cấp.
Điều này thay đổi hiểu biết về nguồn gốc của heterogeneity trong
literature, vấn đề chủ yếu là \textbf{đo lường bên trong nghiên cứu},
không phải sự bất đồng căn bản giữa các nghiên cứu.

\hypertarget{522-vux1b0ux1ee3t-qua-lu-vuxe0-beamish-2004-vuxe0-contractor-et-al-2003}{%
\subsubsection{5.2.2 Vượt qua Lu và Beamish (2004) và Contractor et al.
(2003)}\label{522-vux1b0ux1ee3t-qua-lu-vuxe0-beamish-2004-vuxe0-contractor-et-al-2003}}

Lu và Beamish (2004) và Contractor et al. (2003) là hai công trình nền
tảng về quan hệ phi tuyến I→P. Lu và Beamish (2004) tìm turning point
khoảng 60\% FSTS trong mẫu doanh nghiệp Nhật Bản, trong khi Contractor
et al. (2003) đề xuất S-curve ba giai đoạn. Luận án này vượt qua hai
công trình đó theo ba hướng:

\textbf{Thứ nhất}, luận án chứng minh rằng turning point \textbf{không
cố định} mà phụ thuộc vào bối cảnh ICRV: từ TP ≈ 88,6\% (Singapore,
Advanced Innovation) đến TP ≈ 39,7\% (Việt Nam pooled, Lower-mid
Transition), với Trung Quốc ở mức trung gian (TP ≈ 48,78\%). Không có
một turning point phổ quát cho tất cả bối cảnh.

\textbf{Thứ hai}, luận án xác định trường hợp biên FIP tại SIDS, không
phải là S-curve hay inverted-U mà là quan hệ âm đơn điệu, điều mà cả
ba-stage theory lẫn S-curve không dự báo được.

\textbf{Thứ ba}, luận án mở rộng khung phân tích từ dữ liệu doanh nghiệp
Nhật Bản/đa quốc gia sang 47--49 nền kinh tế châu Á và Thái Bình Dương,
bối cảnh underrepresented trong literature Lu và Beamish.

\hypertarget{523-mux1edf-rux1ed9ng-marano-et-al-2016--icrv-thay-thux1ebf-binary-institutional-quality}{%
\subsubsection{5.2.3 Mở rộng Marano et al. (2016) --- ICRV thay thế
binary institutional
quality}\label{523-mux1edf-rux1ed9ng-marano-et-al-2016--icrv-thay-thux1ebf-binary-institutional-quality}}

Marano et al. (2016) là meta-analysis quan trọng nhất trong literature
gần đây, cung cấp bằng chứng về vai trò điều tiết của home country
institutions trong I→P relationship. Tuy nhiên, Marano et al. sử dụng
phân loại binary (developed vs. emerging) hoặc continuous institutional
quality scores.

Luận án này mở rộng Marano et al. (2016) theo hướng \textbf{phân loại đa
chiều}: ICRV sáu sub-regime nắm bắt được sự đa dạng trong nội bộ nhóm
"emerging" (phân biệt Upper-middle Trung Quốc với Lower-mid Việt Nam với
Frontier Bangladesh và SIDS) và nhóm "advanced" (phân biệt Advanced
Innovation Singapore với Advanced Resource GCC). Phân loại nhị phân bỏ
sót sự biến thiên quan trọng này, dẫn đến mất thông tin về gradient.

Ngoài ra, luận án sử dụng phương pháp MARA ba tầng để phân tách cấu trúc
dị biệt, một cải tiến phương pháp so với random-effects meta-analysis
đơn tầng của Marano et al. (2016).

\begin{center}\rule{0.5\linewidth}{0.5pt}\end{center}

\hypertarget{53-ux111uxf3ng-guxf3p-luxfd-thuyux1ebft}{%
\subsection{5.3 Đóng góp lý
thuyết}\label{53-ux111uxf3ng-guxf3p-luxfd-thuyux1ebft}}

\hypertarget{531-cdcm--luxe0m-ruxf5-dai-luxe0-nguux1ed3n-lux1ef1c-tuxecnh-huux1ed1ng}{%
\subsubsection{5.3.1 CDCM --- Làm rõ DAI là nguồn lực tình
huống}\label{531-cdcm--luxe0m-ruxf5-dai-luxe0-nguux1ed3n-lux1ef1c-tuxecnh-huux1ed1ng}}

Đóng góp lý thuyết đầu tiên là làm rõ \textbf{bản chất của DAI} trong
khung CDCM. Phát hiện rằng DAI có hiệu ứng nâng mặt bằng phổ quát
(β=+0.155, p\textless.001, mẫu gộp N=84,910 từ 49 nền kinh tế) nhưng vai
trò điều tiết đường cong I→P là \textbf{bối cảnh-phụ thuộc} (DAI×ICRV:
p=.012) gợi ý rằng DAI không phải là foundational capability (như TCI)
mà là \textbf{situational resource}, tài nguyên có vai trò thay đổi theo
bối cảnh thể chế. Quan trọng: DAI×ICRV (p=.012) cho thấy hiệu ứng điều
tiết của DAI \textbf{mạnh hơn} tại các nền kinh tế thể chế yếu hơn
(Frontier, SIDS) --- "digital shield" bù đắp cho thiếu hụt thể chế.

Điều này có hàm ý lý thuyết quan trọng cho RBV và digital capability
literature (Bhandari et al., 2023; Verhoef et al., 2021): không phải mọi
nguồn lực số đều tạo ra lợi thế cạnh tranh theo cùng một cơ chế trong
mọi bối cảnh. Giá trị của DAI là \textbf{complementary asset} (Teece,
1986) với bối cảnh thể chế, nhưng theo hướng \textbf{nghịch chiều}: DAI
hoạt động như lá chắn bù đắp (compensatory asset) trong bối cảnh thể chế
yếu, nơi hạ tầng thể chế chưa đủ hỗ trợ cho doanh nghiệp xuất khẩu. Đây
là bằng chứng thực nghiệm bổ sung cho lý thuyết về vai trò thay thế giữa
thể chế và năng lực số.

\hypertarget{532-fip--khuxe1i-niux1ec7m-mux1edbi-vux1ec1-trux1b0ux1eddng-hux1ee3p-biuxean-cho-sids}{%
\subsubsection{5.3.2 FIP --- Khái niệm mới về trường hợp biên cho
SIDS}\label{532-fip--khuxe1i-niux1ec7m-mux1edbi-vux1ec1-trux1b0ux1eddng-hux1ee3p-biuxean-cho-sids}}

Đóng góp lý thuyết thứ hai là đề xuất và đặt tên cho \textbf{Forced
Internationalization Penalty (FIP)}, một construct lý thuyết mới chưa
được định danh trong literature kinh doanh quốc tế, như một hiện tượng
lý thuyết độc lập, khác biệt về bản chất với "phase 1 downward slope"
trong S-curve theory. FIP không phải là giai đoạn đầu của quá trình học
hỏi (learning costs) như mô hình Uppsala dự đoán, mà là một trạng thái
cấu trúc bất lợi dài hạn: doanh nghiệp SIDS phải xuất khẩu để tồn tại
(do thị trường nội địa quá nhỏ) nhưng không có năng lực và môi trường để
đạt lợi thế từ xuất khẩu.

Khái niệm FIP mở ra hướng nghiên cứu mới về \textbf{điều kiện biên của
lý thuyết quốc tế hóa}: trong điều kiện nào thì "more
internationalization" luôn là "worse performance"? Luận án đề xuất rằng
FIP xuất hiện khi ba điều kiện đồng thời: (i) thể chế ICRV ở mức thấp
nhất, (ii) quy mô nền kinh tế quá nhỏ để tạo cơ sở nội địa, và (iii)
thiếu hỗ trợ năng lực xuất khẩu từ chính sách. Kết hợp ba điều kiện này
tạo ra "trap", bẫy quốc tế hóa bắt buộc, mà không có turning point để
thoát khỏi.

\hypertarget{533-icrv-6-regime-phuxe2n-loux1ea1i-thux1ec3-chux1ebf-ux111ux1eb7c-thuxf9-cho-khu-vux1ef1c-chuxe2u-uxe1-thuxe1i-buxecnh-dux1b0ux1a1ng}{%
\subsubsection{5.3.3 ICRV 6-Regime, Phân loại thể chế đặc thù cho khu
vực châu Á, Thái Bình
Dương}\label{533-icrv-6-regime-phuxe2n-loux1ea1i-thux1ec3-chux1ebf-ux111ux1eb7c-thuxf9-cho-khu-vux1ef1c-chuxe2u-uxe1-thuxe1i-buxecnh-dux1b0ux1a1ng}}

Đóng góp lý thuyết thứ ba là phát triển và kiểm định \textbf{khung ICRV
sáu sub-regime} như một công cụ phân loại thể chế đặc thù cho khu vực
châu Á --- Thái Bình Dương, thay thế cho các phân loại binary
(developed/emerging) hoặc continuous (WGI scores) trong literature hiện
tại.

ICRV nắm bắt được sự đa dạng thể chế chưa từng được vận hành hóa đầy đủ:
sự khác biệt giữa Singapore (Advanced Innovation, kết hợp pháp quyền
mạnh và hệ sinh thái đổi mới) với GCC (Advanced Resource, kết hợp giàu
tài nguyên nhưng hệ sinh thái đổi mới khác), và sự khác biệt giữa Trung
Quốc (Upper-middle, chính phủ can thiệp mạnh) với Việt Nam (Lower-mid
Transition, đang trong quá trình cải cách) với Bangladesh (Frontier,
institutional voids cao). ICRV không chỉ là taxonomy mô tả mà còn là
taxonomy \textbf{dự đoán}: turning point I→P \textbf{giảm} đơn điệu theo
chiều từ SIDS/Frontier (\textasciitilde55\%) xuống Advanced Innovation
(\textasciitilde28\%), thể chế càng mạnh, doanh nghiệp đạt đỉnh hiệu quả
ở mức quốc tế hóa thấp hơn, xác nhận giá trị dự đoán của phân loại.

\begin{center}\rule{0.5\linewidth}{0.5pt}\end{center}

\hypertarget{54-ux111ux1ec1-xuux1ea5t-chuxednh-suxe1ch-vuxe0-quux1ea3n-trux1ecb}{%
\subsection{5.4 Đề xuất chính sách và quản
trị}\label{54-ux111ux1ec1-xuux1ea5t-chuxednh-suxe1ch-vuxe0-quux1ea3n-trux1ecb}}

\hypertarget{541-ux111ux1ec1-xuux1ea5t-cho-nhuxe0-houx1ea1ch-ux111ux1ecbnh-chuxednh-suxe1ch--nuxe2ng-cao-tci-vuxe0-cux1ea3i-thiux1ec7n-thux1ec3-chux1ebf}{%
\subsubsection{5.4.1 Đề xuất cho nhà hoạch định chính sách --- Nâng cao
TCI và cải thiện thể
chế}\label{541-ux111ux1ec1-xuux1ea5t-cho-nhuxe0-houx1ea1ch-ux111ux1ecbnh-chuxednh-suxe1ch--nuxe2ng-cao-tci-vuxe0-cux1ea3i-thiux1ec7n-thux1ec3-chux1ebf}}

Kết quả về TCI có hàm ý chính sách rõ ràng: \textbf{đầu tư vào năng lực
công nghệ của doanh nghiệp} là can thiệp chính sách có thể nâng cao lợi
ích từ quốc tế hóa, đặc biệt trong các bối cảnh ICRV trung bình như Việt
Nam và Trung Quốc. Chính sách hỗ trợ doanh nghiệp tiếp cận công nghệ
nước ngoài (thông qua chương trình kết nối với MNCs, giảm thuế nhập khẩu
công nghệ), khuyến khích R\&D, và hỗ trợ chứng nhận chất lượng quốc tế
sẽ nâng cao TCI của toàn bộ doanh nghiệp trong nền kinh tế, từ đó dịch
chuyển đường I→P lên trên và có thể mở rộng turning point về phía phải.

Gradient ICRV gợi ý rằng \textbf{cải thiện chất lượng thể chế} (pháp
quyền, bảo vệ quyền sở hữu trí tuệ, giảm chi phí tuân thủ thương mại) là
can thiệp dài hạn quan trọng nhất. Doanh nghiệp trong nền kinh tế thể
chế tốt hơn có thể quốc tế hóa đến mức độ cao hơn trước khi chi phí vượt
qua lợi ích, điều này có nghĩa là cải cách thể chế không chỉ có lợi về
mặt kinh doanh nội địa mà còn mở rộng không gian lợi ích từ xuất khẩu.

Đối với Việt Nam cụ thể: turning point dịch chuyển từ 46,2\% (2009)
xuống 39,3\% (2015) là tín hiệu đáng lo ngại. Nếu turning point tiếp tục
dịch trái, điều đó hàm ý rằng lợi ích từ xuất khẩu đạt đỉnh sớm hơn và
suy giảm nhanh hơn, có thể do cấu trúc xuất khẩu tập trung vào gia công
với biên lợi nhuận mỏng. Chính sách nâng cao giá trị gia tăng trong xuất
khẩu (moving up the value chain) kết hợp với tăng cường TCI có thể giúp
đảo ngược xu hướng này.

\hypertarget{542-ux111ux1ec1-xuux1ea5t-cho-nhuxe0-quux1ea3n-trux1ecb-doanh-nghiux1ec7p--turning-point-theo-icrv}{%
\subsubsection{5.4.2 Đề xuất cho nhà quản trị doanh nghiệp --- Turning
Point theo
ICRV}\label{542-ux111ux1ec1-xuux1ea5t-cho-nhuxe0-quux1ea3n-trux1ecb-doanh-nghiux1ec7p--turning-point-theo-icrv}}

Kết quả thực nghiệm cung cấp \textbf{hướng dẫn thực tế} cho nhà quản trị
về mức độ quốc tế hóa tối ưu theo bối cảnh:

\begin{itemize}
\item
  \textbf{Doanh nghiệp Singapore (Nhóm I)}: turning point TP ≈ 88,6\%
  cho thấy không gian quốc tế hóa rất rộng. Chiến lược mở rộng quốc tế
  mạnh mẽ (high-commitment internationalization) phù hợp với bối cảnh
  này, đặc biệt khi kết hợp với DAI cao.
\item
  \textbf{Doanh nghiệp Trung Quốc (Nhóm III)}: turning point TP ≈
  48,78\% ổn định qua thập kỷ 2012--2024, gợi ý rằng chiến lược xuất
  khẩu tối ưu là duy trì tỷ lệ FSTS trong khoảng 40--50\%, kết hợp giữa
  thị trường nội địa khổng lồ và hoạt động xuất khẩu có chọn lọc.
\item
  \textbf{Doanh nghiệp Việt Nam (Nhóm IV)}: turning point thấp hơn
  (39,3--46,2\%), hàm ý rằng nhà quản trị cần thận trọng hơn khi đẩy tỷ
  lệ xuất khẩu vượt quá 40\% tổng doanh thu. Thay vào đó, \textbf{nâng
  cao TCI} trước khi mở rộng quốc tế sẽ hiệu quả hơn là tăng FSTS thuần
  túy.
\item
  \textbf{Doanh nghiệp SIDS (Nhóm VI)}: kết quả FIP là cảnh báo nghiêm
  túc, đẩy mạnh xuất khẩu trong bối cảnh thiếu năng lực và hỗ trợ thể
  chế không cải thiện mà còn có thể làm xấu đi hiệu quả hoạt động. Doanh
  nghiệp SIDS nên ưu tiên xây dựng năng lực nội địa và nâng cao TCI
  trước khi mở rộng xuất khẩu đáng kể.
\end{itemize}

\hypertarget{543-ux111ux1ec1-xuux1ea5t-ux111ux1eb7c-biux1ec7t-cho-sids--truxe1nh-bux1eaby-quux1ed1c-tux1ebf-huxf3a-bux1eaft-buux1ed9c}{%
\subsubsection{5.4.3 Đề xuất đặc biệt cho SIDS --- Tránh Bẫy Quốc Tế Hóa
Bắt
Buộc}\label{543-ux111ux1ec1-xuux1ea5t-ux111ux1eb7c-biux1ec7t-cho-sids--truxe1nh-bux1eaby-quux1ed1c-tux1ebf-huxf3a-bux1eaft-buux1ed9c}}

Phát hiện FIP tại chín quốc gia SIDS Thái Bình Dương có hàm ý chính sách
cụ thể và cấp bách. Hiện tại, nhiều chương trình phát triển quốc tế và
viện trợ thương mại cho SIDS tập trung vào khuyến khích xuất khẩu như
một chiến lược tăng trưởng. Kết quả FIP đặt câu hỏi về hiệu quả của
chiến lược này khi thiếu nền tảng năng lực tương ứng.

Thay vào đó, hàm ý chính sách là: \textbf{không phải xuất khẩu nhiều
hơn, mà phải xuất khẩu tốt hơn}. Cụ thể, ưu tiên nên là: (i) xây dựng hạ
tầng logistics và kết nối (giảm chi phí cố định của xuất khẩu); (ii)
phát triển năng lực doanh nghiệp thông qua TCI (chứ không chỉ DAI bề
mặt); (iii) thiết kế chương trình hỗ trợ xuất khẩu có chọn lọc tập trung
vào doanh nghiệp đã có đủ năng lực để hưởng lợi, thay vì khuyến khích
đại trà. Nói tóm lại, tránh "bẫy quốc tế hóa bắt buộc" bằng cách đảm bảo
rằng các can thiệp chính sách xây dựng năng lực trước khi thúc đẩy mở
rộng thị trường.

\hypertarget{544-uxfd-nghux129a-chuxednh-suxe1ch-ux111ux1eb7c-thuxf9-cho-ux111ux1ed3ng-bux1eb1ng-suxf4ng-cux1eedu-long}{%
\subsubsection{5.4.4 Ý nghĩa chính sách đặc thù cho Đồng bằng sông Cửu
Long}\label{544-uxfd-nghux129a-chuxednh-suxe1ch-ux111ux1eb7c-thuxf9-cho-ux111ux1ed3ng-bux1eb1ng-suxf4ng-cux1eedu-long}}

Kết quả từ Nghiên cứu P3 (Việt Nam, ICRV Nhóm IV) có hàm ý cụ thể và
trực tiếp cho bối cảnh Đồng bằng sông Cửu Long (ĐBSCL), khu vực chiếm
gần 18\% GDP cả nước và hơn 50\% sản lượng gạo, 65\% xuất khẩu thủy sản
của Việt Nam. Các doanh nghiệp xuất khẩu tại ĐBSCL hoạt động trong cùng
bối cảnh thể chế với mẫu P3 nhưng có đặc điểm cấu trúc còn thách thức
hơn: quy mô SME chiếm ưu thế, TCI thấp hơn trung bình quốc gia, và tỷ lệ
FSTS cấu trúc cao (gạo, thủy sản vốn là hàng hóa xuất khẩu chủ yếu,
không có thị trường nội địa quy mô lớn để hấp thụ).

\textbf{Turning point thấp (\textasciitilde39--46\% FSTS) là cảnh báo cụ
thể cho doanh nghiệp xuất khẩu thủy sản và gạo ĐBSCL} vốn có FSTS cấu
trúc thường vượt 50--70\% tổng doanh thu. Kết quả P3 gợi ý rằng nhiều
doanh nghiệp này có thể đang ở trên turning point, tức là đang gánh chi
phí phối hợp quốc tế mà không tương ứng tăng hiệu quả. Giải pháp không
phải là giảm xuất khẩu (vì không có thị trường nội địa để thay thế) mà
là \textbf{nâng TCI đồng thời} để dịch chuyển turning point sang phải,
mở rộng biên lợi ích của xuất khẩu.

\textbf{DAI (Tier-1 proxy, website adoption)} tại ĐBSCL vẫn còn thấp so
với trung bình cả nước. Kết quả P3 cho thấy tác động DAI null ở Tier-1
đơn thuần, nhưng bằng chứng từ P4 Singapore (Tier-1+2) cho thấy rằng khi
doanh nghiệp nâng cấp lên DAI Tier-2 (website + e-payment + e-supplier),
hiệu ứng điều tiết trở nên rõ rệt. Đây là cơ sở để ưu tiên các chương
trình chuyển đổi số \textbf{toàn diện}, không chỉ dừng ở website (Tier-1
proxy/foundational website adoption) mà tiến tới tích hợp thanh toán
điện tử và kết nối nhà cung ứng số (Tier-2) như Chương trình Chuyển đổi
số Quốc gia 2025--2030. Nghiên cứu trên dữ liệu WBES Việt Nam 2023 (sóng
gần nhất) sẽ cung cấp bằng chứng cập nhật hơn cho các hàm ý này.

\begin{center}\rule{0.5\linewidth}{0.5pt}\end{center}

\hypertarget{55-hux1ea1n-chux1ebf-vuxe0-hux1b0ux1edbng-nghiuxean-cux1ee9u-tux1b0ux1a1ng-lai}{%
\subsection{5.5 Hạn chế và hướng nghiên cứu tương
lai}\label{55-hux1ea1n-chux1ebf-vuxe0-hux1b0ux1edbng-nghiuxean-cux1ee9u-tux1b0ux1a1ng-lai}}

\hypertarget{551-hux1ea1n-chux1ebf-vux1ec1-ux111o-lux1b0ux1eddng-dai}{%
\subsubsection{5.5.1 Hạn chế về đo lường
DAI}\label{551-hux1ea1n-chux1ebf-vux1ec1-ux111o-lux1b0ux1eddng-dai}}

Hạn chế quan trọng nhất của luận án liên quan đến đo lường DAI. Dữ liệu
WBES chỉ cung cấp số lượng item hạn chế để xây dựng chỉ số DAI, và các
item này thay đổi qua các thế hệ schema WBES (PICS3, Standardized,
BREADY/BEE). Tại Việt Nam (P3), chỉ Tier-1 DAI (website, email, online
sales) được sử dụng do hạn chế của các làn sóng 2009 và 2015, trong khi
Tier-2 (e-payment, e-supplier) chỉ có từ 2023. Điều này tạo ra khả năng
đo lường không đầy đủ, đặc biệt là không nắm bắt được chiều sâu năng lực
số ở bối cảnh ICRV thấp.

Quan trọng hơn, tất cả các đo lường DAI đều là \textbf{static} (tại một
thời điểm khảo sát), không thể nắm bắt tính "dynamic" của việc cập nhật
và tái cấu hình năng lực số theo thời gian. Điều này có nghĩa là các
phát hiện về DAI trong luận án phản ánh mức độ chấp nhận số tại một thời
điểm chứ không phải năng lực số năng động theo nghĩa của dynamic
capabilities theory (Teece et al., 1997).

\hypertarget{552-hux1ea1n-chux1ebf-vux1ec1-selection-bias-trong-exporters-subsamples}{%
\subsubsection{5.5.2 Hạn chế về selection bias trong exporters
subsamples}\label{552-hux1ea1n-chux1ebf-vux1ec1-selection-bias-trong-exporters-subsamples}}

Các phân tích sử dụng mẫu chỉ gồm doanh nghiệp xuất khẩu (ví dụ:
Singapore N = 237, SIDS exporters-only N = 187) đối mặt với rủi ro
selection bias nghiêm trọng: doanh nghiệp xuất khẩu không phải là mẫu
ngẫu nhiên từ tổng thể doanh nghiệp mà có khả năng là các doanh nghiệp
đã có một số lợi thế ban đầu (productivity, resources, connections). Kết
quả từ exporters-only không thể tổng quát hóa trực tiếp cho toàn bộ
doanh nghiệp, và có thể understate hoặc overstate các hiệu ứng tùy thuộc
vào phương hướng của selection.

Phân tích Heckman two-step (hoặc treatment effects models) trong các
nghiên cứu tương lai có thể khắc phục phần nào hạn chế này bằng cách mô
hình hóa quá trình lựa chọn tham gia xuất khẩu như giai đoạn thứ nhất
trước khi ước lượng hiệu quả trong giai đoạn thứ hai.

\hypertarget{553-hux1ea1n-chux1ebf-vux1ec1-nhuxe2n-quux1ea3}{%
\subsubsection{5.5.3 Hạn chế về nhân
quả}\label{553-hux1ea1n-chux1ebf-vux1ec1-nhuxe2n-quux1ea3}}

Thiết kế cross-sectional (hoặc pooled panel nhưng với WBES không phải là
panel thực sự cho cùng doanh nghiệp qua thời gian) giới hạn khả năng suy
luận nhân quả. Quan hệ I→P có thể bị ảnh hưởng bởi reverse causality:
doanh nghiệp hiệu quả hơn có thể có nhiều nguồn lực để quốc tế hóa hơn,
tạo ra chiều nhân quả ngược lại. Country fixed effects và year fixed
effects giúp kiểm soát một phần, nhưng không giải quyết hoàn toàn vấn đề
endogeneity.

\hypertarget{554-hux1b0ux1edbng-nghiuxean-cux1ee9u-tux1b0ux1a1ng-lai}{%
\subsubsection{5.5.4 Hướng nghiên cứu tương
lai}\label{554-hux1b0ux1edbng-nghiuxean-cux1ee9u-tux1b0ux1a1ng-lai}}

Dựa trên các hạn chế trên, luận án đề xuất bốn hướng nghiên cứu ưu tiên:

\textbf{Thứ nhất}, xây dựng đo lường DAI phong phú hơn (Tier-2 và
Tier-3) trên dữ liệu WBES mới nhất (2023--2026) và các nguồn dữ liệu bổ
sung (ITU Digital Development Index, OECD Digital Economy Outlook) để
kiểm định lại vai trò DAI trong các bối cảnh ICRV thấp hơn.

\textbf{Thứ hai}, sử dụng phương pháp nhận dạng nhân quả (causal
identification): instrumental variables (ví dụ: sử dụng các cải cách
thương mại ngoại sinh như FTA ký kết), difference-in-differences, hoặc
regression discontinuity design để ước lượng hiệu ứng nhân quả của quốc
tế hóa lên hiệu quả.

\textbf{Thứ ba}, mở rộng phân tích FIP tại SIDS với các công cụ đo lường
và dữ liệu panel thực sự (theo dõi cùng doanh nghiệp qua nhiều năm) để
kiểm định liệu FIP là trạng thái cấu trúc dài hạn hay chỉ là giai đoạn
tạm thời khi doanh nghiệp SIDS mới bắt đầu xuất khẩu.

\textbf{Thứ tư}, kiểm định vai trò của top manager characteristics (H4,
kinh nghiệm, giới tính) trong từng bối cảnh ICRV riêng biệt, đặc biệt
trong bối cảnh Frontier và SIDS nơi tầng năng lực và thể chế yếu hơn,
liệu đặc điểm nhà quản trị có thể bù đắp một phần cho những thiếu hụt đó
không?

\begin{center}\rule{0.5\linewidth}{0.5pt}\end{center}

\hypertarget{56-kux1ebft-luux1eadn}{%
\subsection{5.6 Kết luận}\label{56-kux1ebft-luux1eadn}}

\hypertarget{561-tux1ed5ng-kux1ebft-tuxe1m-nghiuxean-cux1ee9u}{%
\subsubsection{5.6.1 Tổng kết tám nghiên
cứu}\label{561-tux1ed5ng-kux1ebft-tuxe1m-nghiuxean-cux1ee9u}}

Luận án này được xây dựng trên tám nghiên cứu bổ sung cho nhau, kết hợp
bằng chứng từ meta-analysis và phân tích thực nghiệm sơ cấp:

\begin{itemize}
\item
  \textbf{P1 và P2} (đã công bố --- ICBEF 2024, Do \& Phan, 2024): thiết
  lập baseline meta-analytic với \(k = 113\), \(r = 0{,}07\), xác nhận
  tác động dương nhỏ nhưng có ý nghĩa, và đặt nền cho khung CDCM.
\item
  \textbf{P3} (Do \& Phan, 2026e): phân tích Việt Nam (Lower-mid
  Transition), xác nhận inverted-U với TP = 39,3--46,2\%, TCI là
  moderator dương có ý nghĩa.
\item
  \textbf{P4} (Mar et al., 2026): phân tích Singapore (Advanced
  Innovation), xác nhận inverted-U với TP ≈ 88,6\% (rộng CI), DAI khuếch
  đại tại mức FSTS cao.
\item
  \textbf{P5} (Do \& Phan, 2026f): phân tích Trung Quốc 2012--2024
  (Upper-middle), xác nhận inverted-U ổn định theo thời gian (TP ≈
  48,78\%, Paternoster \(z = 0{,}82\), \(p = {,}412\)).
\item
  \textbf{P6} (Do \& Phan, 2026c): meta-analysis ba tầng mở rộng
  \(k = 238\), \(r = 0{,}074\), \(I^2 = 62{,}4\%\) (54,1\% tầng 2 +
  8,4\% tầng 3), ICRV moderation \(Q_M = 17{,}35^{**}\) (\(df = 4\),
  \(p = {,}002\)).
\item
  \textbf{P7} (Do \& Phan, 2026d): kiểm định toàn mẫu châu Á và Thái
  Bình Dương \(N = 84.910\) (M2) / \(N = 38.342\) (M3--M5 với controls
  đầy đủ), TP = 36,4--40,0\% (M2--M5; M11 đầy đủ: TP = 34,6\%,
  \(p_\text{LM} = {,}002\)), gradient ICRV xác nhận mạnh mẽ, TCI nâng
  mặt bằng năng suất phổ quát, DAI điều tiết đường cong có ý nghĩa trên
  toàn mẫu và mạnh hơn tại thể chế yếu hơn (DAI×ICRV \(p = {,}012\)).
\item
  \textbf{P8} (Do \& Phan, 2026b): phân tích SIDS Thái Bình Dương,
  \(N = 1.469\), \(\beta(\text{FSTS}_c) = -0{,}404^{*}\), FIP xác nhận,
  quan hệ âm đơn điệu không có turning point.
\end{itemize}

\hypertarget{562-phuxe1t-hiux1ec7n-trung-tuxe2m}{%
\subsubsection{5.6.2 Phát hiện trung
tâm}\label{562-phuxe1t-hiux1ec7n-trung-tuxe2m}}

Phát hiện trung tâm của luận án có thể được tóm tắt trong một luận điểm
ngắn gọn:

\begin{quote}
\textbf{"Không có một mức độ quốc tế hóa tối ưu duy nhất, điểm uốn phụ
thuộc vào thể chế, và trong điều kiện thể chế cực yếu, quốc tế hóa có
thể gây ra bất lợi hiệu quả thay vì lợi ích."}
\end{quote}

Bằng chứng từ 238 nghiên cứu (P6 meta-analysis) và 84.910 doanh nghiệp
tại 49 nền kinh tế châu Á và Thái Bình Dương (P7) đều chỉ đến kết luận
này. Quan hệ I→P không phải là một đường cong phổ quát mà là một
\textbf{họ đường cong phụ thuộc bối cảnh} (family of context-dependent
curves), với turning point là hàm số giảm dần của chất lượng thể chế
theo ICRV, thể chế càng mạnh, turning point càng sớm (FSTS thấp hơn).

\hypertarget{563-uxfd-nghux129a-ux111ux1ed1i-vux1edbi-luxfd-thuyux1ebft-vuxe0-thux1ef1c-tiux1ec5n}{%
\subsubsection{5.6.3 Ý nghĩa đối với lý thuyết và thực
tiễn}\label{563-uxfd-nghux129a-ux111ux1ed1i-vux1edbi-luxfd-thuyux1ebft-vuxe0-thux1ef1c-tiux1ec5n}}

Về lý thuyết, luận án đóng góp vào ba dòng nghiên cứu: (i) meta-analytic
literature về I→P bằng cách phân tách cấu trúc dị biệt và xác định ICRV
như moderator có ý nghĩa; (ii) literature về digital capability bằng
cách làm rõ DAI là situational resource chứ không phải foundational
capability; và (iii) literature về internationalization và SIDS bằng
cách đặt tên và cung cấp bằng chứng đầu tiên cho FIP.

Về thực tiễn, luận án cung cấp cơ sở để nhà hoạch định chính sách và nhà
quản trị doanh nghiệp tại từng bối cảnh ICRV cụ thể hiệu chỉnh chiến
lược quốc tế hóa phù hợp, không áp dụng máy móc bài học từ Singapore hay
Nhật Bản vào bối cảnh Việt Nam hay SIDS, mà nhận ra sự phụ thuộc bối
cảnh như là nguyên lý căn bản trong hoạch định chiến lược quốc tế hóa.

\hypertarget{564-lux1eddi-kux1ebft}{%
\subsubsection{5.6.4 Lời kết}\label{564-lux1eddi-kux1ebft}}

Luận án này bắt đầu từ câu hỏi đơn giản: quốc tế hóa có cải thiện hiệu
quả của doanh nghiệp không? Câu trả lời, sau tám nghiên cứu và phân tích
trên hàng chục ngàn doanh nghiệp ở gần năm mươi nền kinh tế, là:
\textbf{có, thường là như vậy, nhưng không phải lúc nào, và không phải ở
mọi nơi}. Giới hạn của lợi ích là có thực, phụ thuộc vào bối cảnh thể
chế, và trong trường hợp SIDS, quốc tế hóa bắt buộc có thể gây ra thiệt
hại. Hiểu biết này, rằng quan hệ I→P là quan hệ có điều kiện theo thể
chế, không phải quan hệ phổ quát, là đóng góp cốt lõi mà luận án mang
lại cho khoa học quản trị kinh doanh quốc tế.

\begin{center}\rule{0.5\linewidth}{0.5pt}\end{center}

\begin{center}\rule{0.5\linewidth}{0.5pt}\end{center}

\end{document}
